% 电机学课程笔记 —— 主文件（正文排版在 0.样式.tex，章节文件只放内容）
% 编译：xelatex 1.main.tex 两遍；成品见 build.ps1
\documentclass[11pt,a4paper,oneside,openany]{book}
\usepackage[UTF8,fontset=windows,scheme=plain]{ctex}
\usepackage{fontspec}
\setmainfont{Times New Roman}
\setsansfont{Microsoft YaHei}
\newcommand{\YHBbooktrue}{}
\newcommand{\YHBcourse}{电机学}
\newcommand{\YHBversion}{v0.1}
\newcommand{\YHBquote}{以磁路为根，以旋转磁场为魂。}
\newcommand{\YHBcoverSize}{40pt}
\newcommand{\YHBcoverLead}{48pt}
\usepackage{xcolor}
\definecolor{yhbAccent}{HTML}{8C3A2B}
\definecolor{yhbAccent2}{HTML}{B0704F}
% ============================================================
% 0.样式.tex —— 课程笔记统一正文排版
% 用法：主文件先定义课程信息与主色，再 % ============================================================
% 0.样式.tex —— 课程笔记统一正文排版
% 用法：主文件先定义课程信息与主色，再 % ============================================================
% 0.样式.tex —— 课程笔记统一正文排版
% 用法：主文件先定义课程信息与主色，再 \input{0.样式.tex}
%   \newcommand{\YHBcourse}{微积分A}      % 课程名（页眉与书签用）
%   \newcommand{\YHBversion}{v1.2}        % 版本号（封面用）
%   \newcommand{\YHBquote}{……}            % 封面题句
%   \definecolor{yhbAccent}{HTML}{24465C} % 主色（深）
%   \definecolor{yhbAccent2}{HTML}{347A80}% 辅色（浅）
%   \YHBbooktrue / \YHBbookfalse          % 是否 book 文档类（有 \chapter）
% 章节文件只写内容，不要重复 \usepackage、不要改字号。
% ============================================================

% ---------- 页面与字体 ----------
\usepackage{geometry}
\geometry{a4paper,left=24mm,right=24mm,top=23mm,bottom=24mm,headheight=15pt,headsep=7mm,footskip=13mm}
\usepackage{microtype}
\usepackage{amsmath,amssymb,amsfonts,mathtools}
\usepackage{esint}
\usepackage{graphicx}
\usepackage{adjustbox}
\usepackage{array,booktabs,multirow,makecell,tabularx,longtable}
\usepackage{xcolor,colortbl}
\usepackage{enumitem}
\usepackage{needspace,etoolbox}
\usepackage{titlesec,titletoc}
\usepackage{fancyhdr}
\usepackage{caption}
\usepackage[most]{tcolorbox}

% ---------- 颜色别名（章节里用的旧名继续有效） ----------
\colorlet{mainblue}{yhbAccent}
\colorlet{secondcyan}{yhbAccent2}
\colorlet{thirdteal}{yhbAccent2}
\colorlet{highlight}{yhbAccent!6!white}
\colorlet{tablehead}{yhbAccent!12!white}
\definecolor{textgray}{HTML}{4A5057}

% ---------- 课程标签 ----------
\providecommand{\YHBdoclabel}{课程笔记}
% ---------- 正文节奏 ----------
\linespread{1.3}
\setlength{\parindent}{0pt}
\setlength{\parskip}{6pt plus 1.5pt minus 1.5pt}
\setlength{\emergencystretch}{2.5em}
\raggedbottom
\widowpenalty=10000
\clubpenalty=10000
\setlength{\tabcolsep}{7pt}
\renewcommand{\arraystretch}{1.3}
\allowdisplaybreaks[2]
\AtBeginDocument{%
  \setlength{\abovedisplayskip}{10pt plus 2pt minus 2pt}%
  \setlength{\belowdisplayskip}{10pt plus 2pt minus 2pt}%
  \setlength{\abovedisplayshortskip}{6pt}%
  \setlength{\belowdisplayshortskip}{6pt}%
}
\captionsetup{font=small,labelfont=bf,skip=6pt}

% ---------- 列表 ----------
\setlist{leftmargin=2em,labelsep=0.6em,itemsep=4pt,parsep=2pt,topsep=6pt,partopsep=0pt}
\setlist[itemize,2]{label=\textbullet,leftmargin=1.6em}

% ---------- 标题层级 ----------
\ifdefined\YHBbooktrue
  \titleformat{\chapter}[display]
    {\sffamily\bfseries\color{yhbAccent}}
    {\normalsize\sffamily 第\,\thechapter\,章}
    {6pt}
    {\LARGE}
    [{\vspace{5pt}{\color{yhbAccent!55}\titlerule[1pt]}}]
  \titlespacing*{\chapter}{0pt}{14pt}{24pt}
  \titleformat{\section}{\Large\sffamily\bfseries\color{yhbAccent}}{\thesection}{0.7em}{}
  \titlespacing*{\section}{0pt}{18pt plus 3pt minus 2pt}{8pt}
\else
  \titleformat{\section}{\Large\sffamily\bfseries\color{yhbAccent}}{\thesection}{0.7em}{}[{\vspace{3pt}{\color{yhbAccent!45}\titlerule[0.7pt]}}]
  \titlespacing*{\section}{0pt}{18pt plus 3pt minus 2pt}{9pt}
\fi
\titleformat{\subsection}{\large\sffamily\bfseries\color{yhbAccent!85!black}}{\thesubsection}{0.7em}{}
\titleformat{\subsubsection}{\normalsize\bfseries\color{black!85}}{\thesubsubsection}{0.7em}{}
\titlespacing*{\subsection}{0pt}{12pt plus 2pt minus 1pt}{5pt}
\titlespacing*{\subsubsection}{0pt}{9pt plus 2pt minus 1pt}{4pt}
\pretocmd{\section}{\Needspace{6\baselineskip}}{}{}
\pretocmd{\subsection}{\Needspace{5\baselineskip}}{}{}
\pretocmd{\subsubsection}{\Needspace{4\baselineskip}}{}{}

\ifdefined\YHBbooktrue
  \setcounter{secnumdepth}{2}
  \setcounter{tocdepth}{1}
\else
  \setcounter{secnumdepth}{3}
  \setcounter{tocdepth}{2}
\fi
\renewcommand{\contentsname}{目录}
% 目录编号颜色由 titletoc 的 titlecontents 控制，不再重定义 \numberline

% ---------- 目录样式 ----------
\ifdefined\YHBbooktrue
  \titlecontents{chapter}[0pt]
    {\addvspace{6pt}\sffamily\bfseries\color{yhbAccent}}
    {\contentslabel{2.2em}}{}{\titlerule*[0.6em]{.}\contentspage}
  \titlecontents{section}[2.2em]
    {\addvspace{2pt}}
    {\contentslabel{2.2em}}{}{\titlerule*[0.6em]{.}\contentspage}
\else
  \titlecontents{section}[0pt]
    {\addvspace{5pt}\sffamily\bfseries\color{yhbAccent}}
    {\contentslabel{2.2em}}{}{\titlerule*[0.6em]{.}\contentspage}
  \titlecontents{subsection}[2.2em]
    {\addvspace{2pt}}
    {\contentslabel{2.4em}}{}{\titlerule*[0.6em]{.}\contentspage}
\fi

% ---------- 页眉页脚 ----------
\pagestyle{fancy}
\fancyhf{}
\fancyhead[L]{\footnotesize\sffamily\color{textgray}\YHBcourse\ \YHBdoclabel}
\ifdefined\YHBbooktrue
  \fancyhead[R]{\footnotesize\sffamily\color{textgray}\nouppercase{\leftmark}}
  \renewcommand{\chaptermark}[1]{\markboth{#1}{}}
\else
  \fancyhead[R]{\footnotesize\sffamily\color{textgray}\nouppercase{\rightmark}}
  \renewcommand{\sectionmark}[1]{\markright{#1}}
\fi
\fancyfoot[C]{\footnotesize\color{textgray}\thepage}
\renewcommand{\headrulewidth}{0.35pt}
\renewcommand{\headrule}{{\color{yhbAccent!45}\hrule width\headwidth height\headrulewidth}}
\fancypagestyle{plain}{%
  \fancyhf{}%
  \fancyfoot[C]{\footnotesize\color{textgray}\thepage}%
  \renewcommand{\headrulewidth}{0pt}%
}

% ---------- 公式与考点框 ----------
\newcommand{\keyformula}[1]{%
  \begin{tcolorbox}[enhanced,breakable,colback=highlight,colframe=yhbAccent!30!white,boxrule=0.5pt,arc=2.5pt,left=11pt,right=11pt,top=9pt,bottom=9pt,boxsep=0pt,before skip=11pt,after skip=11pt,halign=center,width=\linewidth]
  \adjustbox{max width=\linewidth-2em}{$\displaystyle #1$}
  \end{tcolorbox}}
\newcommand{\corepoint}[1]{%
  \begin{tcolorbox}[enhanced,breakable,colback=highlight,colframe=yhbAccent2,boxrule=0pt,borderline west={2.4pt}{0pt}{yhbAccent2},arc=0pt,left=11pt,right=11pt,top=8pt,bottom=8pt,before skip=11pt,after skip=11pt]
  \textbf{\color{yhbAccent}核心考点}\quad #1
  \end{tcolorbox}}
\newcommand{\symb}[2]{\textbf{#1}：#2\par}

% ---------- 超链接（务必最后加载） ----------
\usepackage[unicode=true,colorlinks=true,linkcolor=yhbAccent,citecolor=yhbAccent,urlcolor=yhbAccent!85!black,pdfborder={0 0 0},bookmarksnumbered=true,pdfstartview=FitH]{hyperref}
\usepackage{hypcap}
\hypersetup{pdftitle={\YHBcourse 课程笔记},pdfauthor={浮萍},pdfsubject={课程学习与复习笔记},pdfcreator={XeLaTeX}}

%   \newcommand{\YHBcourse}{微积分A}      % 课程名（页眉与书签用）
%   \newcommand{\YHBversion}{v1.2}        % 版本号（封面用）
%   \newcommand{\YHBquote}{……}            % 封面题句
%   \definecolor{yhbAccent}{HTML}{24465C} % 主色（深）
%   \definecolor{yhbAccent2}{HTML}{347A80}% 辅色（浅）
%   \YHBbooktrue / \YHBbookfalse          % 是否 book 文档类（有 \chapter）
% 章节文件只写内容，不要重复 \usepackage、不要改字号。
% ============================================================

% ---------- 页面与字体 ----------
\usepackage{geometry}
\geometry{a4paper,left=24mm,right=24mm,top=23mm,bottom=24mm,headheight=15pt,headsep=7mm,footskip=13mm}
\usepackage{microtype}
\usepackage{amsmath,amssymb,amsfonts,mathtools}
\usepackage{esint}
\usepackage{graphicx}
\usepackage{adjustbox}
\usepackage{array,booktabs,multirow,makecell,tabularx,longtable}
\usepackage{xcolor,colortbl}
\usepackage{enumitem}
\usepackage{needspace,etoolbox}
\usepackage{titlesec,titletoc}
\usepackage{fancyhdr}
\usepackage{caption}
\usepackage[most]{tcolorbox}

% ---------- 颜色别名（章节里用的旧名继续有效） ----------
\colorlet{mainblue}{yhbAccent}
\colorlet{secondcyan}{yhbAccent2}
\colorlet{thirdteal}{yhbAccent2}
\colorlet{highlight}{yhbAccent!6!white}
\colorlet{tablehead}{yhbAccent!12!white}
\definecolor{textgray}{HTML}{4A5057}

% ---------- 课程标签 ----------
\providecommand{\YHBdoclabel}{课程笔记}
% ---------- 正文节奏 ----------
\linespread{1.3}
\setlength{\parindent}{0pt}
\setlength{\parskip}{6pt plus 1.5pt minus 1.5pt}
\setlength{\emergencystretch}{2.5em}
\raggedbottom
\widowpenalty=10000
\clubpenalty=10000
\setlength{\tabcolsep}{7pt}
\renewcommand{\arraystretch}{1.3}
\allowdisplaybreaks[2]
\AtBeginDocument{%
  \setlength{\abovedisplayskip}{10pt plus 2pt minus 2pt}%
  \setlength{\belowdisplayskip}{10pt plus 2pt minus 2pt}%
  \setlength{\abovedisplayshortskip}{6pt}%
  \setlength{\belowdisplayshortskip}{6pt}%
}
\captionsetup{font=small,labelfont=bf,skip=6pt}

% ---------- 列表 ----------
\setlist{leftmargin=2em,labelsep=0.6em,itemsep=4pt,parsep=2pt,topsep=6pt,partopsep=0pt}
\setlist[itemize,2]{label=\textbullet,leftmargin=1.6em}

% ---------- 标题层级 ----------
\ifdefined\YHBbooktrue
  \titleformat{\chapter}[display]
    {\sffamily\bfseries\color{yhbAccent}}
    {\normalsize\sffamily 第\,\thechapter\,章}
    {6pt}
    {\LARGE}
    [{\vspace{5pt}{\color{yhbAccent!55}\titlerule[1pt]}}]
  \titlespacing*{\chapter}{0pt}{14pt}{24pt}
  \titleformat{\section}{\Large\sffamily\bfseries\color{yhbAccent}}{\thesection}{0.7em}{}
  \titlespacing*{\section}{0pt}{18pt plus 3pt minus 2pt}{8pt}
\else
  \titleformat{\section}{\Large\sffamily\bfseries\color{yhbAccent}}{\thesection}{0.7em}{}[{\vspace{3pt}{\color{yhbAccent!45}\titlerule[0.7pt]}}]
  \titlespacing*{\section}{0pt}{18pt plus 3pt minus 2pt}{9pt}
\fi
\titleformat{\subsection}{\large\sffamily\bfseries\color{yhbAccent!85!black}}{\thesubsection}{0.7em}{}
\titleformat{\subsubsection}{\normalsize\bfseries\color{black!85}}{\thesubsubsection}{0.7em}{}
\titlespacing*{\subsection}{0pt}{12pt plus 2pt minus 1pt}{5pt}
\titlespacing*{\subsubsection}{0pt}{9pt plus 2pt minus 1pt}{4pt}
\pretocmd{\section}{\Needspace{6\baselineskip}}{}{}
\pretocmd{\subsection}{\Needspace{5\baselineskip}}{}{}
\pretocmd{\subsubsection}{\Needspace{4\baselineskip}}{}{}

\ifdefined\YHBbooktrue
  \setcounter{secnumdepth}{2}
  \setcounter{tocdepth}{1}
\else
  \setcounter{secnumdepth}{3}
  \setcounter{tocdepth}{2}
\fi
\renewcommand{\contentsname}{目录}
% 目录编号颜色由 titletoc 的 titlecontents 控制，不再重定义 \numberline

% ---------- 目录样式 ----------
\ifdefined\YHBbooktrue
  \titlecontents{chapter}[0pt]
    {\addvspace{6pt}\sffamily\bfseries\color{yhbAccent}}
    {\contentslabel{2.2em}}{}{\titlerule*[0.6em]{.}\contentspage}
  \titlecontents{section}[2.2em]
    {\addvspace{2pt}}
    {\contentslabel{2.2em}}{}{\titlerule*[0.6em]{.}\contentspage}
\else
  \titlecontents{section}[0pt]
    {\addvspace{5pt}\sffamily\bfseries\color{yhbAccent}}
    {\contentslabel{2.2em}}{}{\titlerule*[0.6em]{.}\contentspage}
  \titlecontents{subsection}[2.2em]
    {\addvspace{2pt}}
    {\contentslabel{2.4em}}{}{\titlerule*[0.6em]{.}\contentspage}
\fi

% ---------- 页眉页脚 ----------
\pagestyle{fancy}
\fancyhf{}
\fancyhead[L]{\footnotesize\sffamily\color{textgray}\YHBcourse\ \YHBdoclabel}
\ifdefined\YHBbooktrue
  \fancyhead[R]{\footnotesize\sffamily\color{textgray}\nouppercase{\leftmark}}
  \renewcommand{\chaptermark}[1]{\markboth{#1}{}}
\else
  \fancyhead[R]{\footnotesize\sffamily\color{textgray}\nouppercase{\rightmark}}
  \renewcommand{\sectionmark}[1]{\markright{#1}}
\fi
\fancyfoot[C]{\footnotesize\color{textgray}\thepage}
\renewcommand{\headrulewidth}{0.35pt}
\renewcommand{\headrule}{{\color{yhbAccent!45}\hrule width\headwidth height\headrulewidth}}
\fancypagestyle{plain}{%
  \fancyhf{}%
  \fancyfoot[C]{\footnotesize\color{textgray}\thepage}%
  \renewcommand{\headrulewidth}{0pt}%
}

% ---------- 公式与考点框 ----------
\newcommand{\keyformula}[1]{%
  \begin{tcolorbox}[enhanced,breakable,colback=highlight,colframe=yhbAccent!30!white,boxrule=0.5pt,arc=2.5pt,left=11pt,right=11pt,top=9pt,bottom=9pt,boxsep=0pt,before skip=11pt,after skip=11pt,halign=center,width=\linewidth]
  \adjustbox{max width=\linewidth-2em}{$\displaystyle #1$}
  \end{tcolorbox}}
\newcommand{\corepoint}[1]{%
  \begin{tcolorbox}[enhanced,breakable,colback=highlight,colframe=yhbAccent2,boxrule=0pt,borderline west={2.4pt}{0pt}{yhbAccent2},arc=0pt,left=11pt,right=11pt,top=8pt,bottom=8pt,before skip=11pt,after skip=11pt]
  \textbf{\color{yhbAccent}核心考点}\quad #1
  \end{tcolorbox}}
\newcommand{\symb}[2]{\textbf{#1}：#2\par}

% ---------- 超链接（务必最后加载） ----------
\usepackage[unicode=true,colorlinks=true,linkcolor=yhbAccent,citecolor=yhbAccent,urlcolor=yhbAccent!85!black,pdfborder={0 0 0},bookmarksnumbered=true,pdfstartview=FitH]{hyperref}
\usepackage{hypcap}
\hypersetup{pdftitle={\YHBcourse 课程笔记},pdfauthor={浮萍},pdfsubject={课程学习与复习笔记},pdfcreator={XeLaTeX}}

%   \newcommand{\YHBcourse}{微积分A}      % 课程名（页眉与书签用）
%   \newcommand{\YHBversion}{v1.2}        % 版本号（封面用）
%   \newcommand{\YHBquote}{……}            % 封面题句
%   \definecolor{yhbAccent}{HTML}{24465C} % 主色（深）
%   \definecolor{yhbAccent2}{HTML}{347A80}% 辅色（浅）
%   \YHBbooktrue / \YHBbookfalse          % 是否 book 文档类（有 \chapter）
% 章节文件只写内容，不要重复 \usepackage、不要改字号。
% ============================================================

% ---------- 页面与字体 ----------
\usepackage{geometry}
\geometry{a4paper,left=24mm,right=24mm,top=23mm,bottom=24mm,headheight=15pt,headsep=7mm,footskip=13mm}
\usepackage{microtype}
\usepackage{amsmath,amssymb,amsfonts,mathtools}
\usepackage{esint}
\usepackage{graphicx}
\usepackage{adjustbox}
\usepackage{array,booktabs,multirow,makecell,tabularx,longtable}
\usepackage{xcolor,colortbl}
\usepackage{enumitem}
\usepackage{needspace,etoolbox}
\usepackage{titlesec,titletoc}
\usepackage{fancyhdr}
\usepackage{caption}
\usepackage[most]{tcolorbox}

% ---------- 颜色别名（章节里用的旧名继续有效） ----------
\colorlet{mainblue}{yhbAccent}
\colorlet{secondcyan}{yhbAccent2}
\colorlet{thirdteal}{yhbAccent2}
\colorlet{highlight}{yhbAccent!6!white}
\colorlet{tablehead}{yhbAccent!12!white}
\definecolor{textgray}{HTML}{4A5057}

% ---------- 课程标签 ----------
\providecommand{\YHBdoclabel}{课程笔记}
% ---------- 正文节奏 ----------
\linespread{1.3}
\setlength{\parindent}{0pt}
\setlength{\parskip}{6pt plus 1.5pt minus 1.5pt}
\setlength{\emergencystretch}{2.5em}
\raggedbottom
\widowpenalty=10000
\clubpenalty=10000
\setlength{\tabcolsep}{7pt}
\renewcommand{\arraystretch}{1.3}
\allowdisplaybreaks[2]
\AtBeginDocument{%
  \setlength{\abovedisplayskip}{10pt plus 2pt minus 2pt}%
  \setlength{\belowdisplayskip}{10pt plus 2pt minus 2pt}%
  \setlength{\abovedisplayshortskip}{6pt}%
  \setlength{\belowdisplayshortskip}{6pt}%
}
\captionsetup{font=small,labelfont=bf,skip=6pt}

% ---------- 列表 ----------
\setlist{leftmargin=2em,labelsep=0.6em,itemsep=4pt,parsep=2pt,topsep=6pt,partopsep=0pt}
\setlist[itemize,2]{label=\textbullet,leftmargin=1.6em}

% ---------- 标题层级 ----------
\ifdefined\YHBbooktrue
  \titleformat{\chapter}[display]
    {\sffamily\bfseries\color{yhbAccent}}
    {\normalsize\sffamily 第\,\thechapter\,章}
    {6pt}
    {\LARGE}
    [{\vspace{5pt}{\color{yhbAccent!55}\titlerule[1pt]}}]
  \titlespacing*{\chapter}{0pt}{14pt}{24pt}
  \titleformat{\section}{\Large\sffamily\bfseries\color{yhbAccent}}{\thesection}{0.7em}{}
  \titlespacing*{\section}{0pt}{18pt plus 3pt minus 2pt}{8pt}
\else
  \titleformat{\section}{\Large\sffamily\bfseries\color{yhbAccent}}{\thesection}{0.7em}{}[{\vspace{3pt}{\color{yhbAccent!45}\titlerule[0.7pt]}}]
  \titlespacing*{\section}{0pt}{18pt plus 3pt minus 2pt}{9pt}
\fi
\titleformat{\subsection}{\large\sffamily\bfseries\color{yhbAccent!85!black}}{\thesubsection}{0.7em}{}
\titleformat{\subsubsection}{\normalsize\bfseries\color{black!85}}{\thesubsubsection}{0.7em}{}
\titlespacing*{\subsection}{0pt}{12pt plus 2pt minus 1pt}{5pt}
\titlespacing*{\subsubsection}{0pt}{9pt plus 2pt minus 1pt}{4pt}
\pretocmd{\section}{\Needspace{6\baselineskip}}{}{}
\pretocmd{\subsection}{\Needspace{5\baselineskip}}{}{}
\pretocmd{\subsubsection}{\Needspace{4\baselineskip}}{}{}

\ifdefined\YHBbooktrue
  \setcounter{secnumdepth}{2}
  \setcounter{tocdepth}{1}
\else
  \setcounter{secnumdepth}{3}
  \setcounter{tocdepth}{2}
\fi
\renewcommand{\contentsname}{目录}
% 目录编号颜色由 titletoc 的 titlecontents 控制，不再重定义 \numberline

% ---------- 目录样式 ----------
\ifdefined\YHBbooktrue
  \titlecontents{chapter}[0pt]
    {\addvspace{6pt}\sffamily\bfseries\color{yhbAccent}}
    {\contentslabel{2.2em}}{}{\titlerule*[0.6em]{.}\contentspage}
  \titlecontents{section}[2.2em]
    {\addvspace{2pt}}
    {\contentslabel{2.2em}}{}{\titlerule*[0.6em]{.}\contentspage}
\else
  \titlecontents{section}[0pt]
    {\addvspace{5pt}\sffamily\bfseries\color{yhbAccent}}
    {\contentslabel{2.2em}}{}{\titlerule*[0.6em]{.}\contentspage}
  \titlecontents{subsection}[2.2em]
    {\addvspace{2pt}}
    {\contentslabel{2.4em}}{}{\titlerule*[0.6em]{.}\contentspage}
\fi

% ---------- 页眉页脚 ----------
\pagestyle{fancy}
\fancyhf{}
\fancyhead[L]{\footnotesize\sffamily\color{textgray}\YHBcourse\ \YHBdoclabel}
\ifdefined\YHBbooktrue
  \fancyhead[R]{\footnotesize\sffamily\color{textgray}\nouppercase{\leftmark}}
  \renewcommand{\chaptermark}[1]{\markboth{#1}{}}
\else
  \fancyhead[R]{\footnotesize\sffamily\color{textgray}\nouppercase{\rightmark}}
  \renewcommand{\sectionmark}[1]{\markright{#1}}
\fi
\fancyfoot[C]{\footnotesize\color{textgray}\thepage}
\renewcommand{\headrulewidth}{0.35pt}
\renewcommand{\headrule}{{\color{yhbAccent!45}\hrule width\headwidth height\headrulewidth}}
\fancypagestyle{plain}{%
  \fancyhf{}%
  \fancyfoot[C]{\footnotesize\color{textgray}\thepage}%
  \renewcommand{\headrulewidth}{0pt}%
}

% ---------- 公式与考点框 ----------
\newcommand{\keyformula}[1]{%
  \begin{tcolorbox}[enhanced,breakable,colback=highlight,colframe=yhbAccent!30!white,boxrule=0.5pt,arc=2.5pt,left=11pt,right=11pt,top=9pt,bottom=9pt,boxsep=0pt,before skip=11pt,after skip=11pt,halign=center,width=\linewidth]
  \adjustbox{max width=\linewidth-2em}{$\displaystyle #1$}
  \end{tcolorbox}}
\newcommand{\corepoint}[1]{%
  \begin{tcolorbox}[enhanced,breakable,colback=highlight,colframe=yhbAccent2,boxrule=0pt,borderline west={2.4pt}{0pt}{yhbAccent2},arc=0pt,left=11pt,right=11pt,top=8pt,bottom=8pt,before skip=11pt,after skip=11pt]
  \textbf{\color{yhbAccent}核心考点}\quad #1
  \end{tcolorbox}}
\newcommand{\symb}[2]{\textbf{#1}：#2\par}

% ---------- 超链接（务必最后加载） ----------
\usepackage[unicode=true,colorlinks=true,linkcolor=yhbAccent,citecolor=yhbAccent,urlcolor=yhbAccent!85!black,pdfborder={0 0 0},bookmarksnumbered=true,pdfstartview=FitH]{hyperref}
\usepackage{hypcap}
\hypersetup{pdftitle={\YHBcourse 课程笔记},pdfauthor={浮萍},pdfsubject={课程学习与复习笔记},pdfcreator={XeLaTeX}}

\usepackage{tikz}
\usepackage{circuitikz}
\xeCJKDeclareCharClass{CJK}{"2460->"24FF, "2160->"217F, "2190->"21FF, "2100->"214F}

\begin{document}
% ============================================================
% 2.cover.tex —— 统一封面（信息全部来自主文件定义）
% 需要的宏：\YHBcourse、\YHBversion、\YHBquote、yhbAccent、yhbAccent2
% 可选：\YHBcoverSize / \YHBcoverLead（课程名字号与行距）、\YHBcoversub（副标题）
% ============================================================
\providecommand{\YHBcoverSize}{38pt}
\providecommand{\YHBcoverLead}{46pt}
\providecommand{\YHBcoversub}{课程笔记}
\begin{titlepage}
\thispagestyle{empty}
\pdfbookmark[0]{封面}{cover}

\begin{tikzpicture}[remember picture,overlay]
  % 顶部深色块
  \fill[yhbAccent] (current page.north west) rectangle ([yshift=-132mm]current page.north east);
  % 分隔细条（辅色）
  \fill[yhbAccent2] ([yshift=-132mm]current page.north west) rectangle ([yshift=-135mm]current page.north east);
  % 极细横线，弱装饰，不压字
  \foreach \y in {20,34,48,62,76,90,104,118}{%
    \draw[white,opacity=0.08,line width=0.5pt]
      ([xshift=16mm,yshift=-\y mm]current page.north west) --
      ([xshift=-16mm,yshift=-\y mm]current page.north east);}
\end{tikzpicture}

\vspace*{32mm}
\noindent
{\xeCJKsetup{CJKecglue={}}\color{white}\sffamily\bfseries\fontsize{\YHBcoverSize}{\YHBcoverLead}\selectfont \YHBcourse\par}
\vspace{6mm}
{\color{white}\sffamily\fontsize{25}{32}\selectfont \YHBcoversub\par}
\vspace{9mm}
{\color{white}\fontsize{12.5}{18}\selectfont 作者\quad 浮萍\qquad 版本\quad \YHBversion\par}

\vspace{20mm}
{\sffamily\bfseries\fontsize{13}{20}\selectfont\color{yhbAccent}\YHBquote\par}

\vfill
{\color{yhbAccent}\rule{\textwidth}{1.4pt}}\par
\vspace{7mm}
{\fontsize{11.5}{17}\selectfont\sffamily
\textbf{QQ 邮箱}\quad\href{mailto:3071974740@qq.com}{3071974740@qq.com}\par
\vspace{2mm}
\textbf{GitHub}\quad\href{https://github.com/Duckweed-yhb}{github.com/Duckweed-yhb}\par}
\vspace{7mm}
{\fontsize{10.5}{16}\selectfont\color{black!78}
本笔记为个人学习整理，内容仍在持续完善。如发现内容、公式、文字或排版问题，欢迎通过上述 QQ 邮箱反馈；请尽量注明版本、页码与问题描述，便于核对和修订。\par}
\vspace*{12mm}
\end{titlepage}

\newpage
\tableofcontents
\newpage
% ===== 章节目录：在此按教材章节顺序 \input =====
% ============================================================
% 绪论（教材：戈宝军《电机学》第三版，绪论，印刷页 1--2）
% 教材中绪论不编章号，故用 \chapter* 保持"第X章"与教材一致
% ============================================================
\chapter*{绪论}
\markboth{绪论}{}
\phantomsection\addcontentsline{toc}{chapter}{绪论}

电机学研究的对象是\textbf{电机}——以电磁感应定律为基础、实现机械能与电能之间转换以及电能特性变换的机械装置。它包含两大类：做相对运动的\textbf{旋转电机}与不动的\textbf{变压器}。本课程的任务不是记住一堆结构图，而是建立一套"磁场—磁路—电路—能量"的分析框架，并用它去解释四类电机为什么这样运行、性能由什么决定、参数如何测定。

这门课的特点，教材在前言里说得直白：理论性强、概念多、电磁与时空关系复杂、与工程实际联系紧密。因此在读每一章时都要反复问三个问题：\textbf{磁场从哪来}（励磁）、\textbf{磁场怎么动}（相对运动或交变）、\textbf{能量怎么过去}（电磁功率与损耗）。这三个问题一旦答不上来，后面的等值电路就只是符号游戏。

% ------------------------------------------------------------
\section*{一、电机的概念}
\addcontentsline{toc}{section}{一、电机的概念}

电机是以电磁感应定律为基础，实现机械能与电能之间的转换，以及电能特性变换的机械装置，包括\textbf{旋转电机}和\textbf{变压器}两大类。

\begin{itemize}
  \item \textbf{旋转电机}是机电能量转换装置：作为\textbf{发电机}运行时把机械能转换成电能；作为\textbf{电动机}运行时把电能转换成机械能；有的电机还用作\textbf{调相机}，用于改善电网的功率因数；此外还有大量\textbf{微特电机}，广泛用于自动控制系统中。
  \item \textbf{变压器}是各部件间无相对运动的电能变换装置，广泛用于电能传输，以及电压、电流、阻抗的变换和电路隔离。
\end{itemize}

\begin{center}
\begin{tabular}{@{}lll@{}}
\toprule
装置 & 有无相对运动 & 能量/信号的作用 \\
\midrule
发电机 & 有 & 机械能 $\rightarrow$ 电能 \\
电动机 & 有 & 电能 $\rightarrow$ 机械能 \\
调相机 & 有 & 只交换无功，改善功率因数 \\
变压器 & 无 & 电压、电流、阻抗变换，电路隔离 \\
控制电机 & 有/无 & 控制系统中作执行、检测、放大、解算元件 \\
\bottomrule
\end{tabular}
\end{center}

\corepoint{发电机与电动机只是同一台电机的两种\textbf{运行方式}，电机本身是\textbf{可逆}的。判断某台电机处于发电机态还是电动机态，看的不是它的构造，而是电磁转矩与转速的方向关系、以及功率的流向（第\,12、16、20\,章会反复用到这条判据）。}

% ------------------------------------------------------------
\section*{二、电机的主要作用}
\addcontentsline{toc}{section}{二、电机的主要作用}

教材把电机的作用归为三个方面。

\subsection*{1. 电能的生产、传输和分配}
在发电厂，通过汽轮机、燃气轮机、柴油机、水轮机、风（轮）机等的驱动，把燃料燃烧和原子核裂变产生的能量，或水的势能、风的动能等转化为机械能，由\textbf{发电机}完成机电能量转换生产电能；然后经\textbf{变压器}升高电压，通过输电线路送到各用电地区，再经变压器降低电压供用户使用。

这一条链上，电机学要回答的是：为什么输电要升压（同样的功率下，电压升高则电流减小、线路损耗按电流平方下降）；为什么发电机出口电压受限于绝缘与制造（第\,16 章）；为什么变压器容量与铁心截面、绕组匝数直接相关（第\,3--4 章）。

\subsection*{2. 驱动各种生产机械和装备}
在工农业、交通运输、国防等部门和日常生活中，电动机广泛用于驱动生产机械、设备和器具：机床驱动、电力给排水、农副产品加工、矿石采掘和传送、电驱动车的牵引、鼓风设备、起重设备、轧钢机械、造纸设备、化工机械及家用电器的驱动等。全社会的用电量中，电动机（尤其是感应电动机）占绝大部分，这也是本课程把感应电机放在核心位置的原因。

\subsection*{3. 自动控制系统中的控制元件}
控制电机在控制系统、自动化系统中作为执行、检测、放大和解算元件。这类电机功率一般较小，但品种繁多、用途广泛（步进电机、伺服电机、测速发电机、旋转变压器等）。

% ------------------------------------------------------------
\section*{三、电机的分类}
\addcontentsline{toc}{section}{三、电机的分类}

电机种类很多，结构和用途各不相同，分类方法也很多。电机学中常采用以下三种分类方法。

\subsection*{（1）按应用的电流种类}
分为\textbf{直流电机}和\textbf{交流电机}。

\subsection*{（2）按能量转换功能和用途}
\begin{enumerate}
  \item 将机械能转换成电能的\textbf{发电机}；
  \item 将电能转换成机械能的\textbf{电动机}；
  \item 将一种交流电压等级改变成另一种交流电压等级的\textbf{变压器}；
  \item 在控制系统中起执行和测量作用的\textbf{控制电机}。
\end{enumerate}
另外还有\textbf{变流机}、\textbf{变频机}和\textbf{移相机}，分别用来改变电流、频率和相位。

\subsection*{（3）按运行转速的特点}
\begin{enumerate}
  \item \textbf{静止设备}的变压器；
  \item 没有固定同步转速的\textbf{直流电机}；
  \item 转子转速与同步转速有差异的\textbf{感应电机}；
  \item 转速恒为同步转速的\textbf{同步电机}。
\end{enumerate}

随着电力电子技术和电工材料的发展，出现了不属于上述传统类型的电机，如步进电机、无刷电机、超声波电机等，通常称为\textbf{特种电机}。本书按\textbf{工作原理}介绍电机，主要包括变压器、感应电机、同步电机和直流电机四大类电机的基本原理与应用。

\begin{center}
\begin{tabular}{@{}p{2.4cm}p{2.2cm}p{2.6cm}p{5.0cm}@{}}
\toprule
大类 & 电流种类 & 转速特点 & 本课程关注的核心问题 \\
\midrule
变压器 & 交流 & 静止 & 变比、等值电路、参数测定、联结组、并联运行 \\
感应电机 & 交流 & $n \neq n_1$（异步） & 转差率、等值电路、转矩—转差特性、起动与调速 \\
同步电机 & 交流 & $n = n_1$（同步） & 电枢反应、功角特性、并网、功率调节 \\
直流电机 & 直流 & 无固定同步转速 & 电枢反应、感应电动势与电磁转矩、机械特性 \\
\bottomrule
\end{tabular}
\end{center}

\corepoint{分类的目的不是背名词，而是确定"用哪一套方程"。看到"转速恒等于同步转速"就要想到同步电机方程（$E_0$ 与 $U$ 的关系、功角 $\delta$）；看到"转速随负载变化"就要想到感应电机的转差率 $s$；看到"直流"就要想到换向器把交流整成直流、以及 $E_a=C_e\Phi n$ 这条线性关系。}

% ------------------------------------------------------------
\section*{四、电机的发展}
\addcontentsline{toc}{section}{四、电机的发展}

英国物理学家法拉第（M. Faraday）于 1821 年发现了载流导体在磁场中会受到力的作用的现象，1831 年又发现了电磁感应定律，并很快就出现了原始模型电机。从此，电机的研究和应用迅速发展起来，至今已近 200 年。电机发展的历史大体上可以分成四个时期：

\begin{enumerate}
  \item \textbf{直流电机的产生和发展时期}；
  \item \textbf{交流电机的发展时期}；
  \item \textbf{电机理论、设计和制造工艺逐步达到完善化的时期}；
  \item \textbf{新技术、新材料等与电机相结合发展的新时期}。
\end{enumerate}

\begin{center}
\begin{tabular}{@{}llp{8.2cm}@{}}
\toprule
年份 & 人物/事件 & 意义 \\
\midrule
1821 & 法拉第发现载流导体受力 & 电动机原理的起点 \\
1831 & 法拉第发现电磁感应定律 & 发电机与变压器的共同理论基础 \\
1869 & 格拉姆（Z. T. Gramme）制成实用直流发电机 & 电力发展的重要推动 \\
1882 & 爱迪生（T. A. Edison）建成纽约曼哈顿直流照明系统 & 直流供电进入商业运行 \\
1883 & 特斯拉（N. Tesla）研制出第一台两相感应电机 & 交流电机进入实用阶段 \\
1888 & 多利沃-多勃罗夫斯基制成三相感应电机 & 三相制的基础 \\
20 世纪初 & 三相交流制在电力工业中占据统治地位 & 同步电机分高速（汽轮）与低速（水轮）两类 \\
\bottomrule
\end{tabular}
\end{center}

电机发展初期主要是直流电机发展的历史。随着直流电机的广泛应用，它很快暴露出固有缺点：不能解决远距离输电以及电压高低变换的问题；随着单机容量日益增大，换向也越来越困难，于是交流电机获得了迅速发展。19 世纪 90 年代初开始，三相同步电机按结构逐渐划分为高速和低速两类，高速的以汽轮发电机为代表，低速的以水轮发电机为代表。到 19 世纪末，直流电机、感应电机、同步电机、变压器等常规电机都已得到迅速发展和应用，相应的基本理论和设计方法也已初步建立。

20 世纪是电机发展的新时期。就单机容量而言，水轮发电机的最大单机容量在 20 世纪初不超过 1\,MW，现在已达到 1000\,MW；汽轮发电机的单机容量在 20 世纪初不超过 5\,MW，1930 年提高到 100\,MW，20 世纪 40—50 年代由于采用了氢冷、氢内冷、油冷和水冷等冷却方法，单机容量进一步提高，目前已超过 1000\,MW。由于新的永磁材料的出现，制成了许多高效节能、维护简单的永磁电机；由于电机和电力电子装置、单片机相组合，出现了各种性能和形态迥异的"一体化电机"。交流调速电机、风力发电机、核电电机和高效节能电机等方面的发展最为令人瞩目。

% ------------------------------------------------------------
\section*{五、本课程的任务}
\addcontentsline{toc}{section}{五、本课程的任务}

本课程是一门专业基础课。通过学习可获得电机的基本理论、基本知识和基本技能，为学习专业课做好准备，也为今后从事有关的专业工作打下基础。学习本课程后应达到下列基本要求。

\begin{enumerate}
  \item 对\textbf{磁路}的基本概念、基本定律、计算方法及交流铁芯线圈的励磁特性有基本的了解。
  \item 对变压器和三种主要电机（感应电机、同步电机和直流电机）的\textbf{基本结构}，交流三相单层、双层绕组和直流单叠、单波绕组的连接规律，要有一定认识；对各种电机中定、转子的\textbf{磁动势和气隙磁场}的性质及时间、空间关系，要有深入了解。
  \item 对变压器和三种主要电机正常\textbf{稳态运行}时的分析方法和运行性能，要牢固掌握：能正确地建立基本方程，弄清电机中的能量传递和能量转换关系；对稳态运行时的参数要有清晰的物理概念，能熟练运用\textbf{等值电路}计算变压器和交流电机的性能。
  \item 了解电机的有关\textbf{工程性问题}，如发热、冷却、温升、励磁系统和各种电机的应用范围；掌握电机的额定值、变压器和电机的主要性能数据及范围。
  \item 了解常用的特殊变压器、特殊感应电机和同步永磁电动机的基本结构、特点与应用。
  \item 通过\textbf{实验}，熟练掌握电机的基本实验方法和操作技能。
\end{enumerate}

\subsection*{学习方法的四条建议}
\begin{enumerate}
  \item \textbf{先建方程，再画图。}四类电机的分析套路几乎相同：写电压方程 $\rightarrow$ 写磁动势（磁通）方程 $\rightarrow$ 归算 $\rightarrow$ 画等值电路 $\rightarrow$ 画相量图 $\rightarrow$ 由等值电路求性能。顺序颠倒就会变成"背图"。
  \item \textbf{抓住"磁动势平衡"这一条主线。}变压器是一次、二次电流产生的磁动势共同建立主磁通；感应电机是定子、转子磁动势合成建立气隙磁场；同步电机是励磁磁动势与电枢磁动势合成。三者形式不同，实质相同（第\,4、13、17\,章）。
  \item \textbf{把"归算"当成一种坐标变换。}归算不改变功率和损耗，只改变电压、电流、阻抗的数值，目的是把两个不同电压等级（或不同频率的转子）接到同一个电路里（第\,4.3 节、第\,13.3 节）。
  \item \textbf{每个结论都要能说出前提。}如"$T \propto U^2$"只在忽略定子阻抗、转差率不变的近似下成立；"$U_1\approx E_1$"只在忽略一次漏阻抗时成立。写公式时把条件一并写上，考试与工程都不会吃亏。
\end{enumerate}

\corepoint{本课程的计算题几乎全部落在三件事上：\textbf{参数}（由空载、短路试验求 $R_m$、$X_m$、$R_k$、$X_k$ 等）、\textbf{特性}（由等值电路求 $I$、$P$、$T$、$\eta$、$\Delta U$）、\textbf{运行方式}（起动、调速、制动、并联、并网）。复习时按这三条把每章内容重新归类一遍，比按页码顺序重读有效得多。}

% ============================================================
% 第 1 章 磁路和磁性材料（教材第一篇，印刷页 4--20）
% 节名与教材一致：第一节 磁场中基本物理量和磁路的概念；第二节 磁路的基本定律；
% 第三节 磁性材料及其特性；第四节 简单磁路计算；第五节 交流磁路中的励磁电流和磁通
% ============================================================
\chapter{磁路和磁性材料}

以电磁感应定律为基础进行能量转换或传递的机电装置主要包括：将机械能转换为电能的发电机，将电能转换为机械能的电动机，以及实现电能传递的变压器。这些装置在能量传递或转换过程中均以\textbf{电磁场}为媒介，因此电磁场的强弱和分布不仅影响电机与变压器的性能，也决定着它们的体积和质量。

电磁场问题的求解一般有两种方法：一是\textbf{场的分析方法}，在给定区域用离散点物理量的数值定义其电磁特性，计算准确但计算难度与计算量较大，结构复杂时往往难以完成；二是\textbf{路的分析方法}，从宏观角度出发，把给定区域的场量与电路中的物理量作类比，用集总参数等效离散参数，把磁场问题简化成磁路问题。路的方法准确度低于场的方法，但计算简单，工程上足够用，因此本书主要采用路的分析方法。

\section{磁场中基本物理量和磁路的概念}

\subsection{磁场中常用物理量}

\subsubsection{磁感应强度}
当闭合导体中通入电流时，导体周围就会产生磁场，磁场的强弱和方向用\textbf{磁感应强度} $\boldsymbol{B}$ 来描述，其大小可由毕奥—萨伐尔定律求得。设导体中电流为 $i$，每一线元 $\mathrm{d}\boldsymbol{l}$ 在距其 $R$ 处的点 $P$ 产生的磁感应强度元为
\begin{equation}
\mathrm{d}\boldsymbol{B}=\frac{\mu_0}{4\pi}\cdot\frac{i\,\mathrm{d}\boldsymbol{l}\times \boldsymbol{a}_R}{R^{2}},
\end{equation}
式中 $\boldsymbol{a}_R$ 为由电流元指向点 $P$ 的单位矢量，$R$ 为从电流元到点 $P$ 的距离，$\mu_0$ 为真空磁导率。

$\boldsymbol{B}$ 的方向也可用闭合的磁感应矢量线（磁力线）描述：通电导体中的电流与所产生的磁场磁力线之间符合\textbf{右手螺旋定则}。$\boldsymbol{B}$ 是矢量，国际单位制中单位为 T（特斯拉）。

\subsubsection{磁通量}
$\boldsymbol{B}$ 描述的是空间每一点上磁场的大小和方向；要描述一个给定面上的磁场情况，需引入\textbf{磁通量}（简称磁通）$\Phi$。当磁感应强度在给定开表面 $A$ 上不均匀分布时
\begin{equation}
\Phi=\int_{A}\boldsymbol{B}\cdot \mathrm{d}\boldsymbol{A}.
\end{equation}
当 $\boldsymbol{B}$ 均匀分布且磁力线垂直穿过表面时
\begin{equation}
\Phi=BA.
\end{equation}

$\Phi$ 的单位为 Wb（韦伯），$A$ 的单位为 $\mathrm{m^2}$，故 $1\,\mathrm{T}=1\,\mathrm{Wb/m^2}$。工程上有时用 Gs（高斯）与 Mx（麦克斯韦）：
\begin{equation}
1\,\mathrm{Mx}=10^{-8}\,\mathrm{Wb},\qquad 1\,\mathrm{Gs}=10^{-4}\,\mathrm{T}.
\end{equation}
由 $\Phi=BA$ 可见，磁感应强度可看成穿过单位面积的磁通量，故又称\textbf{磁通密度}（磁密）。

\subsubsection{磁导率}
通电导体所产生磁场的强弱与导体周围介质的磁性能密切相关：有些介质使磁场显著增强，另一些则使其略有削弱。表征介质磁性能的物理量称为\textbf{磁导率} $\mu$，单位 H/m。真空磁导率
\begin{equation}
\mu_0=4\pi\times10^{-7}\ \mathrm{H/m},
\end{equation}
铁磁材料的磁导率 $\mu\gg\mu_0$。磁场计算时常采用\textbf{相对磁导率} $\mu_r=\mu/\mu_0$：各种硅钢片材料 $\mu_r$ 可达 5000、6000；铸钢约为 1000。

\subsubsection{磁场强度}
表征磁场性能的另一个基本物理量是\textbf{磁场强度} $\boldsymbol{H}$，定义为
\begin{equation}
\boldsymbol{H}=\frac{\boldsymbol{B}}{\mu}\qquad\text{或}\qquad \boldsymbol{B}=\mu\boldsymbol{H}.
\end{equation}
磁场强度反映了单位长度磁路上所消耗的磁动势，又称单位长度上的磁压降。$\boldsymbol{H}$ 也是矢量，单位为 A/m。

\subsection{磁路的概念}
电路中把电流流过的路径称为电路；与之相类比，研究磁场问题时把\textbf{磁通所经过的路径}称为\textbf{磁路}。电机和变压器中，为使磁通按设定路径通过，磁路一般由铁磁材料组成。

磁路可类比于电路，但两者有本质区别：
\begin{itemize}
  \item 电路中的电流由带电粒子的定向运动产生；磁路中\textbf{没有任何物质沿着闭合回路流动}。
  \item 物质世界存在电的导体与绝缘体，电流只在导体中流动；但\textbf{不存在磁的绝缘体}，磁通可以存在于任何介质之中。
\end{itemize}

由于铁磁材料磁导率远大于空气，绝大部分磁通在铁芯内通过——铁芯路径是设定的磁路，其中的磁通称为\textbf{主磁通}，其路径称为\textbf{主磁路}；铁芯外非铁磁材料磁导率不为零，总有少量磁通经铁芯外路径闭合，称为\textbf{漏磁通}，其路径称为\textbf{漏磁路}。电机和变压器中能量的转换与传递通过主磁通完成，故主磁通又称\textbf{工作磁通}。

产生磁通的载流线圈称为\textbf{励磁线圈}（励磁绕组），其中电流称为\textbf{励磁电流}。直流励磁电流产生的磁通不随时间变化，相应的磁路称为\textbf{直流磁路}；交流励磁电流产生的磁通随时间变化，称为\textbf{交流磁路}。

\begin{center}
\begin{tabular}{@{}p{2.4cm}p{4.9cm}p{5.6cm}@{}}
\toprule
对比项 & 电路 & 磁路 \\
\midrule
“流动”的物理量 & 电流 $I$（电荷定向运动） & 磁通 $\Phi$（无物质流动） \\
“源” & 电动势 $E$ & 磁动势 $F=Ni$ \\
“阻力” & 电阻 $R=l/(\gamma A)$ & 磁阻 $R_m=l/(\mu A)$ \\
介质界限 & 有导体与绝缘体之分 & 无磁绝缘体，任何介质均可通磁 \\
漏泄 & 漏电通常可忽略 & 漏磁严重，精确计算必须考虑 \\
线性 & 金属电阻基本为常数 & $\mu$ 随饱和程度变化，$R_m$ 非线性 \\
叠加原理 & 线性电路可用 & 非线性磁路\textbf{不可}用 \\
\bottomrule
\end{tabular}
\end{center}

\corepoint{磁路与电路只是"形式类比"，不是"物理同构"。最容易错的两处：① 磁路中 $\mu$ 不是常数，因此不能像电路那样对多个磁动势用叠加原理；② 漏磁不能忽略，尤其在气隙附近。做题时若发现要用叠加，先检查磁路是否处于线性段（$\mu_r$ 恒定）——只有线性化之后才可以。}

\section{磁路的基本定律}

\subsection{安培环路定律}
沿着任何一条闭合回路 $l$，磁场强度 $\boldsymbol{H}$ 的线积分等于该闭合回路所包围的总电流（代数和）：
\begin{equation}
\oint_{l}\boldsymbol{H}\cdot \mathrm{d}\boldsymbol{l}=\sum i .
\end{equation}
其中，若电流正方向与回路上磁场强度的方向符合右手螺旋关系，电流取正值，否则取负值。若沿回路 $l$ 磁场强度大小处处相等，且闭合回路所包围的总电流由通有电流 $i$ 的 $N$ 匝线圈提供，当 $Hl$ 的回路方向与电流方向符合右手螺旋定则时，可简写为
\begin{equation}
Hl=Ni .
\end{equation}

\subsection{磁路的欧姆定律}
设一等截面无分支铁芯磁路，铁芯截面积 $A$、磁路平均长度 $l$、材料磁导率 $\mu$，励磁绕组匝数 $N$、电流 $i$。若不计漏磁通，并认为各截面上磁密均匀分布且垂直于截面，则
\begin{equation}
\Phi=BA=\mu H A=\mu\frac{Ni}{l}A=\frac{Ni}{l/(\mu A)}=\frac{F}{R_m},
\end{equation}
其中
\begin{equation}
F=Ni,\qquad R_m=\frac{l}{\mu A},\qquad \Lambda=\frac{1}{R_m}.
\end{equation}

\keyformula{\Phi=\frac{F}{R_m}=\frac{Ni}{l/(\mu A)}}

式中：$F$ 为作用在铁芯磁路上的安匝数，称为\textbf{磁动势}，单位 A；$R_m$ 为\textbf{磁阻}，单位 A/Wb 或 $1/\mathrm{H}$；$\Lambda$ 为磁阻的倒数，称为\textbf{磁导}，单位 Wb/A 或 H。该式与电路欧姆定律形式相似，故称\textbf{磁路的欧姆定律}：$F$ 类比于 $E$，$\Phi$ 类比于 $I$，$R_m$ 类比于 $R$。

磁阻主要取决于磁路的几何尺寸和材料磁导率。由于铁磁材料 $\mu$ 不是常数，\textbf{铁磁材料的磁阻是非线性的}，这正是磁路与电路的本质区别之一。

\subsection{磁路的基尔霍夫第一定律}
若铁芯是带分支的磁路，不计漏磁时，任取一闭合面 $A$，由磁通连续性原理，穿过闭合面磁通的代数和必为零——进入该闭合面的磁通等于离开该闭合面的磁通：
\begin{equation}
\Phi_1=\Phi_2+\Phi_3
\quad\Longleftrightarrow\quad
-\Phi_1+\Phi_2+\Phi_3=0
\quad\Longleftrightarrow\quad
\sum\Phi=0 .
\end{equation}
规定穿出闭合面的磁通为正、进入为负，即得磁路的基尔霍夫第一定律。

\subsection{磁路的基尔霍夫第二定律}
磁路计算中通常把磁路分段：凡是材料和横截面积均相同、且磁通也相等的磁路作为一段。每一段内各点磁场强度相同。对由 $k$ 段组成的闭合磁路，由安培环路定律
\begin{equation}
H_1l_1+H_2l_2+H_3l_3=Ni
\quad\Longleftrightarrow\quad
\sum_{k}H_k l_k=Ni .
\end{equation}
式中，若 $H_k$ 的正方向与闭合回路方向一致，$H_kl_k$ 前取“$+$”号，否则取“$-$”号；若电流 $i$ 的正方向与 $l$ 的方向符合右手螺旋关系，$Ni$ 前取“$+$”号，否则取“$-$”号。

\keyformula{\sum_{k}H_k l_k=Ni}

$H_kl_k$ 称为第 $k$ 段磁路上的\textbf{磁位降}，$Ni$ 为作用在磁路上的\textbf{总磁动势}。该式表明：作用在任何闭合磁路的总磁动势恒等于各段磁路磁压降的代数和——即磁路的基尔霍夫第二定律，它实质上是安培环路定律的另一种表达形式。

\begin{center}
\begin{tabular}{@{}lll@{}}
\toprule
磁路定律 & 表达式 & 电路中的对应式 \\
\midrule
安培环路定律 & $\oint \boldsymbol{H}\cdot\mathrm{d}\boldsymbol{l}=\sum i$ & 全电流定律 \\
磁路欧姆定律 & $\Phi=F/R_m$，$R_m=l/(\mu A)$ & $I=E/R$ \\
基尔霍夫第一定律 & $\sum\Phi=0$ & $\sum I=0$ \\
基尔霍夫第二定律 & $\sum H_k l_k=Ni$ & $\sum U=\sum E$ \\
\bottomrule
\end{tabular}
\end{center}

\corepoint{用第二定律时符号最易错。稳妥做法：先任意规定回路绕行正方向，再逐段判断 $H_k$ 与绕行方向是否一致、$Ni$ 与绕行方向是否成右螺旋，同号取正、异号取负，最后解代数方程。切忌"全取正号"硬算。}

\section{磁性材料及其特性}

磁性材料按磁性能分为\textbf{顺磁材料}、\textbf{反磁材料}和\textbf{铁磁材料}。顺磁材料 $\mu$ 略大于 $\mu_0$，反磁材料 $\mu$ 略小于 $\mu_0$；工程上二者的磁导率一般按等于 $\mu_0$ 计算，如空气、变压器油、铜和铝等。铁磁材料 $\mu\gg\mu_0$，一般大到几千倍，如铁、镍、钴及其合金。电机和变压器的磁路通常都由铁磁材料制成。

\subsection{铁磁材料的磁化}
铁磁材料放入磁场中会呈现很强的磁性，这种现象称为\textbf{磁化}，能被磁化是铁磁材料的基本特性之一。普遍认为，铁磁材料内部存在许多很小的、被称为\textbf{磁畴}的天然磁化区，每个磁畴可看做一个微型磁针。无外磁场时各磁畴任意排列，材料不显磁性；放入外磁场后，磁畴受磁力作用发生旋转并沿磁场方向排列整齐，材料显示较强磁性。由此形成的磁化磁场与外磁场叠加，使合成磁场显著增强。非铁磁材料中不含磁畴，故在同样磁场强度下磁场不会显著增强——这就是铁磁材料磁导率大得多的原因。

\subsection{铁磁材料的磁化曲线}

\subsubsection{原始磁化曲线}
将未磁化的铁磁材料磁化，$H$ 由零逐渐增加时 $B$ 随之增加，所得 $B=f(H)$ 曲线称为\textbf{原始磁化曲线}。其形状分四段：起始段（$B$ 随 $H$ 增加较慢）、直线段（$B$ 迅速增长）、饱和段（增长放慢并逐渐趋于饱和）、过饱和段（曲线基本成为与非铁磁材料 $B=\mu_0H$ 平行的直线）。$b$ 点称\textbf{膝点}，$c$ 点称\textbf{饱和点}。

由于磁化曲线不是直线，$\mu=B/H$ 也随 $H$ 变化：直线段内 $\mu$ 随 $H$ 增加迅速增大；进入饱和区后 $\mu$ 急剧下降；至过饱和区 $\mu$ 基本不变，保持为常数。因此\textbf{铁磁材料的磁阻不是常数，而随饱和程度增加而增大}。

在电机和变压器的磁路设计中，为了用较少的磁动势获得较大的磁通，通常把铁磁材料的工作点选在\textbf{膝点附近}。

\subsubsection{磁滞回线}
将铁磁材料在 $-H_m$ 到 $+H_m$ 之间反复周期性磁化，得到对称于原点的闭合曲线，即\textbf{磁滞回线}。当 $H$ 从零增到 $+H_m$，$B$ 增到 $B_m$；再逐渐减小 $H$，$B$ 沿曲线下降，当 $H=0$ 时 $B$ 并不为零，而是等于 $B_r$——去掉外加磁场后材料内部仍保留一定磁感应强度的现象称为\textbf{顽磁性}，$B_r$ 称为\textbf{剩余磁感应强度}（剩磁）。要使 $B$ 从 $B_r$ 减小到零，必须加反向磁场，该反向磁场强度称为\textbf{矫顽力} $H_c$。$B_r$ 与 $H_c$ 是铁磁材料的两个重要参数。

$B$ 的变化总是落后于 $H$ 的变化，这种现象称为\textbf{磁滞}，相应的 $B$—$H$ 闭合曲线称为磁滞回线。磁滞现象是铁磁材料的另一个基本特性。

\subsubsection{基本磁化曲线}
对同一材料用不同的 $H_m$ 周期性磁化，可得一系列大小不同的磁滞回线；将各磁滞回线的\textbf{顶点}连接起来所得的曲线称为\textbf{基本磁化曲线}。工程上进行磁路计算时采用的都是基本磁化曲线。

\subsection{铁磁材料的分类}
按磁滞回线形状，铁磁材料分为两类：
\begin{itemize}
  \item \textbf{软磁材料}：磁滞回线窄、剩磁和矫顽力都小，如铸铁、铸钢、硅钢片和坡莫合金等。磁滞损耗小、磁导率高，常用来制造变压器和电机的铁芯。
  \item \textbf{硬磁材料}：磁滞回线宽、剩磁和矫顽力都大，如铝镍钴合金和稀土合金等。因剩磁大，常用来制造永久磁铁，故又称\textbf{永磁材料}。
\end{itemize}

\subsection{铁芯损耗}
\textbf{铁芯损耗}（铁耗）是铁芯磁路中由于磁滞和涡流引起的功率损耗的总称，\textbf{只存在于交流磁路中}，直流磁路没有这项功率损耗。

\subsubsection{磁滞损耗}
外磁场作用下，磁畴为保持磁场方向与外磁场一致会发生转动；外磁场交变时磁畴来回翻转、彼此摩擦而引起损耗，称为\textbf{磁滞损耗}。实验证明，磁滞损耗正比于交变磁场频率 $f$、磁滞回线面积和铁磁材料体积 $V$，而磁滞回线面积又正比于最大磁感应强度 $B_m$ 的 $n$ 次方（$n$ 取决于材料，一般硅钢片 $n=1.6\sim2.3$）：
\begin{equation}
P_h=C_h f B_m^{n} V .
\end{equation}

\subsubsection{涡流损耗}
磁通交变时，由电磁感应定律，铁芯内将产生感应电动势和电流，这些电流在铁芯内部围绕磁通形成环流，称为\textbf{涡流}；涡流在铁芯中引起的电阻损耗称为\textbf{涡流损耗}：
\begin{equation}
P_e=C_e f^{2}B_m^{2}\delta^{2} V ,
\end{equation}
式中 $C_e$ 为涡流损耗系数，反比于材料电阻率；$\delta$ 为硅钢片厚度。该式表明：为减小涡流损耗，可在钢材中加入少量硅以增加电阻率，并用相互绝缘的许多薄硅钢片叠成铁芯代替整块铁芯——变压器和电机的铁芯就是采用硅钢片制作的。

\subsubsection{铁芯损耗}
将磁滞损耗与涡流损耗相加即得铁芯损耗：
\begin{equation}
P_{Fe}=P_h+P_e=\left(C_h f B_m^{n}+C_e f^{2}B_m^{2}\delta^{2}\right)V .
\end{equation}
对于一般硅钢片，在正常工作磁感应强度范围内（$1\,\mathrm{T}<B_m<1.8\,\mathrm{T}$），可近似写成
\begin{equation}
P_{Fe}\approx C_{Fe}f^{1.3}B_m^{2}G ,
\end{equation}
式中 $C_{Fe}$ 为铁芯损耗系数，$G$ 为铁芯质量。

\keyformula{P_{Fe}\approx C_{Fe}f^{1.3}B_m^{2}G}

该式表明：铁芯损耗近似与频率的 1.3 次方成正比，与铁芯中最大磁感应强度的二次方及铁芯质量成正比。

\begin{center}
\begin{tabular}{@{}llll@{}}
\toprule
损耗 & 成因 & 表达式 & 减小措施 \\
\midrule
磁滞损耗 & 磁畴反复翻转、相互摩擦 & $P_h=C_hfB_m^{n}V$ & 选软磁材料（硅钢片、坡莫合金） \\
涡流损耗 & 交变磁通在铁芯内感生环流 & $P_e=C_ef^{2}B_m^{2}\delta^{2}V$ & 加硅提高电阻率、薄片叠装、片间绝缘 \\
铁芯损耗 & 以上两者之和 & $P_{Fe}\approx C_{Fe}f^{1.3}B_m^{2}G$ & 降低 $B_m$、减薄片厚、提高电阻率 \\
\bottomrule
\end{tabular}
\end{center}

\corepoint{铁耗 $\propto B_m^2$、$\propto f^{1.3}$，这条比例关系贯穿全书：变压器空载损耗近似等于铁耗（与负载无关），而铜耗随负载电流平方变化；感应电机里铁耗基本为常数、被归入空载损耗。凡遇到"某量随电压/频率如何变化"的判断，先想 $B_m\propto U/f$ 这一层关系。}

\section{简单磁路计算}

磁路计算问题一般分为两类。

\begin{enumerate}
  \item \textbf{给定磁通求磁动势}（正问题）。磁路材料和尺寸给定，要求产生给定磁通所需的励磁磁动势。由于铁磁材料非线性，必须先用 $\Phi$ 与各段有效面积求出各段 $B$，再对每个 $B$ 在基本磁化曲线上查出相应 $H$，最后由 $\sum H_kl_k=Ni$ 求出 $F$ 与励磁电流。
  \item \textbf{给定磁动势求磁通}（反问题）。由于 $\mu\neq$ 常数，通常先\textbf{假设}磁通，按正问题算出所需磁动势，与给定值比较后修正假设，用\textbf{迭代法}逐步逼近。
\end{enumerate}

\subsection{计算要点与步骤}
\begin{enumerate}
  \item \textbf{分段}。按"材料相同、截面积相同且磁通相等"把磁路分成若干段；有气隙时气隙单独成段。
  \item \textbf{求各段中心线长度 $l_k$ 与有效截面积 $A_k$}。硅钢片叠装时须乘\textbf{叠片系数}（除去片间绝缘厚度），即 $A=A_{\text{几何}}\times K_{Fe}$。
  \item \textbf{气隙有效面积}。考虑气隙磁场的边缘效应，计算气隙有效面积时通常在长、宽方向各增加一个气隙长度，即 $A_\delta=(a+\delta)(b+\delta)$。
  \item \textbf{求各段磁密} $B_k=\Phi/A_k$。
  \item \textbf{求各段磁场强度} $H_k$：线性段可用 $H_k=B_k/(\mu_r\mu_0)$；非线性段须查基本磁化曲线。
  \item \textbf{求磁压降与磁动势}：$F_k=H_kl_k$（或 $F_k=R_{mk}\Phi$），$F=\sum F_k$，励磁电流 $I=F/N$。
\end{enumerate}

\subsection{两种算式的选用}
\begin{itemize}
  \item \textbf{$F=Hl$ 法}：适用于材料磁导率非线性、需查磁化曲线的场合，是工程计算的主方法。
  \item \textbf{$F=R_m\Phi$ 法}：只在 $\mu$ 为常数的线性磁路（如气隙、或明确给出 $\mu_r$ 恒定的铁芯）中方便。
\end{itemize}

一个必须记住的结论：\textbf{气隙虽短，磁压降却占绝大部分}。教材算例中气隙仅占磁路平均长度的 0.36\%，其磁阻与磁压降却占整个磁路的 92\%；即使铁芯部分已进入饱和、该段磁阻增大，气隙磁压降仍占励磁磁动势的一半以上。

\corepoint{三处高频错误：① 忘记乘叠片系数，导致 $B$ 偏小、$H$ 查表偏小、$F$ 偏小；② 气隙面积直接取几何面积，忘了长宽各加一个 $\delta$，使气隙磁阻偏大；③ 铁芯进入饱和后仍用常数 $\mu_r$ 计算，结果严重偏小（饱和段 $\mu$ 急剧下降，$H$ 会跳升一到两个数量级）。}

\section{交流磁路中的励磁电流和磁通}

交流铁芯磁路中，由于铁芯磁化曲线的非线性，励磁电流 $i_m$ 与主磁通的波形会产生\textbf{畸变}。以单一铁芯材料的交流磁路为例，为便于图解，把磁化特性写成 $\Phi=f(i)$ 的形式。

\subsection{磁化特性为基本磁化曲线}
磁路内磁通较小时，磁路不饱和，$\Phi$ 与 $i$ 基本为线性关系：若磁通随时间按正弦规律变化，励磁电流也按正弦规律变化，两者\textbf{同相位}，波形相同。

磁通较大时磁路饱和，$\Phi$ 与 $i$ 呈非线性，需用图解法分析，分两种情况。

\subsubsection{磁通为正弦波时励磁电流的波形}
设磁通随时间作正弦变化，取若干时刻的磁通瞬时值，由 $\Phi=f(i)$ 曲线查出产生该磁通所需的励磁电流，逐点描出即得 $i_m=f(t)$。可见：磁通正弦且磁路饱和时，由于非线性，\textbf{励磁电流波形发生畸变，成为尖顶波}。分解后除基波外还含其他奇次谐波，其中\textbf{3 次谐波最大}。磁路越饱和，尖顶越严重，谐波越显著。但无论尖顶多严重，\textbf{其基波相位始终与磁通同相}。

\subsubsection{励磁电流为正弦波时磁通的波形}
当 $i_m$ 随时间正弦变化且磁路饱和时，磁通波形也发生畸变，成为\textbf{平顶波}；分解后除基波外亦含奇次谐波，同样以 3 次谐波最大。磁路越饱和，平顶越严重。但无论平顶多严重，\textbf{其基波相位始终与励磁电流同相}。

综上：励磁电流与磁通波形是否畸变，决定于磁路是否饱和。不饱和时磁化曲线基本线性，两者波形非常接近；饱和后，磁通为正弦时电流畸变成尖顶波，电流为正弦时磁通畸变成平顶波。无论尖顶波还是平顶波，分解后的主要成分都是\textbf{基波和 3 次谐波}；磁通基波与励磁电流基波的相位始终相同，与是否饱和及饱和程度无关。

\subsection{磁化特性为磁滞回线}
考虑磁滞现象时，磁化特性为磁滞回线。当磁通为正弦波时，通过磁滞回线用图解法求得 $i_m=f(t)$，它为非正弦波，可分解为两个分量：
\begin{equation}
i_m=i_\mu+i_{Fe},\qquad i_{Fe}=i_h+i_e ,
\end{equation}
其中 $i_\mu$ 为与磁通同相位的尖顶波，与忽略磁滞、只考虑饱和时得到的尖顶波完全一致，即不考虑磁滞时的励磁电流，与基本磁化曲线相对应，称为\textbf{磁化电流}；$i_{Fe}$ 波形近似正弦，相位\textbf{超前磁通 $90^\circ$}，是反映磁滞损耗和涡流损耗的\textbf{有功电流}。其中 $i_h$ 是考虑了磁滞现象才存在的\textbf{磁滞损耗电流}（不计磁滞时为零），$i_e$ 是\textbf{涡流损耗电流}，两者合称\textbf{铁耗电流}。

\subsection{三点结论}
\begin{enumerate}
  \item 激励产生磁通的电流是\textbf{磁化电流} $i_\mu$，其基波与磁通同相位；由于磁路非线性，$i_\mu$ 与磁通波形不同：磁通为正弦时 $i_\mu$ 为尖顶波，电流为正弦时磁通为平顶波。
  \item 铁芯中存在磁滞与涡流损耗，励磁电流除磁化电流外还存在铁芯损耗电流，后者接近正弦，相位\textbf{超前磁通 $90^\circ$}。
  \item 励磁电流是磁化电流与铁芯损耗电流的合成，即 $i_m=i_{Fe}+i_\mu$。变压器与电机铁芯损耗很小，$i_{Fe}$ 通常仅为 $i_m$ 的 10\% 左右，故 $i_m$ 主要由 $i_\mu$ 决定。
\end{enumerate}

\corepoint{三条常考结论要能脱口而出：① 磁通正弦 $\Rightarrow$ 励磁电流尖顶波；② 励磁电流正弦 $\Rightarrow$ 磁通平顶波；③ 无论哪种情况，\textbf{磁通基波与励磁电流基波同相位}，而铁耗电流超前磁通 $90^\circ$。第 ④ 章变压器的相量图里，$I_0$ 与 $\dot\Phi_m$ 的夹角 $\alpha$（铁耗角）就是从这里来的。}

\begin{center}
\begin{tabular}{@{}llll@{}}
\toprule
磁路状态 & 给定波形 & 另一量波形 & 主要谐波 \\
\midrule
不饱和 & 磁通正弦 & 电流正弦 & 无（近正弦） \\
饱和 & 磁通正弦 & 电流尖顶波 & 3 次最大 \\
饱和 & 电流正弦 & 磁通平顶波 & 3 次最大 \\
计及磁滞 & 磁通正弦 & 电流非正弦（含 $i_\mu+i_{Fe}$） & 3 次最大；$i_{Fe}$ 超前 $90^\circ$ \\
\bottomrule
\end{tabular}
\end{center}

% ============================================================
% 第 2 章 能量转换基本定律（教材第一篇，印刷页 21--23）
% 节名：第一节 电磁感应定律；第二节 毕奥—萨伐尔电磁力定律；第三节 能量守恒原理
% ============================================================
\chapter{能量转换基本定律}

各种机电能量转换装置的用途和结构虽有差别，其基本原理是相同的：\textbf{载流导体在磁场中要受到力的作用}是电动机工作的基本原理；\textbf{变化的磁场在导体中要感应电动势}是发电机工作的基本原理。本章阐述机电能量转换过程中的基本定律与能量关系。后续各章的分析，全部建立在这三条定律之上。

\section{电磁感应定律}

电磁感应定律描述"磁变生电"的电磁感应现象，1831 年由法拉第提出，又称法拉第电磁感应定律。它指出：交变的磁场会产生电场，并在导体中感应电动势。

\subsection{变压器电动势（磁变生电）}
当规定感应电动势正方向与磁通的正方向符合右手螺旋关系时，感应电动势与磁通对时间的变化率的负值成正比：
\begin{equation}
e=-N\frac{\mathrm{d}\Phi}{\mathrm{d}t}=-\frac{\mathrm{d}\Psi}{\mathrm{d}t},
\end{equation}
式中 $e$ 为回路中的感应电动势，$N$ 为绕组匝数，$\Phi$ 为通过该回路面积的磁通，$\Psi=N\Phi$ 为交链回路的\textbf{磁链}（$N$ 匝线圈所匝链的总磁链）。

\keyformula{e=-N\frac{\mathrm{d}\Phi}{\mathrm{d}t}=-\frac{\mathrm{d}\Psi}{\mathrm{d}t}}

\subsubsection{推导要点}
\begin{enumerate}
  \item 磁链 $\Psi=N\Phi$ 是"把 $N$ 个线匝的磁通串起来"的总量；对集中绕制的线圈，各匝磁通相同，故 $e=-N\mathrm{d}\Phi/\mathrm{d}t$。
  \item 负号是\textbf{楞次定律}的数学表达：感应电动势产生的电流总是阻碍磁通的变化。工程计算中常先按右手螺旋规定正方向、写出大小，再由楞次定律判断实际方向。
  \item 若磁通按正弦变化 $\Phi=\Phi_m\sin\omega t$，则
  \begin{equation}
  e=-N\Phi_m\omega\cos\omega t=E_m\sin\!\left(\omega t-90^\circ\right),
  \end{equation}
  即感应电动势\textbf{滞后磁通 $90^\circ$}，其有效值为 $E=E_m/\sqrt2=2\pi fN\Phi_m/\sqrt2=4.44fN\Phi_m$。这一"滞后 $90^\circ$ + $4.44$ 系数"是第 4 章变压器电动势公式的来源，务必在此处打牢。
\end{enumerate}

\subsection{切割电动势（动生生电）}
磁场恒定而导体在磁场中运动时，导体中的感应电动势为
\begin{equation}
e=\int_{l}\left(\boldsymbol{v}\times\boldsymbol{B}\right)\cdot \mathrm{d}\boldsymbol{l},
\end{equation}
若磁场方向、导体长度方向和导体运动方向三者相互垂直，则简化为
\begin{equation}
e=Blv .
\end{equation}

\keyformula{e=Blv}

式中 $B$ 为磁感应强度，$l$ 为导体长度，$v$ 为导体运动速度。感应电动势的方向可用\textbf{右手定则}确定（磁力线穿入手心，拇指指向导体运动方向，四指指向感应电动势方向）。

\subsubsection{两种电动势的统一}
切割电动势与变压器电动势不是两种不同的物理规律，而是同一法拉第定律在不同参考系下的表现：导体切割磁力线等价于导体所交链的磁通发生变化。因此在电机分析中，定子绕组看到的是"磁通交变"（变压器电动势），转子绕组看到的是"切割磁通"（速度电动势），两者形式不同、本质相同。

\corepoint{三个高频混淆点：① 变压器电动势 $e=-N\mathrm{d}\Phi/\mathrm{d}t$ 与切割电动势 $e=Blv$ 的正方向规定不同（前者用右手螺旋，后者用右手定则），混用会得出"电动势反了"的错误结论；② $e=Blv$ 中 $v$ 是导体\textbf{相对磁场}的速度，不是绝对速度；③ 三相绕组中 $e$ 的相位差来自空间位置差，而不是时间差，这条在第 9 章导体电动势星形图里会用到。}

\section{毕奥—萨伐尔电磁力定律}

磁场最基本的特性是对场域中的载流导体有力的作用。电磁力定律是描述电与磁之间相互作用产生力的基本定律。

取有效长度为 $l$ 的导体，其中流过电流 $i$，所处磁场为 $\boldsymbol{B}$，则导体所受的作用力为
\begin{equation}
\boldsymbol{f}=\int_{l}i\,\mathrm{d}\boldsymbol{l}\times\boldsymbol{B}.
\end{equation}
若磁场均匀且与导体相互垂直，则
\begin{equation}
f=Bli .
\end{equation}

\keyformula{f=Bli}

电磁力的方向可用\textbf{左手定则}判定（磁力线穿入手心，四指指向电流方向，拇指指向电磁力方向）。

\subsection{电磁力与电磁转矩}
电机中导体沿圆周分布，电磁力对转轴形成转矩。设导体所在处的磁密为 $B$、导体电流为 $i$、导体有效长度为 $l$、转子半径为 $R$，则该导体产生的转矩为 $TR=BliR$；若共有 $N$ 根导体、每极磁通为 $\Phi$、极对数为 $p$，经归并可得电磁转矩与电磁功率的关系
\begin{equation}
T=\frac{P_{em}}{\Omega},
\end{equation}
其中 $\Omega=2\pi n/60$ 为机械角速度。这一关系在第 12、16、20 章分别以不同形式出现（$T=\dfrac{P_{em}}{\Omega_1}$、$T=\dfrac{P_{em}}{\Omega_s}$ 等），其本源就是本节。

\subsection{电磁功率的恒等性}
把两条定律放在一起看，会得到一个非常重要的结论：导体既感应电动势 $e=Blv$、又受电磁力 $f=Bli$，于是
\begin{equation}
ei=\left(Blv\right)i=\left(Bli\right)v=fv .
\end{equation}
左边是\textbf{电功率}，右边是\textbf{机械功率}。这说明：电机中通过磁场耦合的电功率与机械功率在数值上相等，磁场只起"媒介"作用，不消耗能量。这就是机电能量转换的物理基础，也是"电磁功率"这一概念的由来。

\corepoint{$ei=fv$ 这条恒等式是全书最有价值的一句话：变压器只有 $ei$（没有 $v$），故只传递电能；旋转电机同时有 $ei$ 与 $fv$，故能实现机电能量转换。判断某台电机是发电机还是电动机，看的就是 $ei$ 与 $fv$ 谁输入、谁输出。}

\section{能量守恒原理}

能量守恒原理是物理学的基本原理，可表述为：在质量不变的物理系统内，能量是守恒的——能量不会凭空产生，亦不会凭空消灭，仅能变换其存在形式。

\subsection{电机中的四种能量形式}
电机中能量转换包括四种能量形式：\textbf{电能}、\textbf{机械能}、\textbf{磁场储能}和\textbf{热能}。

\subsubsection{电动机的能量关系}
\begin{equation}
\text{电源输入的电能}=\text{磁场储能的增量}+\text{转换为热能的能量损耗}+\text{输出的机械能}.
\end{equation}

\subsubsection{发电机的能量关系}
\begin{equation}
\text{输入的机械能}=\text{磁场储能的增量}+\text{转换为热能的能量损耗}+\text{输出的电能}.
\end{equation}

稳态运行时磁场储能不变（增量项为零），上面两式就退化为熟悉的功率平衡方程；瞬态过程中磁场储能参与吞吐，这一项不可忽略（本书主要讨论稳态，故通常略去）。

\subsection{能量损耗的三类}
转换为热能的能量主要包括三部分：
\begin{enumerate}
  \item \textbf{电气损耗}（铜耗）：导体中的电流在电阻中引起的损耗，$P_{Cu}=I^2R$；
  \item \textbf{机械损耗}：消耗于轴承摩擦及通风中的损耗；
  \item \textbf{铁芯损耗}（铁耗）：电机内磁场在铁芯中产生的磁滞损耗和涡流损耗。
\end{enumerate}
这些能量均转换为热能，并且是\textbf{不可逆}的——这部分能量很难或根本不可能再转换为电能或机械能。

\begin{center}
\begin{tabular}{@{}llll@{}}
\toprule
损耗类别 & 与什么有关 & 典型表达式 & 在等效电路中的位置 \\
\midrule
铜耗（电气损耗） & 电流平方 & $P_{Cu}=mI^2R$ & 串联电阻 $R_1$、$R_2'$ \\
铁耗 & 磁密平方、频率 1.3 次方 & $P_{Fe}\approx C_{Fe}f^{1.3}B_m^2G$ & 励磁支路中的等效电阻 \\
机械损耗 & 转速（摩擦、通风） & 由实验测定 & 一般不进入电路，从 $P_{em}$ 中扣除 \\
杂散损耗 & 漏磁、齿槽、谐波 & 按额定功率的百分数估算 & 归入附加损耗 \\
\bottomrule
\end{tabular}
\end{center}

\subsection{效率}
由能量守恒，效率定义为
\begin{equation}
\eta=\frac{P_2}{P_1}=\frac{P_1-\sum p}{P_1}\times100\%,
\end{equation}
其中 $P_1$ 为输入功率，$P_2$ 为输出功率，$\sum p$ 为总损耗。电机效率在负载率为额定值的 70\%—80\% 附近最高，因为此时不变损耗（铁耗、机械损耗）与可变损耗（铜耗）大致相等——这一"最大效率条件"在第 6 章变压器效率分析中会严格证明。

\subsection{能量转换的可逆性与条件}
\begin{itemize}
  \item \textbf{可逆性}：同一台电机既可作发电机、也可作电动机，取决于外部条件（原动机驱动还是电网供电）。
  \item \textbf{必要条件}：必须存在\textbf{相对运动}（旋转电机）或\textbf{交变磁通}（变压器），并且磁场与绕组之间存在相对切割或磁链交变。静止不动的导体放在恒定磁场中，既不感应电动势、也不产生持续转矩。
  \item \textbf{能量流向的判据}：$ei$ 与 $fv$ 同号为正（吸收电功率、输出机械功率 $\Rightarrow$ 电动机）；反之为发电机。
\end{itemize}

\corepoint{本章三条定律构成全书分析方法的骨架：写电压方程用 $e=-\mathrm{d}\Psi/\mathrm{d}t$（或 $e=Blv$），写转矩方程用 $f=Bli$（或 $T=P_{em}/\Omega$），检查结果用能量守恒（输入 $=$ 输出 $+$ 损耗）。任何一道电机计算题，最后都可以用"功率平衡表"验算：四大类电机的功率流程图虽然不同，但"输入 $=$ 输出 $+$ 各项损耗"永远成立。}

\begin{center}
\begin{tabular}{@{}lll@{}}
\toprule
定律 & 表达式 & 在电机分析中的用途 \\
\midrule
电磁感应定律 & $e=-N\dfrac{\mathrm{d}\Phi}{\mathrm{d}t}$，$e=Blv$ & 建立绕组电动势方程（变压器、定子、转子） \\
电磁力定律 & $f=Bli$，$T=P_{em}/\Omega$ & 建立转矩方程、机械特性 \\
能量守恒原理 & $P_1=P_2+\sum p$ & 功率平衡、效率计算、结果校验 \\
\bottomrule
\end{tabular}
\end{center}

% ============================================================
% 第 3 章 变压器概述（教材第二篇，印刷页 24--29）
% 节名：第一节 变压器的分类；第二节 变压器的基本结构；第三节 变压器的型号与额定值
% ============================================================
\chapter{变压器概述}

变压器是电力系统输电与配电的主要设备，也是测量、控制及通信等装置的重要元件，被广泛应用在各个领域，因此种类很多。本章只解决三个问题：有哪些变压器（分类）、它由什么构成（结构）、它的容量电压电流是怎么规定的（额定值）。这三件事是后面所有分析计算的"边界条件"——等值电路里的 $U_{1N}$、$I_{1N}$、$S_N$ 就是从这里来的。

\section{变压器的分类}

\subsection{按用途分类}
\begin{enumerate}
  \item \textbf{电力变压器}：在电力系统中传送和分配电能。容量大（从几十千伏安到百万伏安），电压高（从几百伏到兆伏），应用范围广。
  \item \textbf{调压变压器}：在电网或实验室中用于调节电压。
  \item \textbf{测量变压器}：如电流互感器和电压互感器等。
  \item \textbf{专用变压器}：如电炉变压器、电焊变压器、整流变压器、高压试验变压器，以及无线电通信、自动控制系统和各种仪器用的小功率变压器等。
\end{enumerate}

\subsection{按其他方式分类}
\begin{itemize}
  \item 按\textbf{绕组数目}：双绕组变压器、三绕组变压器等；
  \item 按\textbf{铁芯结构}：心式变压器、壳式变压器；
  \item 按\textbf{相数}：单相变压器、三相变压器和多相变压器（如电力电子装置中的六相、十二相变压器）；
  \item 按\textbf{冷却方式}：油浸式变压器、干式变压器。
\end{itemize}

变压器的种类虽多，但基本工作原理一致。本书着重讨论\textbf{电力变压器}。

\begin{center}
\begin{tabular}{@{}llp{7.6cm}@{}}
\toprule
分类依据 & 类别 & 特点与典型用途 \\
\midrule
用途 & 电力变压器 & 输电与配电，容量几十 kV·A～百万 kV·A \\
 & 调压变压器 & 电网或实验室调压 \\
 & 测量变压器 & 电流互感器、电压互感器 \\
 & 专用变压器 & 电炉、电焊、整流、高压试验、小功率通信与控制 \\
铁芯结构 & 心式 & 绕组包围铁芯柱，结构简单、绝缘易做，应用最广 \\
 & 壳式 & 铁芯包围绕组，机械强度高但制造复杂，用于低压大电流 \\
绕组数目 & 双/三绕组 & 三绕组用于三种电压等级的交汇点 \\
相数 & 单相/三相/多相 & 三相最常用；多相用于电力电子装置 \\
冷却 & 油浸式/干式 & 油浸式散热好、容量大；干式防火性能好 \\
\bottomrule
\end{tabular}
\end{center}

\section{变压器的基本结构}

变压器通过套装在同一铁芯柱上的两个或两个以上的绕组工作，因此基本结构件是\textbf{铁芯}和\textbf{绕组}。对于油浸式电力变压器，还有油箱、变压器油、绝缘套管和其他结构附件。铁芯和绕组合称为变压器的\textbf{器身}。

\subsection{铁芯}
铁芯是变压器的\textbf{磁路}，又是变压器的\textbf{机械骨架}。它由\textbf{铁芯柱}和\textbf{铁轭}两部分组成：铁芯柱上套装绕组，铁轭把各个铁芯柱连接起来，形成闭合铁芯磁路。

为减少铁芯中的磁滞和涡流损耗，铁芯通常选用厚度为 0.23、0.35\,mm（目前多为 0.30\,mm 及以下）的两面涂绝缘漆的\textbf{硅钢片}叠成。硅钢片的铁芯损耗小、磁导率高。

\subsubsection{心式与壳式}
\begin{itemize}
  \item \textbf{心式}：绕组包围铁芯柱。结构简单，绕组套装和绝缘比较容易，应用比较广泛。
  \item \textbf{壳式}：铁芯包围绕组。铁芯机械强度高，但制造复杂，通常只有低压大电流的变压器或小容量的电信变压器采用。
\end{itemize}

\subsubsection{叠片与接缝}
铁芯叠片通常将硅钢片剪成一定尺寸的长方形，再按要求叠装成铁芯。叠片\textbf{交错叠装}，使相邻两层的接缝互相错开，这样叠成的铁芯接缝气隙小、导磁性能好。

硅钢片有热轧和冷轧两种：
\begin{itemize}
  \item \textbf{热轧硅钢片}导磁方向性不明显，叠积铁芯可采用\textbf{直接缝}；
  \item \textbf{冷轧硅钢片}导磁性更好、损耗更小，变压器几乎都采用冷轧硅钢片；但它具有很强的\textbf{方向性}，为避免磁通转向时在拐角处增加磁阻，通常采用\textbf{斜接缝}。
\end{itemize}
由于直接缝比斜接缝的材料利用率高，变压器中直接缝与斜接缝混合使用也较普遍。

\subsubsection{铁芯柱截面}
小型变压器铁芯柱的横截面是\textbf{正方形}；容量较大的变压器，为充分利用绕组的内圆空间，一般采用\textbf{多级阶梯形}截面。当铁芯柱横截面外接圆直径大于 500\,mm 时，中间还应留出\textbf{油道}以改善铁芯内部的散热条件。铁轭的横截面形状一般为矩形，面积比铁芯柱的横截面积大 5\%～10\%，以减少励磁电流和铁芯损耗。

\subsection{绕组}
绕组是变压器的\textbf{电路部分}，通常用纸包的绝缘扁导线或漆包的圆导线绕制而成，导线多为铜线。根据高、低压绕组在铁芯柱上的相对位置，绕组的排列方式分为两种。

\subsubsection{同心式}
高、低压绕组同心地套装在铁芯柱上，\textbf{低压绕组在里面，高压绕组在外面}。这样可使绝缘距离小，绕组和铁芯的尺寸都可以做得小些，同时也便于引出高压绕组的分接头。同心式绕组制造方便，国产变压器多采用这种结构。

\subsubsection{交叠式}
高、低压线圈都做成饼状，沿铁芯柱高度方向交替放置；靠近上下铁轭的位置一般放置低压绕组。交叠式绕组\textbf{漏抗较小，机械强度较高}，易于接成多条并联支路，主要用于低电压、大电流的电焊变压器和电炉变压器。

绕组有电流流通时要发热，中大型变压器的绕组之间、一个绕组的各层之间，有时甚至各匝之间都垫上绝缘块，形成\textbf{油路}以利散热。

在配电变压器中，高压绕组一般都引出几个\textbf{分接头}，在电网电压有波动时，可通过分接开关接通不同的分接头，调节高压绕组的匝数，改变变压器的实际匝数比，使低压侧输出电压保持相对稳定。

变压器输入电能一侧的绕组称为\textbf{一次绕组}，输出电能一侧的绕组称为\textbf{二次绕组}。两侧匝数不等，通过电磁感应作用，一次绕组把电能传递到二次绕组，二次绕组就能得到不同的电压和电流。

\begin{center}
\begin{tabular}{@{}p{1.7cm}p{4.4cm}p{3.3cm}p{3.6cm}@{}}
\toprule
排列方式 & 结构 & 优点 & 主要用途 \\
\midrule
同心式 & 低压在内、高压在外，同心套在铁芯柱上 & 绝缘距离小、尺寸小、便于引出分接头、制造方便 & 国产电力变压器普遍采用 \\
交叠式 & 高、低压饼状线圈沿柱高交替放置，靠铁轭处放低压 & 漏抗小、机械强度高、易多支路并联 & 低电压大电流（电焊、电炉变压器） \\
\bottomrule
\end{tabular}
\end{center}

\subsection{油箱及变压器油}
油箱是变压器的外壳，由钢板焊接而成。油箱内盛变压器油，器身浸在油中。变压器油有\textbf{绝缘}和\textbf{散热}两个作用：绕组和铁芯产生的热量由变压器油传递给油箱壁、散热管或散热器，从而冷却器身。

\subsection{绝缘套管}
绝缘套管是变压器绕组的引出线之间以及引出线与油箱之间的绝缘，并起固定引出线的作用，用陶瓷制成，有多种形式。

除此之外，油浸式电力变压器中通常还有\textbf{储油柜}、\textbf{散热器}、\textbf{气体继电器}、\textbf{安全气道}、\textbf{吸湿器}、\textbf{分接开关}、\textbf{油位计}和继电保护装置等附件。

\subsection{干式变压器}
自 20 世纪 80 年代开始，干式变压器在我国得到迅速发展。干式变压器的铁芯和绕组不浸在任何液体中，具有运行安全可靠、维护简单、可深入负荷中心等优点。但由于空气的绝缘强度和散热性能都比油差，以空气绝缘的干式变压器有效材料消耗比油浸式多，承受冲击电压的能力也比油浸式差。干式变压器主要适用于高层建筑、机场和地铁等安全防火要求较高的配电场所，主要由铁芯和绕组组成，铁芯及其夹件充当机壳；绕组通常用纱包的绝缘扁铜导线绕制，为增强绝缘性能通常进行\textbf{树脂浇注}；较大容量的干式变压器有时配装风冷装置和罩壳。

\corepoint{结构与原理的两条对应关系必须记牢：① 铁芯为什么要闭合、要叠片、要斜接缝——都是为了"用最小的励磁磁动势建立主磁通，同时把铁耗压到最低"；② 低压绕组为什么放里面——为了减小高压绕组对铁芯的绝缘距离，从而缩小整体尺寸。凡是"结构为什么这样"的问题，答案通常落在\textbf{绝缘、散热、损耗、机械强度}这四个词上。}

\section{变压器的型号与额定值}

\subsection{变压器的型号}
变压器的型号包括说明结构特点的基本代号、额定容量和额定电压。例如型号 \textbf{S-500/10} 表示：三相油浸自冷式、双绕组电力变压器，额定容量 500\,kV·A，额定电压 10\,kV。

\begin{center}
\begin{tabular}{@{}ll@{}}
\toprule
代号 & 含义 \\
\midrule
S & 三相（D 为单相） \\
油浸自冷 & 冷却方式（干式、强迫油循环风冷等另有代号） \\
双绕组 & 绕组数目 \\
500 & 额定容量，kV·A \\
10 & 高压侧额定电压等级，kV \\
\bottomrule
\end{tabular}
\end{center}

\subsection{额定值}
变压器的额定值是制造厂对变压器正常运行时所作的使用规定，也是设计、使用和试验变压器的依据。在额定状态下运行，可以保证变压器长期可靠地工作，并具有优良的性能。

\subsubsection{额定容量 $S_N$}
在规定的额定状态下变压器输出的\textbf{视在功率}，单位 kV·A。对于双绕组电力变压器，一、二次绕组的额定容量设计值相同。

\subsubsection{额定电压 $U_{1N}$ 和 $U_{2N}$}
在规定的额定状态下，一次绕组所加的电压值称为一次绕组的额定电压 $U_{1N}$；当一次绕组加额定电压 $U_{1N}$ 时，二次绕组所得到的\textbf{空载电压}称为二次绕组的额定电压 $U_{2N}$。单位 V 或 kV。\textbf{三相变压器的额定电压是指线电压}。

\subsubsection{额定电流 $I_{1N}$ 和 $I_{2N}$}
根据额定容量和额定电压算出的\textbf{线电流}，单位 A。

\keyformula{\text{单相：}I_{1N}=\frac{S_N}{U_{1N}},\quad I_{2N}=\frac{S_N}{U_{2N}};\qquad
\text{三相：}I_{1N}=\frac{S_N}{\sqrt3\,U_{1N}},\quad I_{2N}=\frac{S_N}{\sqrt3\,U_{2N}}}

\subsubsection{额定频率 $f_N$}
我国额定工频为 50\,Hz。此外，额定工作状态下变压器的效率、温升等数据均属额定值。由于额定值都标注在变压器外壳的铭牌上，所以额定值也称为\textbf{铭牌值}。

\begin{center}
\begin{tabular}{@{}lll@{}}
\toprule
额定值 & 定义要点 & 单位 \\
\midrule
$S_N$ & 额定状态下输出的视在功率；双绕组一、二次侧设计值相同 & kV·A \\
$U_{1N}$ & 一次绕组应加的电压（三相指线电压） & V 或 kV \\
$U_{2N}$ & 一次加 $U_{1N}$ 时二次绕组的\textbf{空载}电压（三相指线电压） & V 或 kV \\
$I_{1N}$、$I_{2N}$ & 由 $S_N$ 与 $U_N$ 算出的线电流 & A \\
$f_N$ & 我国工频 50\,Hz & Hz \\
\bottomrule
\end{tabular}
\end{center}

\corepoint{三个易错点：① $U_{2N}$ 是\textbf{空载}电压，不是额定负载下的端电压——带负载后二次电压会因漏阻抗压降而低于 $U_{2N}$，这正是第 6 章"电压变化率"要算的东西；② 三相变压器的 $U_N$、$I_N$ 都是\textbf{线值}，换算到相值要除以 $\sqrt3$（星接）或乘以 $\sqrt3$（角接），做参数计算时务必先统一到相值；③ $S_N$ 是视在功率，不是有功功率，故 $I_N=S_N/(\sqrt3U_N)$ 与负载功率因数无关。}

% ============================================================
% 第 4 章 变压器的基本理论（教材第二篇，印刷页 30--47）
% 节名：第一节 变压器的空载运行；第二节 变压器的负载运行；
%       第三节 变压器的归算；第四节 变压器的等效电路和相量图；第五节 变压器的参数测定
% ============================================================
\chapter{变压器的基本理论}

本章介绍变压器的基本原理和分析方法：先从变压器内部的电磁关系出发得出基本方程式，再通过归算得到等效电路和相量图，最后给出由试验测定参数的方法。这一整套"物理过程 $\rightarrow$ 基本方程 $\rightarrow$ 归算 $\rightarrow$ 等效电路 $\rightarrow$ 参数测定"的流程，是后面分析感应电机、同步电机的通用模板，务必在本章彻底打通。

\section{变压器的空载运行}

变压器空载运行是指一次绕组接到交流电源上、二次绕组开路、二次绕组没有电流流通的运行情况。

\subsection{空载运行的物理情况和变压原理}
设一、二次绕组匝数分别为 $N_1$、$N_2$。二次绕组开路、一次绕组接上交流电源时，在外加电压作用下一次绕组中流过交变电流，称为\textbf{空载电流} $i_0$。$i_0$ 产生磁动势 $F_0=N_1i_0$，建立变压器的空载磁场。

由于铁芯磁路的磁导率远大于周围空气或油的磁导率，$F_0$ 产生的交变磁通绝大部分通过铁芯，并与一、二次绕组同时交链。这部分磁通是联系一、二次绕组的媒介，称为\textbf{主磁通} $\Phi$。另有一小部分仅与一次绕组交链、主要经铁芯外空气或油闭合的磁通，称为一次绕组的\textbf{漏磁通} $\Phi_{1\sigma}$。

按图示正方向（磁通正方向与感应电动势正方向符合右手螺旋定则），
\begin{equation}
e_1=-N_1\frac{\mathrm{d}\Phi}{\mathrm{d}t},\qquad e_2=-N_2\frac{\mathrm{d}\Phi}{\mathrm{d}t}.
\end{equation}
漏磁通仅与一次绕组交链，只在该绕组中感应漏磁电动势 $e_{1\sigma}$。由于漏磁通主要沿非铁磁材料闭合，基本不受铁芯饱和影响，故漏磁链与 $i_0$ 成正比：
\begin{equation}
\Psi_{1\sigma}=L_{1\sigma}i_0,\qquad
e_{1\sigma}=-L_{1\sigma}\frac{\mathrm{d}i_0}{\mathrm{d}t}=-N_1\frac{\mathrm{d}\Phi_{1\sigma}}{\mathrm{d}t},
\end{equation}
其中 $L_{1\sigma}$ 为一次绕组的漏磁电感（漏感），是一常值。

由基尔霍夫定律，一次绕组的电压方程为
\begin{equation}
u_1=-e_1-e_{1\sigma}+R_1i_0=-e_1+L_{1\sigma}\frac{\mathrm{d}i_0}{\mathrm{d}t}+R_1i_0 .
\end{equation}
由于漏磁通远小于主磁通，$|e_{1\sigma}|\ll|e_1|$；电阻压降 $R_1i_0$ 更小，在电力变压器中通常约占一次绕组电压的 0.1\%。二者均可忽略，故
\begin{equation}
u_1\approx-e_1,\qquad u_{20}=e_2 .
\end{equation}
于是空载时一、二次电压之比
\begin{equation}
\frac{U_1}{U_{20}}\approx\frac{E_1}{E_2}=\frac{N_1}{N_2}=k ,
\end{equation}
$k$ 称为变压器的\textbf{变比}。

\keyformula{\frac{U_1}{U_{20}}\approx\frac{E_1}{E_2}=\frac{N_1}{N_2}=k}

该式表明：空载运行时，变压器一、二次绕组的电压之比等于匝数比。因此要使一、二次绕组具有不同电压，只要使它们具有不同的匝数即可——这就是变压器变压的原理。

\subsection{主磁通与感应电动势}
一次绕组漏抗压降与电阻压降之和小于 $0.4\%U_{1N}$，可以忽略，故 $U_1\approx E_1$，即一次绕组的感应电动势近似与外加电源电压相平衡。若外加电压按正弦规律变化，$e_1$ 也按正弦变化；由 $e_1=-N_1\mathrm{d}\Phi/\mathrm{d}t$ 可知，产生正弦电动势的磁通必须按正弦规律变化。设
\begin{equation}
\Phi=\Phi_m\sin\omega t,
\end{equation}
代入 $e_1=-N_1\mathrm{d}\Phi/\mathrm{d}t$ 得
\begin{equation}
e_1=-\omega N_1\Phi_m\cos\omega t=\sqrt2E_1\sin\!\left(\omega t-90^\circ\right),
\end{equation}
故一次绕组感应电动势的有效值与主磁通的关系为
\begin{equation}
E_1=\frac{\omega N_1\Phi_m}{\sqrt2}=4.44fN_1\Phi_m .
\end{equation}

\keyformula{E_1=4.44fN_1\Phi_m,\qquad E_2=4.44fN_2\Phi_m}

同理可得二次绕组感应电动势 $E_2=4.44fN_2\Phi_m$，其相位与 $E_1$ 相同，亦滞后主磁通 $90^\circ$。

\subsubsection{两个重要推论}
\begin{enumerate}
  \item \textbf{电动势大小}与绕组匝数以及主磁通变化的频率与幅值成正比。
  \item 由 $U_1\approx E_1=4.44fN_1\Phi_m$，当外加电源电压不变时，对已制好的变压器，
  \begin{equation}
  \Phi_m\approx\frac{U_1}{4.44fN_1}=\text{常数}.
  \end{equation}
  \textbf{空载运行时的主磁通基本不变}。这条"恒磁通"结论是后面磁动势平衡方程成立的前提（第 4.2 节），也是判断"加电压还是加频率会改变铁耗"的依据。
\end{enumerate}

\subsection{励磁电流与励磁阻抗}
空载时 $i_0$ 产生的磁动势 $F_0=N_1i_0$ 用来建立空载磁场，故\textbf{空载电流基本就是励磁电流}，即 $i_0=i_m$。由第 1 章第五节，主磁通为正弦波时励磁电流为非正弦波，可分解为产生主磁通的\textbf{磁化电流} $i_\mu$ 与\textbf{铁芯损耗电流} $i_{Fe}$ 两个分量：
\begin{equation}
i_m=i_\mu+i_{Fe}.
\end{equation}
$i_\mu$ 是尖顶波，其数值与额定电流相比很小，工程上常用\textbf{等效正弦波}代替——即该正弦波的有效值、频率和相位都与尖顶波相同。经过等效后，励磁电流可以用相量表示：
\begin{equation}
\dot I_m=\dot I_\mu+\dot I_{Fe}.
\end{equation}

相量关系为：$\dot I_\mu$ 与 $\dot\Phi_m$ 同相，$\dot E_1$ 超前于 $\dot\Phi_m$ $90^\circ$，故 $\dot I_\mu$ 与 $\dot E_1$ 反相，为一\textbf{无功电流}；$\dot I_{Fe}$ 超前于 $\dot\Phi_m$ $90^\circ$，与 $\dot E_1$ 同相，为一\textbf{有功电流}。励磁电流 $\dot I_m$ 超前于磁通 $\dot\Phi_m$ 的角度很小，用 $\alpha_{Fe}$ 表示，通常称为\textbf{铁耗角}。因此铁芯损耗还可以表示为
\begin{equation}
P_{Fe}=E_1I_m\cos\!\left(90^\circ-\alpha_{Fe}\right)=E_1I_m\sin\alpha_{Fe}=E_1I_{Fe}.
\end{equation}

若将励磁电流与感应电动势 $E_1$ 直接联系起来，其关系为
\begin{equation}
\dot E_1=-\dot I_m\left(R_{Fe}+\mathrm{j}X_\mu\right)=-\dot I_mZ_m,
\end{equation}
式中：$R_{Fe}$ 为表征铁耗的参数，称为\textbf{铁耗电阻}；$X_\mu$ 为表征铁芯磁化性能的参数，称为\textbf{磁化电抗}；$Z_m$ 为变压器的\textbf{励磁阻抗}，是用串联阻抗形式表征铁芯磁化性能和铁芯损耗的一个综合参数，$Z_m=R_m+\mathrm{j}X_m$。其中 $R_m$ 为励磁电阻（表征铁耗的等效参数），$X_m$ 为励磁电抗（表征磁化性能的等效参数）。

励磁回路的并联形式由铁耗电阻 $R_{Fe}$ 与磁化电抗 $X_\mu$ 并联构成；进一步用一个等效串联阻抗 $Z_m$ 代替这两个并联支路，则由串并联等效关系
\begin{equation}
R_m=\frac{R_{Fe}X_\mu^{2}}{R_{Fe}^{2}+X_\mu^{2}},\qquad
X_m=\frac{R_{Fe}^{2}X_\mu}{R_{Fe}^{2}+X_\mu^{2}} .
\end{equation}
通常 $R_{Fe}\gg X_\mu$，故近似有 $R_m\approx X_\mu^{2}/R_{Fe}$、$X_m\approx X_\mu$。

将励磁电流等效成正弦波后，漏磁电动势可写成相量形式 $\dot E_{1\sigma}=-\mathrm{j}\dot I_0X_{1\sigma}$，其中 $X_{1\sigma}=\omega L_{1\sigma}$ 为一次绕组的\textbf{漏抗}。于是电压方程的相量形式为
\begin{equation}
\dot U_1=-\dot E_1+\dot I_0\left(R_1+\mathrm{j}X_{1\sigma}\right)=-\dot E_1+\dot I_0Z_{1\sigma},
\end{equation}
其中 $Z_{1\sigma}=R_1+\mathrm{j}X_{1\sigma}$ 为一次绕组的\textbf{漏阻抗}。

\corepoint{由于磁化曲线非线性，饱和情况下 $X_m$ 是一个变值，随饱和增加而减小。但通常电源电压变化很小，磁通变化范围很小，故可近似认为 $Z_m$ 为常值。这就是"非线性问题局部线性化"的处理手法：把受饱和影响的主磁路作为一个局部非线性环节，单独用 $Z_m$ 表征；把不受饱和影响的漏磁通分离出来，用常值 $X_{1\sigma}$、$X_{2\sigma}$ 表征。}

\section{变压器的负载运行}

\subsection{负载运行时的物理过程}
空载时由 $i_0$ 产生磁动势 $F_0=N_1i_0$ 建立主磁通 $\Phi$，在一、二次绕组中分别感应 $e_1$、$e_2$；电网电压 $u_1$ 与 $E_1$ 以及一次漏阻抗压降相平衡，维持 $i_0$ 在一次绕组内流通，使变压器中的电磁关系处于平衡状态。

二次绕组接上负载 $Z_L$ 后，在 $e_2$ 作用下二次绕组中有电流 $i_2$ 流通。$i_2$ 建立磁动势 $F_2=N_2i_2$ 作用在铁芯磁路上，引起主磁通变化，主磁通的变化又引起 $e_1$、$e_2$ 变化。在电网电压不变的情况下，$e_1$ 的变化将引起一次绕组电流从空载的 $i_0$ 过渡到负载时的 $i_1$；这时磁动势 $N_1i_1$ 和 $N_2i_2$ 同时作用在铁芯磁路上共同建立主磁通，在一、二次绕组中产生新的感应电动势，以使一、二次绕组达到新的电压平衡，维持一、二次绕组电流流通。

为便于分析，假设一次绕组漏阻抗 $Z_{1\sigma}\approx0$，则 $U_1\approx E_1$。当电网电压有效值 $U_1$ 保持不变时，主磁通 $\Phi_m$ 亦保持不变。因此一次绕组电流从空载 $i_0$ 过渡到负载 $i_1$ 时，\textbf{励磁磁动势 $N_1i_0$ 将保持不变}，故一次绕组电流增加部分产生的磁动势应完全与二次绕组磁动势 $N_2i_2$ 抵消，即
\begin{equation}
N_1\dot I_1+N_2\dot I_2=N_1\dot I_0 .
\end{equation}
考虑到 $e_1/e_2=N_1/N_2$，可得
\begin{equation}
e_1i_1+e_2i_2=e_1i_0,
\end{equation}
即一次绕组从电网增加输入的电功率传递到二次绕组，变成二次绕组获得的电功率。

综上：变压器负载运行时，通过电磁感应关系，一、二次绕组的电流紧密联系在一起——二次绕组电流的增加或减小必然同时引起一次绕组电流的增加或减小；相应地，二次绕组输出功率增加或减小时，一次绕组从电网吸收的功率必然同时增加或减小。

\subsection{磁动势方程}
作用在磁路上的合成磁动势，也就是励磁磁动势。若用一次绕组电流所产生的磁动势来表征励磁磁动势，则磁动势方程为
\begin{equation}
N_1\dot I_m=N_1\dot I_1+N_2\dot I_2
\quad\Longleftrightarrow\quad
N_1\dot I_1=N_1\dot I_m+\left(-N_2\dot I_2\right).
\end{equation}
若用电流表示，
\begin{equation}
\dot I_1=\dot I_m+\dot I_{1L},\qquad
\dot I_{1L}=-\frac{N_2}{N_1}\dot I_2=-\frac{\dot I_2}{k},
\end{equation}
其中 $\dot I_{1L}$ 称为一次绕组的\textbf{负载分量}。

该式表明：变压器负载运行时，一次绕组电流有两个分量——一个是\textbf{励磁电流} $\dot I_m$，用于建立负载运行时的主磁通；另一个是\textbf{负载分量} $\dot I_{1L}$，产生磁动势 $N_1\dot I_{1L}$，用以抵消二次绕组磁动势 $N_2\dot I_2$，即 $N_1\dot I_{1L}+N_2\dot I_2=0$。由于一次绕组漏阻抗很小，从空载到额定负载，$E_1$ 变化很小，因此主磁通和励磁磁动势 $N_1\dot I_m$ 变化都不大，故 $\dot I_0$ 与 $\dot I_m$ 差别不大。

\subsection{电压方程}
负载运行时磁通包括三部分：励磁电流产生的主磁通 $\Phi$、$N_1\dot I_1$ 在一次绕组中产生的漏磁通 $\Phi_{1\sigma}$、$N_2\dot I_2$ 在二次绕组中产生的漏磁通 $\Phi_{2\sigma}$。主磁通在一、二次绕组中感应 $E_1$、$E_2$；漏磁通分别感应漏磁电动势 $E_{1\sigma}$、$E_{2\sigma}$。漏磁通与产生它的电流成正比，故
\begin{equation}
\dot E_{1\sigma}=-\mathrm{j}\dot I_1X_{1\sigma},\qquad
\dot E_{2\sigma}=-\mathrm{j}\dot I_2X_{2\sigma},
\end{equation}
其中 $X_{2\sigma}$ 为二次绕组的漏抗。

从一次电路看，负载运行与空载运行的区别仅是把电流 $i_0$ 改变为 $i_1$，故一次电压方程为
\begin{equation}
\dot U_1=-\dot E_1+\dot I_1\left(R_1+\mathrm{j}X_{1\sigma}\right)=-\dot E_1+\dot I_1Z_{1\sigma}.
\end{equation}
从二次电路看，二次绕组中有了电流 $i_2$，感应电动势不再等于端电压，由基尔霍夫第二定律
\begin{equation}
\dot E_2=\dot U_2+\dot I_2\left(R_2+\mathrm{j}X_{2\sigma}\right)=\dot U_2+\dot I_2Z_{2\sigma},
\end{equation}
即 $\dot U_2=\dot E_2-\dot I_2Z_{2\sigma}$，其中 $Z_{2\sigma}=R_2+\mathrm{j}X_{2\sigma}$ 称为二次绕组的漏阻抗。

归纳起来，变压器负载运行时的基本方程组为
\begin{equation}
\left\{
\begin{aligned}
\dot U_1&=-\dot E_1+\dot I_1Z_{1\sigma},\\
\dot E_1&=k\dot E_2,\\
\dot I_1&=\dot I_m+\dot I_{1L}=\dot I_m-\dot I_2/k,\\
\dot E_1&=-\dot I_mZ_m,\\
\dot U_2&=\dot E_2-\dot I_2Z_{2\sigma},\\
\dot U_2&=\dot I_2Z_L .
\end{aligned}
\right.
\end{equation}

\corepoint{这组方程就是变压器的"牛顿定律"。五个方程、五个未知量（$\dot E_1$、$\dot E_2$、$\dot I_1$、$\dot I_2$、$\dot U_2$），给定 $U_1$ 与 $Z_L$ 即可解。注意第三个方程是\textbf{磁动势平衡}（不是电流平衡），第四个方程把主磁路的非线性"打包"进 $Z_m$——这两个处理是全书通用套路。}

\section{变压器的归算}

原则上，应用基本方程组即可进行变压器的运行分析与计算。但由于一、二次绕组匝数不等，给计算带来麻烦，尤其当 $N_1$ 与 $N_2$ 相差较大时，绕组相对应的量值相差较大，运行计算和相量图绘制都比较困难。解决这个问题通常采用\textbf{归算法}。

所谓归算法，就是把各绕组的匝数变换成相同匝数的方法。\textbf{归算仅仅是分析变压器的一种方法，不改变变压器原有的电磁关系和电磁过程}——归算前后的磁平衡关系、功率传递、损耗和漏磁场储能均应保持不变。通常的做法是用一个匝数亦为 $N_1$ 的新绕组代替原来的二次绕组，而一次绕组匝数不变，称为将二次绕组\textbf{归算到一次侧}。归算后二次侧各物理量的数值称为\textbf{归算值}，在原来符号右上角加"$'$"表示。

\subsection{二次绕组电流的归算}
归算前后二次绕组的磁动势不能改变，即
\begin{equation}
I_2'N_1=I_2N_2\quad\Longrightarrow\quad
I_2'=\frac{N_2}{N_1}I_2=\frac{I_2}{k}.
\end{equation}

\subsection{二次绕组电动势的归算}
归算前后作用在磁路上的磁动势没变，主磁通不会改变，而感应电动势与匝数成正比，于是
\begin{equation}
\frac{E_2'}{E_2}=\frac{N_1}{N_2}=k\quad\Longrightarrow\quad
E_2'=kE_2=E_1 .
\end{equation}
归算后从一次绕组传输到二次侧的视在功率为 $E_2'I_2'=kE_2\cdot I_2/k=E_2I_2$，与归算前相等，满足\textbf{功率不变}原则。

\subsection{二次绕组端电压的归算}
根据功率不变原则，归算前后二次绕组输出的视在功率也不应改变，即 $U_2'I_2'=U_2I_2$，故
\begin{equation}
U_2'=kU_2 .
\end{equation}

\subsection{二次绕组阻抗的归算}
根据归算前后\textbf{损耗不变}原则，$I_2'^2R_2'=I_2^2R_2$，故
\begin{equation}
R_2'=k^{2}R_2 ;
\end{equation}
根据归算前后\textbf{漏磁场储能不变}原则，$\tfrac12L_{2\sigma}'I_2'^2=\tfrac12L_{2\sigma}I_2^2$，故
\begin{equation}
X_{2\sigma}'=k^{2}X_{2\sigma},\qquad Z_{2\sigma}'=k^{2}Z_{2\sigma},\qquad Z_L'=k^{2}Z_L .
\end{equation}

\begin{center}
\begin{tabular}{@{}lll@{}}
\toprule
物理量 & 归算规律 & 依据 \\
\midrule
电动势、电压（V） & 乘 $k$ & 电动势 $\propto$ 匝数；功率不变 \\
电流（A） & 除 $k$ & 磁动势不变 \\
电阻、电抗、阻抗（$\Omega$） & 乘 $k^{2}$ & 损耗不变、漏磁场储能不变 \\
\bottomrule
\end{tabular}
\end{center}

综合上述分析：若把二次绕组归算到一次侧，归算后的二次侧各量，凡单位为 V 的量（电动势和电压）归算值等于原值乘以 $k$；凡单位为 $\Omega$ 的量（电阻、电抗、阻抗）归算值等于原值乘以 $k^2$；电流的归算值则等于原值乘以 $1/k$。

将基本方程组中第二个方程式两边同乘以 $k$、第四个方程式两边同除以 $N_1$，便得到归算后变压器的基本方程组：
\begin{equation}
\left\{
\begin{aligned}
\dot U_1&=-\dot E_1+\dot I_1Z_{1\sigma},\\
\dot U_2'&=\dot E_2'-\dot I_2'Z_{2\sigma}',\\
\dot E_1&=\dot E_2',\\
\dot I_1&=\dot I_m+\left(-\dot I_2'\right),\\
\dot E_1&=-\dot I_mZ_m .
\end{aligned}
\right.
\end{equation}

在变压器中有时也将一次侧各量归算到二次侧：二次侧各量保持不变，一次侧电动势、电压的归算值等于原值乘以 $1/k$；一次侧电阻、电抗和阻抗的归算值等于原值乘以 $1/k^2$；电流的归算值等于原值乘以 $k$——与从二次侧归算到一次侧正好相反。

\corepoint{归算的三个"不变量"是判错的标准答案：\textbf{磁动势不变、功率不变、损耗与漏磁场储能不变}。做题时如果归算后算出的铜耗或输出功率与归算前不同，一定是 $k$ 的幂次用错了。记忆口诀："电压乘 $k$、电流除 $k$、阻抗乘 $k^2$"，但\textbf{三相变压器要用相电压比}，用线电压比会把 $\sqrt3$ 带进 $k^2$，这是最常见的失分点。}

\section{变压器的等效电路和相量图}

归算的目的不仅是使计算与相量图绘制方便，更重要的是可以把既有电关系又有磁关系的实际变压器，用一个\textbf{单纯的电路}来表示——这种电路称为\textbf{等效电路}。

\subsection{T 型等效电路}
分析归算后的基本方程组：第一个方程对应一次侧的串联支路 $Z_{1\sigma}$，第二个方程对应二次侧的串联支路 $Z_{2\sigma}'$ 与负载 $Z_L'$，第四个方程对应励磁支路 $Z_m$。把三个电路的尾端连起来不会改变各自特性；又因 $\dot E_1=\dot E_2'$，三个对应节点等电位，将这三个节点直接相连，即完全满足基本方程组——所得到的电路就是变压器的 \textbf{T 型等效电路}：一次漏阻抗 $Z_{1\sigma}$ 与励磁支路 $Z_m$ 串联后，再与二次漏阻抗 $Z_{2\sigma}'$ 及负载 $Z_L'$ 串联。

运用 T 型等效电路可以很精确地进行变压器的运行分析与计算。

\subsection{相量图}
相量图根据基本方程组画出，特点是可以直观看出各物理量的大小和相位关系。画相量图时认为 $R_1$、$R_2'$、$X_{1\sigma}$、$X_{2\sigma}'$ 和 $Z_m$ 均为已知，负载阻抗 $Z_L'$（或负载电流 $\dot I_2'$）亦给定。

\subsubsection{作图步骤（感性负载）}
\begin{enumerate}
  \item 以负载端电压 $\dot U_2'$ 为参考相量；
  \item 由给定负载阻抗 $Z_L'$ 画出负载电流 $\dot I_2'=\dot U_2'/Z_L'$；
  \item 由 $\dot U_2'=\dot E_2'-\dot I_2'Z_{2\sigma}'$ 反求 $\dot E_2'=\dot E_1$：在 $\dot U_2'$ 上加阻抗压降 $\dot I_2'R_2'+\mathrm{j}\dot I_2'X_{2\sigma}'$，其中 $\dot I_2'R_2'$ 平行于 $\dot I_2'$，$\mathrm{j}\dot I_2'X_{2\sigma}'$ 超前 $\dot I_2'$ $90^\circ$；
  \item 主磁通 $\dot\Phi_m$ 超前 $\dot E_1$（即 $\dot E_2'$）$90^\circ$；
  \item 由 $\dot E_1=-\dot I_mZ_m$ 画励磁电流 $\dot I_m$，它应滞后于 $-\dot E_1$ 角度 $\alpha_{Fe}=\arctan\left(X_m/R_m\right)$（铁耗角）；
  \item 由 $\dot I_1=\dot I_m+\left(-\dot I_2'\right)$ 得一次电流 $\dot I_1$；
  \item 由 $\dot U_1=-\dot E_1+\dot I_1Z_{1\sigma}$ 得一次电压 $\dot U_1$：在 $-\dot E_1$ 上加 $\dot I_1R_1$（平行 $\dot I_1$）与 $\mathrm{j}\dot I_1X_{1\sigma}$（超前 $\dot I_1$ $90^\circ$）。
\end{enumerate}

相量图主要用以直观反映各量之间的大小和相位关系，具体计算一般还是利用等效电路。

\subsection{近似和简化等效电路}
对于普通电力变压器，励磁电流与额定电流相比很小（通常占额定电流的千分之几至百分之几），它在 $Z_{1\sigma}$ 上引起的压降很小；即使是额定电流在 $Z_{1\sigma}$ 上引起的压降 $I_{1N}Z_{1\sigma}$ 也不超过额定电压的百分之几。因此把 T 型等效电路中的励磁分支从中间位置移到\textbf{电源端}，对运行计算仍然很精确，这就是\textbf{近似等效电路}。

如果进一步忽略励磁电流，则得到\textbf{简化等效电路}——此时变压器表现为一个串联阻抗
\begin{equation}
Z_k=R_k+\mathrm{j}X_k,\qquad R_k=R_1+R_2',\qquad X_k=X_{1\sigma}+X_{2\sigma}' ,
\end{equation}
其中 $Z_k$ 为变压器的\textbf{短路阻抗}，$R_k$ 为\textbf{短路电阻}，$X_k$ 为\textbf{短路电抗}。

工程上常用简化等效电路进行定性分析和简化计算；在电力系统计算中，变压器的等效电路几乎都采用简化等效电路。

\begin{center}
\begin{tabular}{@{}p{2.4cm}p{3.6cm}p{4.2cm}p{3.4cm}@{}}
\toprule
等效电路 & 结构 & 适用条件 & 典型用途 \\
\midrule
T 型 & 励磁支路接在两侧漏阻抗之间 & 精确分析，任何负载 & 教学、精确计算 \\
近似 & 励磁支路移到电源端 & 励磁电流远小于额定电流 & 一般运行计算 \\
简化 & 只有 $Z_k=R_k+\mathrm{j}X_k$ & 忽略励磁电流，负载较大 & 电力系统计算、短路计算 \\
\bottomrule
\end{tabular}
\end{center}

\section{变压器的参数测定}

变压器的参数可以通过\textbf{空载试验}和\textbf{短路试验}测定。

\subsection{空载试验}
空载试验的目的是测定变压器的\textbf{变比}、空载电流特性、空载损耗特性和\textbf{励磁阻抗}。为了安全和便于试验，通常\textbf{在低压端加电压，高压端开路}。

\subsubsection{变比}
低压绕组加电源电压，用电压表分别测出低压端的外加电压 $U_1$ 和高压端的开路电压 $U_{20}$，便可算出变比
\begin{equation}
k=\frac{N(\text{高压})}{N(\text{低压})}=\frac{U_{20}}{U_1}.
\end{equation}

\subsubsection{空载特性}
用调压器把电源电压调到变压器低压端额定电压的 1.2 倍左右，然后逐步降低电压，同时测量空载电流 $I_0$、外加电压 $U_1$ 和相应的空载功率 $p_0$，得到空载特性曲线 $I_0=f(U_1)$ 和 $p_0=f(U_1)$。

空载时漏阻抗压降相当小，电压 $U_1$ 与主磁通成正比，且空载电流就是励磁电流，所以 $I_0=f(U_1)$ 实质上反映了变压器铁芯的\textbf{基本磁化曲线}：电压较低时磁路不饱和，$I_0=f(U_1)$ 是线性关系；电压增高时磁路逐渐饱和，$I_0$ 迅速增加。

空载时变压器从电网吸收的功率 $p_0$ 为铁芯损耗 $p_{Fe}$ 与空载绕组损耗 $p_{Cu0}$ 之和。$p_{Cu0}=I_0^2R_1$ 与空载电流二次方成正比；铁芯损耗正比于铁芯磁密（或 $U_1$）的二次方。空载时 $I_0=I_m$ 非常小，可认为 $p_{Cu0}\approx0$，所以 $p_0\approx p_{Fe}$，与外加电压 $U_1$ 的二次方成正比——故\textbf{空载损耗特性 $p_0=f(U_1)$ 基本上是抛物形曲线}。

\subsubsection{励磁参数}
在空载特性曲线上查出额定电压时的 $I_0$ 和 $p_0$，便得到额定电压下的励磁参数：
\begin{equation}
Z_m=\frac{U_{1N}}{I_0},\qquad
R_m=\frac{p_0}{I_0^{2}},\qquad
X_m=\sqrt{Z_m^{2}-R_m^{2}} ,
\end{equation}
式中 $U_{1N}$ 为一次绕组（此处的加压侧）的额定相电压。

\textbf{注意}：此时的励磁参数为\textbf{低压侧的值}，若归算到高压侧，各参数均应乘以 $k^2$。

\subsection{短路试验}
短路试验时把二次绕组短路，一次绕组外加可调的低电压。为便于测量，通常将\textbf{高压侧作为一次侧}，低压侧作为二次侧。通过短路试验可以确定变压器的\textbf{短路特性}和\textbf{短路参数}。

\subsubsection{短路特性}
用调压器升高外加电压 $U_k$，使短路电流达到高压绕组额定电流的 1.2 倍为止，然后再逐渐降低电压，同时测量 $U_k$、$I_k$ 和相应的短路功率 $p_k$，得到短路特性曲线。

短路试验时外加电压很低，一般小于额定电压的 10\%，磁路不饱和，$I_k=f(U_k)$ 为直线。由于电压低，短路时的铁芯损耗很小可以忽略，故短路功率约等于一、二次绕组的电阻损耗之和：
\begin{equation}
p_k\approx I_1^{2}R_1+I_2'^{2}R_2'\approx I_k^{2}R_k ,
\end{equation}
且 $p_k$ 与 $U_k$ 为线性关系，于是 $p_k=f(U_k)$ 为抛物形曲线（$p_k\propto U_k^2$ 的变形写法）。

\subsubsection{短路参数}
在短路特性曲线上查出短路电流为额定电流时的 $U_k$ 和 $p_k$，可得
\begin{equation}
Z_k=\frac{U_k}{I_k},\qquad
R_k=\frac{p_k}{I_k^{2}},\qquad
X_k=\sqrt{Z_k^{2}-R_k^{2}} .
\end{equation}
所得到的结果是\textbf{折算到高压侧}的值。

由于 $X_{1\sigma}$ 与 $X_{2\sigma}'$ 从 $X_k$ 中分离开是极其困难的，对于近似和简化等效电路可直接使用上式求得的短路参数，无需将各参数分开；对于 T 型等效电路，工程上假设
\begin{equation}
X_{1\sigma}=X_{2\sigma}'=\frac{X_k}{2},\qquad
R_1=R_2'=\frac{R_k}{2},
\end{equation}
这对于工程计算已足够精确。

\subsubsection{热态电阻的换算}
绕组电阻随温度变化，而短路试验通常在室温下进行，故测得的电阻必须换算到规定工作温度。我国国家标准规定，油浸式电力变压器的短路电阻须换算到 75\,℃ 时的值。对于铜线绕组
\begin{equation}
R_k(75\,^\circ\mathrm{C})=R_k(\theta)\cdot\frac{235+75}{235+\theta},
\end{equation}
式中 $\theta$ 为试验时的室温（℃），$R_k(\theta)$ 为室温 $\theta$ 时的绕组电阻；铝线绕组把 235 换成 228。
换算后再算短路阻抗：
\begin{equation}
Z_k(75\,^\circ\mathrm{C})=\sqrt{R_k^{2}(75\,^\circ\mathrm{C})+X_k^{2}} .
\end{equation}

\begin{center}
\begin{tabular}{@{}p{1.9cm}p{3.4cm}p{3.9cm}p{4.0cm}@{}}
\toprule
项目 & 接线与条件 & 求得的主要参数 & 关键判据 \\
\midrule
空载试验 & 低压侧加压、高压侧开路 & $k$、$Z_m$、$R_m$、$X_m$（低压侧值） & $p_0\approx p_{Fe}$，$p_0\propto U_1^2$ \\
短路试验 & 高压侧加压、低压侧短路 & $Z_k$、$R_k$、$X_k$（高压侧值） & $p_k\approx p_{Cu}$，$p_k\propto I_k^2$ \\
热态换算 & 铜线 235、铝线 228 & 75\,℃ 下的 $R_k$、$Z_k$ & 国标规定油浸式换算到 75\,℃ \\
\bottomrule
\end{tabular}
\end{center}

\corepoint{六个易错点：① 空载试验在\textbf{低压侧}加压，所得 $Z_m$ 是低压侧值，归算到高压侧要乘 $k^2$；② 短路试验在\textbf{高压侧}加压，所得 $Z_k$ 是高压侧值；③ 空载损耗近似等于铁耗，但严格说含少量铜耗；短路损耗近似等于铜耗，但严格说含少量铁耗；④ $X_{1\sigma}$ 与 $X_{2\sigma}'$ 无法由试验分离，T 型电路只能各取一半；⑤ 短路试验数据必须换算到 75\,℃，否则 $R_k$ 偏小；⑥ 做参数计算前先把\textbf{线值换成相值}，三相变压器尤其容易漏除 $\sqrt3$。}

% ============================================================
% 第 5 章 三相变压器（教材第二篇，印刷页 48--59）
% 节名：第一节 三相变压器组和三相心式变压器；第二节 三相变压器的联结组；
%       第三节 三相变压器中励磁电流、主磁通和感应电动势波形；第四节 标幺值
% ============================================================
\chapter{三相变压器}

由于电能的产生、传输和分配绝大部分采用三相制，三相变压器的应用十分广泛。实际运行时三相变压器的电压和电流基本上是对称的，对称三相电压或电流大小相等、相位互差 $120^\circ$，因此可取三相中的一相来分析计算，即将三相问题简化成单相问题。这样，第 4 章的基本方程、等效电路和相量图对三相变压器仍然适用，本章只介绍三相变压器的几个特殊问题：磁路结构、联结组、波形与标幺值。

\section{三相变压器组和三相心式变压器}

三相变压器按铁芯结构分类，通常有三相变压器组和三相心式变压器两种。

\subsection{三相变压器组}
三相变压器组由三个相同的单相变压器按一定的连接方式组合而成。三相变压器组的\textbf{三相磁路各自独立}。当一次绕组外加三相对称电压时，三相磁路中的主磁通和三相绕组中的励磁电流基波都是三相对称的。

\subsection{三相心式变压器}
如果把三个单相变压器的铁芯合并，每相的高、低压绕组放在一起，并将三相绕组分别套装在外侧的三个铁芯柱上，则磁路仍保持三相对称。当三相绕组外接对称三相电压时，三相磁通亦为三相对称，且三相磁通之和为零：
\begin{equation}
\dot\Phi_A+\dot\Phi_B+\dot\Phi_C=0 .
\end{equation}
也就是说，中间的铁芯柱中没有磁通经过，可以将它省去；如果再使三个铁芯柱布置在同一平面上，便得到三相心式变压器的铁芯。

三相心式变压器的磁路是\textbf{相互联系}的，相当于电路中的 Y 形连接，任一相的磁通都要借助另外两相的磁路闭合。由于三相磁路的长度不同（中间芯柱的磁路较短、磁阻较小），三相磁路的磁阻不同：在三相对称电压作用下，三相磁路中的主磁通相同，但产生主磁通所需的励磁电流却不同，中间的励磁电流稍小一些。工程上常取它们的平均值作为实际的励磁电流。由于磁路的不对称程度很小，三相励磁电流差别不大，对运行性能影响也很小，通常可不予考虑。

\begin{center}
\begin{tabular}{@{}p{2.3cm}p{3.5cm}p{4.0cm}p{3.6cm}@{}}
\toprule
项目 & 三相变压器组 & 三相心式变压器 & 说明 \\
\midrule
磁路 & 三相独立 & 三相相互联系（Y 形） & 决定 3 次谐波磁通能否流通 \\
磁阻 & 三相等 & 三相不等（中柱小） & 励磁电流略有差异，取平均值 \\
硅钢片用量 & 多 & 少（省去中间柱） & 心式质量轻、造价低 \\
制造与运输 & 便于制造、运输，可降低备用容量 & 大容量时运输受限 & 组式多用于巨型大容量 \\
\bottomrule
\end{tabular}
\end{center}

心式变压器的优点是：硅钢片用量少，质量轻，造价便宜，应用比较广泛。三相变压器组便于制造和运输，也可使用户的备用容量降低，故多用于巨型大容量的变压器。

\section{三相变压器的联结组}

\subsection{绕组标记和三相绕组的连接方法}
变压器高、低压绕组出线端的标记，国家标准有统一规定。单相变压器的高压绕组出线端以大写字母 $A$、$X$ 标记，低压绕组出线端以小写字母 $a$、$x$ 标记，其中 $A$、$a$ 表示首端，$X$、$x$ 表示末端。三相变压器的高压绕组首端标以 $A$、$B$、$C$，末端标以 $X$、$Y$、$Z$；低压绕组的首端标以 $a$、$b$、$c$，末端标以 $x$、$y$、$z$。

三相绕组连接方法很多，最常用的有两种：
\begin{itemize}
  \item \textbf{星形接法}：把三相绕组的末端连在一起，把它们的首端引出来；
  \item \textbf{三角形接法}：又分两种——按 $ax$—$by$—$cz$ 顺序连接称为\textbf{顺序三角形接法}；按 $ax$—$cz$—$by$ 顺序连接称为\textbf{逆序三角形接法}。
\end{itemize}

GB 1094.1—1996《电力变压器 第 1 部分：总则》规定：星形接法高压绕组用大写 $Y$ 表示、低压绕组用小写 $y$ 表示，有中性线引出时用 $YN$ 或 $yn$ 表示；三角形接法高压绕组用 $D$ 表示、低压绕组用 $d$ 表示。\textbf{高压侧接法写在前面，低压侧接法写在后面}。将高、低压绕组作不同接法组合，可得到变压器的四种基本联结组：$Yy$、$Yd$、$Dy$、$Dd$，其中 $Y(y)$ 联结组也包括有中性线引出的 $YN(yn)$ 联结组。前两种是国产变压器最常用的联结组，最后一种在国产变压器中基本不用。

\subsection{单相变压器高、低压绕组电动势之间的相位关系}
单相变压器中，高、低压绕组的电动势要么同相、要么反相，取决于两个绕组的绕制方向和标记方法。由于高、低压绕组被同一磁通 $\Phi$ 交链，其感应电动势有一定的极性关系：当高压绕组某一端瞬时电位为正时，低压绕组也有一电位为正的对应端，这两个对应的同极性端点称为\textbf{同名端}，用符号"$\cdot$"表示。

高、低压绕组的绕制可以同向也可以反向，任一出线端可以作为首端也可以作为末端。设绕组电动势正方向均规定为从首端指向末端，则根据绕制方向和标记方法有四种组合：
\begin{itemize}
  \item 绕向相同、标记相同：同名端在相同端，$\dot E_A$ 与 $\dot E_a$ 同相；
  \item 绕向相同、标记相反：同名端在不同端，$\dot E_A$ 与 $\dot E_a$ 反相；
  \item 绕向相反、标记相同：同名端在不同端，反相；
  \item 绕向相反、标记相反：同名端在相同端，同相。
\end{itemize}

规则可概括为：\textbf{若首—首（或末—末）为同名端，则 $\dot E_A$ 与 $\dot E_a$ 同相；若首—末（或末—首）为同名端，则反相}。

变压器中通常用\textbf{时钟法}表示高、低压绕组电动势的相位关系：高压绕组电动势大，代表时钟的\textbf{长针}，并认为它始终指向 0 点钟位置；低压绕组电动势小，代表\textbf{短针}。同相时相当于 0 点钟，反相时相当于 6 点钟。因此单相变压器只有两种情况，分别用 $Ii0$ 和 $Ii6$ 表示（$I$ 或 $i$ 表示单相变压器，数字表示相位关系）。国家标准规定 $Ii0$ 为标准联结组。

\subsection{三相变压器高、低压绕组线电动势之间的相位关系}
三相变压器除相电动势的相位关系外，还有线电动势的相位关系。后者除决定于绕向和标记方法外，还决定于\textbf{三相绕组的连接方法}。

\subsubsection{Yy 联结组}
若高压绕组 $A$、$B$、$C$ 相首端分别与低压绕组 $a$、$b$、$c$ 相首端为同名端，则高压绕组的相电动势 $\dot E_A$、$\dot E_B$、$\dot E_C$ 分别与低压绕组相对应的相电动势 $\dot E_a$、$\dot E_b$、$\dot E_c$ 同相位。由于两侧都是 Y 接，相对应的线电动势也同相。

分析线电动势时用 $\dot E_{BC}=\dot E_B-\dot E_C$，方向由 $\dot E_C$ 的箭头指向 $\dot E_B$ 的箭头；同理 $\dot E_{bc}=\dot E_b-\dot E_c$。两电动势相量平行，夹角为零，具有相同相位。用时钟法表示时，在三相系统中高压绕组的线电动势表示长针并指向 0（或 12），低压绕组的线电动势表示短针，故相位关系相当于 0 点钟，该联结组别为 \textbf{Yy0}。

如果将高压绕组 $A$、$B$、$C$ 相首分别与低压绕组的 $a$、$b$、$c$ 相\textbf{尾}为同名端，则联结组别变成 \textbf{Yy6}。

\subsubsection{Yd 联结组}
若高压绕组的相电动势 $\dot E_A$、$\dot E_B$、$\dot E_C$ 分别与低压绕组相对应的相电动势同相位，但由于高压绕组为 Y 接，\textbf{线电动势与相电动势之间有 $30^\circ$ 的相位移}；低压绕组为 d 接，线电动势等于相电动势。因此虽然相电动势相位相同，但由于接法不同，两侧相对应线电动势的相位不会相同。

以 $\dot E_{BC}$ 与 $\dot E_{bc}$ 为例：高压侧为 Y 接，$\dot E_{BC}$ 与图 5-5 相同；低压侧由 $\dot E_{bc}=\dot E_b-\dot E_c$，其相量与 $\dot E_b$ 相反，超前 $30^\circ$，相当于超前一个钟点。按时钟表示法，相位关系为 11 点钟，故该联结组别为 \textbf{Yd11}。

如果将同名端改为非同名端，则联结组别变成 \textbf{Yd5}；如果其他不变、仅将低压绕组由逆接改为顺接，则变成 \textbf{Yd1}。

\subsubsection{12 个组别}
三相变压器的联结组共有 12 个，相邻两个联结组相位差 $30^\circ$，12 个联结组构成圆的一周，和时钟的 12 个整点相同——这实质上是联结组用时钟法表示的原因。

在 Yy0（或 Yy6）基础上，若高压绕组三相标号 $AX$、$BY$、$CZ$ 不变，把低压绕组三相标号 $ax$、$by$、$cz$ 按顺时针方向旋转 $120^\circ$（$ax$ 被 $by$ 取代、$by$ 被 $cz$ 取代、$cz$ 被 $ax$ 取代），相当于短针顺时针转动 4 个钟点，得到 Yy4（或 Yy10）；逆时针旋转 $120^\circ$ 则得到 Yy8（或 Yy2）。故 \textbf{Yy 联结组具有 0、4、8、6、10、2 这 6 个偶数组别}；同理 \textbf{Yd 联结组具有 11、3、7、5、9、1 这 6 个奇数组别}。

Yy 联结组的某一组号通常只有一种连接方式；而 Yd（或 Dy）联结组的某一组别通常有\textbf{顺接和逆接两种方式}。为制造和使用方便，国家标准规定我国只生产 \textbf{Yyn0、Yd11、YNd11、YNy0 以及 Yy0} 共五种标准联结组，其中前三种最常用。

\begin{center}
\begin{tabular}{@{}lll@{}}
\toprule
联结组 & 钟点数 & 特点与用途 \\
\midrule
Yyn0 & 0 & 配电变压器最常用，可引出中性线供三相四线制 \\
Yd11 & 11 & 大容量输电、降压变压器常用（一侧三角形可抑制 3 次谐波） \\
YNd11 & 11 & 高压侧中性点可接地，用于系统联络 \\
YNy0 & 0 & 高压侧中性点接地 \\
Yy0 & 0 & 仅中小容量 \\
\bottomrule
\end{tabular}
\end{center}

\corepoint{判断联结组的通用步骤（考场上照做即可）：① 按同名端标出一次侧三相相电动势相量（互差 $120^\circ$）；② 按同名端定出二次侧相电动势与一次侧相电动势的同相/反相关系；③ 按接法（Y 或 D）由相电动势合成线电动势；④ 比较同一对线电动势（如 $E_{AB}$ 与 $E_{ab}$）的夹角，除以 $30^\circ$ 即为钟点数。\textbf{不要背结论，把四步画一遍}——组别题一旦背结论必错。}

\section{三相变压器中励磁电流、主磁通和感应电动势波形}

由第 1 章第五节：磁路饱和时，若磁通为正弦波，则励磁电流为尖顶波；尖顶波的主要成分是基波与 3 次谐波。3 次谐波电流频率为基波 3 倍，在三相系统中\textbf{三相同相位}：
\begin{equation}
i_{m3A}=I_{m3}\sin3\omega t,\quad
i_{m3B}=I_{m3}\sin3\left(\omega t-120^\circ\right)=I_{m3}\sin3\omega t,\quad
i_{m3C}=I_{m3}\sin3\omega t .
\end{equation}
反之，当励磁电流为正弦波时，铁芯中的磁通为平顶波，其主要成分同样是基波与 3 次谐波，且 3 次谐波磁通也同相位。

对单相变压器，3 次谐波电流可以流通，因此在正弦电源电压作用下主磁通及感应电动势都接近正弦波。但在三相变压器中，3 次谐波电流和 3 次谐波磁通能否流通，分别要由\textbf{绕组的联结方式}和\textbf{磁路结构}决定。

\subsection{YNy 联结组}
高压绕组有中性线时，中性线上的电流为 $i_N=i_{mA}+i_{mB}+i_{mC}$。由于励磁电流中三相基波分量之和为零（基波电流仅在三相绕组内流通，不经过中性线），所以中性线上的电流只有 3 次谐波电流：
\begin{equation}
i_N=i_{m3A}+i_{m3B}+i_{m3C}=3I_{m3}\sin3\omega t .
\end{equation}
即励磁电流中的 3 次谐波分量在三相绕组与中性线之间形成闭合回路，可以在绕组中流通。此时情况与单相变压器相似：外接正弦三相电压时，即使磁路饱和，铁芯中的磁通波形和绕组中的感应电动势波形仍接近正弦波。由于中性线上的谐波电流会产生损耗和通信干扰，YNy 联结组很少采用。

\subsection{Yy 联结组}
去掉中性线后，励磁电流中的 3 次谐波分量不能流通，励磁电流波形接近正弦波；于是主磁通被迫畸变，成为\textbf{平顶波}。3 次谐波磁通能否流通由磁路结构决定。

\subsubsection{三相变压器组}
三相磁路分开、各自独立，3 次谐波磁通可以和基波磁通一样在铁芯内闭合。基波磁通产生基波电动势 $e_1$，3 次谐波磁通产生 3 次谐波电动势 $e_3$，两者的相位都滞后于各自的磁通 $90^\circ$。虽然 $\Phi_{m3}<\Phi_{m1}$，但 $e_3\propto3f\Phi_{m3}$，所以 3 次谐波电动势相当大，可达基波电动势的 45\%～60\%。因此 $e_1$ 与 $e_3$ 叠加后，每相绕组的合成电动势发生较大畸变，成为\textbf{尖顶波}；尖顶波使相电动势出现相当大的尖峰值，对绕组绝缘不利。因此\textbf{三相变压器组不宜采用 Yy 联结组}。在三相线电动势中，3 次谐波电动势相互抵消，故仍为正弦波。

\subsubsection{三相心式变压器}
铁芯磁路相互联系，相当于星形接法的磁路，3 次谐波磁通在铁芯磁路中受到抑制，不得不从空气（或油）、油箱壁等\textbf{漏磁路}中闭合。漏磁路磁阻较大，$\Phi_{m3}$ 很小，因而 $e_3$ 也很小，主磁通和感应电动势仍接近正弦波。但 3 次谐波磁通流经漏磁路的油箱壁等铁磁部件，会引起附加涡流损耗、局部发热并降低效率，所以采用 Yy 联结组的三相心式变压器容量不宜太大，通常不超过 1800\,kV·A。

\subsection{Dy 和 Yd 联结组}
\begin{itemize}
  \item \textbf{Dy}：一次绕组为三角形连接，3 次谐波电流可以流通，故在正弦三相电压下励磁电流为尖顶波，主磁通及 $e_1$、$e_2$ 皆为正弦波，与单相变压器相似。
  \item \textbf{Yd}：一次绕组为星形连接，3 次谐波电流不能流通，一次绕组中励磁电流接近正弦波，故铁芯磁路中的主磁通含有 3 次谐波分量；该分量在二次绕组中感应 3 次谐波相电动势 $e_{23}$，而二次绕组为三角形连接，三相的 $e_{23}$ 将在闭合三角形绕组中产生 3 次谐波\textbf{环流} $i_{23}$。
\end{itemize}
由于一次绕组为星形连接、3 次谐波电流不能流通，不可能产生 3 次谐波磁动势与 $N_2\dot I_{23}$ 相平衡，故二次绕组中的环流就起到了励磁电流的作用。磁路中的主磁通由与它相交链的\textbf{合成磁动势}产生，而合成磁动势不必非由一个绕组提供；产生正弦主磁通所必须的基波励磁电流与 3 次谐波电流由一、二次绕组共同提供，与单独由一次绕组提供效果相同。因此 Yd 与 Dy 效果一样，主磁通和每相绕组感应电动势仍接近正弦波形。

三角形联结组中的 3 次谐波环流甚小，对运行影响不大，且能改善主磁通和感应电动势波形，因此较大容量的变压器常将其中一侧绕组接成三角形。

\begin{center}
\begin{tabular}{@{}p{1.5cm}p{2.5cm}p{3.4cm}p{4.2cm}@{}}
\toprule
联结组 & 3 次谐波电流 & 3 次谐波磁通 & 相电动势波形 \\
\midrule
YNy & 经中性线流通 & 很小 & 接近正弦 \\
Yy（组式） & 不能流通 & 可在铁芯内闭合 & 尖顶波（$e_3$ 达 45\%～60\%），不宜采用 \\
Yy（心式） & 不能流通 & 受抑制，经漏磁路闭合 & 接近正弦，但油箱附加损耗，容量 $\le1800$\,kV·A \\
Dy & 一次 D 内流通 & 很小 & 正弦 \\
Yd & 二次 D 内环流 & 由两侧合成磁动势建立 & 正弦，大容量常用 \\
\bottomrule
\end{tabular}
\end{center}

\corepoint{这一节的判断链只有三步：① 有没有三角形接法或中性线 $\Rightarrow$ 3 次谐波电流能否流通；② 磁路是组式（独立）还是心式（联系）$\Rightarrow$ 3 次谐波磁通能否在铁芯内闭合；③ 3 次谐波磁通大不大 $\Rightarrow$ 相电动势是否尖顶。记住"线电动势中 3 次谐波互相抵消、永远正弦"，这是常考的一句话陷阱。}

\section{标幺值}

变压器和电机中，电压、电流、阻抗和功率等物理量可用实际值，也常用\textbf{标幺值}表示与计算。标幺值也称相对值，是\textbf{同单位}的实际物理量与其基值之比：
\begin{equation}
\text{标幺值}=\frac{\text{实际值}}{\text{基值}} .
\end{equation}
本书中基值用下角标"b"表示，标幺值用右上角标"*"表示。标幺值乘以 100\% 便得百分值。

\subsection{基值}
应用标幺值时先要选定基值。变压器和电机的工程计算中，通常选定额定值作为基值。

对单相系统：
\begin{equation}
S_b=P_b=Q_b=U_bI_b=U_NI_N=S_N,\qquad
Z_b=R_b=X_b=\frac{U_N}{I_N}=\frac{U_N^{2}}{S_N}.
\end{equation}
对三相系统：
\begin{equation}
U_{bL}=U_{NL}=U_N,\qquad
S_b=P_b=Q_b=\sqrt3U_{NL}I_{NL}=S_N,\qquad
Z_b=\frac{U_{N\phi}}{I_{N\phi}}=\frac{U_{NL}^{2}}{S_N},
\end{equation}
其中下标 $\phi$ 表示相值、$L$ 表示线值；$Z_b$ 在星形接法时为 $U_{NL}^{2}/S_N$，三角形接法时相应取相值形式。

将各物理量的实际值与对应基值之比，即得各物理量的标幺值。注意：\textbf{一次绕组和二次绕组的基值分别选取各自的基值}。例如单相变压器一次绕组电压和电流的标幺值为 $U_1^{*}=U_1/U_{1N}$、$I_1^{*}=I_1/I_{1N}$；三相变压器一次绕组线电压的标幺值为 $U_{1L}^{*}=U_{1L}/U_{1N}$。

\subsection{标幺值的特点}

\subsubsection{额定电压、电流和容量的标幺值为 1}
由于电压、电流和容量的基值都选为额定值，故
\begin{equation}
U_N^{*}=I_N^{*}=S_N^{*}=1 .
\end{equation}

\subsubsection{在三相系统中，相、线标幺值相等}
无论电压还是电流，其相、线标幺值相等。以星形接法为例，
\begin{equation}
U_L^{*}=\frac{U_L}{U_N}=\frac{\sqrt3U_\phi}{\sqrt3U_{N\phi}}=\frac{U_\phi}{U_{N\phi}}=U_\phi^{*},\qquad
I_L^{*}=I_\phi^{*}.
\end{equation}
这一条极其实用：\textbf{用标幺值计算时不必区分相值与线值}，也就不用反复乘除 $\sqrt3$。

\subsubsection{单、三相功率计算公式相同}
用标幺值表示时，单相与三相系统的视在功率、有功功率和无功功率表达式相同，且不存在"线"的概念：
\begin{equation}
S^{*}=U^{*}I^{*},\qquad P^{*}=U^{*}I^{*}\cos\varphi,\qquad Q^{*}=U^{*}I^{*}\sin\varphi,
\end{equation}
式中 $\varphi$ 为功率因数角。

\subsubsection{归算前后的值相等}
以二次绕组电压和阻抗为例：
\begin{equation}
U_2'^{*}=\frac{U_2'}{U_{1N}}=\frac{kU_2}{kU_{2N}}=\frac{U_2}{U_{2N}}=U_2^{*},\qquad
Z_2'^{*}=\frac{k^{2}Z_2}{k^{2}Z_{2b}}=Z_2^{*}.
\end{equation}
即\textbf{用标幺值表示时，归算前后的值相等}——这意味着用标幺值计算可以\textbf{不必再进行归算}，这是标幺值最大的便利。

\subsubsection{变压器基本方程组的表达形式不变}
用标幺值表示，变压器基本方程组的表达形式保持不变，并且归算值上角标的"$'$"可以保留也可以去掉。

\subsubsection{便于比较不同容量的电机}
对于不同容量的变压器或电机，用标幺值表示时各参数和典型性能数据通常都在一定范围内，便于比较和分析。例如电力变压器漏阻抗的标幺值约为 4\%～8\%，空载电流的标幺值约为 1\%～5\%。

\corepoint{标幺值最容易错的三处：① 基值必须"成对"选取——$U_b$ 与 $I_b$ 要满足 $U_bI_b=S_b$，只改一个就会使阻抗基值错；② 一次侧、二次侧各有自己的基值，$U_{2b}=U_{2N}$（不是 $kU_{1N}$ 的换算值），这样才能自动消去 $k$；③ 三相系统选\textbf{线值}作基值时，$Z_b=U_{NL}^2/S_N$；若选相值作基值，$Z_b=U_{N\phi}^2/(S_N/3)=U_{N\phi}/I_{N\phi}$，两者结果相同，但不能混着用。}

% ============================================================
% 第 6 章 变压器的运行（教材第二篇，印刷页 60--67）
% 节名：第一节 变压器的电压变化率和效率；第二节 变压器的并联运行
% ============================================================
\chapter{变压器的运行}

电压变化率和效率是变压器的主要运行性能指标；并联运行则是电力系统中变压器最常见的运行方式。本章把前面建立的等值电路与参数，用到"输出电压稳不稳、运行经不经济、几台能不能一起带负载"这三个工程问题上。

\section{变压器的电压变化率和效率}

\subsection{电压变化率}
将变压器一次绕组外加额定电压，二次绕组空载时端电压为空载电压，且 $U_{20}=U_{2N}$。接上负载以后，变压器内部的漏阻抗产生电阻和漏抗压降，使二次绕组端电压随负载电流变化。

在\textbf{一次绕组端电压保持为额定、负载功率因数为常值}时，空载与负载时二次绕组相电压的代数差 $U_{20}-U_2$ 与二次绕组额定相电压 $U_{2N}$ 的比值，称为变压器的\textbf{电压变化率}，用 $\Delta U$ 表示：
\begin{equation}
\Delta U=\frac{U_{20}-U_2}{U_{2N}}\times100\%=\frac{U_{2N}-U_2}{U_{2N}}\times100\% ,
\end{equation}
其标幺值形式为 $\Delta U^{*}=1-U_2^{*}$（因 $U_{20}^{*}=1$）。

$\Delta U$ 与变压器的参数和负载性质有关。工程上通常应用\textbf{简化等效电路}及其相量图来计算。由于短路阻抗 $Z_k$ 很小，$\dot U_1$ 与 $\dot U_2'$ 之间的夹角 $\theta$ 很小，故有
\begin{equation}
\Delta U^{*}\approx I_2^{*}\left(R_k^{*}\cos\varphi_2+X_k^{*}\sin\varphi_2\right).
\end{equation}

\keyformula{\Delta U^{*}\approx I_2^{*}\left(R_k^{*}\cos\varphi_2+X_k^{*}\sin\varphi_2\right)}

该式表明：电压变化率既与负载电流和功率因数角有关，也与变压器的短路阻抗有关。

\begin{itemize}
  \item \textbf{感性负载}：$\varphi_2>0$，$\sin\varphi_2>0$、$\cos\varphi_2>0$，$\Delta U^{*}$ 为正值。负载电流越大，二次端电压越低。
  \item \textbf{容性负载}：$\varphi_2<0$，$\sin\varphi_2<0$，$\Delta U^{*}$ 可能为正也可能为负。当为负值时，随负载电流增加，二次端电压不但不下降反而\textbf{升高}。
\end{itemize}

当二次负载为额定值（$I_2^{*}=1$）时，
\begin{equation}
\Delta U_N^{*}=R_k^{*}\cos\varphi_N+X_k^{*}\sin\varphi_N .
\end{equation}
电压变化率的大小反映变压器负载时的供电质量：$\Delta U$ 越小，输出电压越平稳。通常在 $\cos\varphi_2=0.8$（滞后）、$I_2=I_{2N}$ 时，电压变化率为 3\%～6\%。

\subsection{效率}
变压器运行时内部产生损耗，可分为两类：
\begin{itemize}
  \item \textbf{铁耗} $p_{Fe}$：主磁通在铁芯中产生的损耗。由 $P_{Fe}\propto B_m^2$，在已制成的变压器中近似正比于 $U_1^2$。由于一次绕组端电压通常保持 $U_1=U_{1N}$ 不变，故铁耗又称\textbf{不变损耗}，可由空载试验确定（$p_{Fe}\approx p_0$）。
  \item \textbf{铜耗} $p_{Cu}$：电流通过一、二次绕组产生的电阻损耗。
  \begin{equation}
  p_{Cu}=mI_1^{2}R_1+mI_2'^{2}R_2'=mI_2'^{2}R_k' .
  \end{equation}
  由于铜耗随负载电流变化，故称\textbf{可变损耗}，可由短路试验确定（$p_{Cu}\approx\beta^2p_{kN}$，$\beta=I_2/I_{2N}$ 为负载系数）。
\end{itemize}

考虑内部损耗后，输出功率 $P_2$ 与输入功率 $P_1$ 的关系为 $P_1=P_2+p_{Cu}+p_{Fe}$，其中 $p=p_{Cu}+p_{Fe}$ 为总损耗。变压器的输出功率为
\begin{equation}
P_2=mU_2I_2\cos\varphi_2 .
\end{equation}
由于电压变化率很小，若忽略二次端电压的变化，认为 $U_2\approx U_{20}$，则 $P_2\approx mU_{20}I_2\cos\varphi_2$，故输入功率
\begin{equation}
P_1\approx mU_{20}I_2\cos\varphi_2+p_{Fe}+mI_2'^{2}R_k' .
\end{equation}

变压器的效率是输出功率与输入功率之比：
\begin{equation}
\eta=\frac{P_2}{P_1}\times100\%=\frac{mU_{20}I_2\cos\varphi_2}{mU_{20}I_2\cos\varphi_2+p_{Fe}+mI_2'^{2}R_k'}\times100\% .
\end{equation}
用额定容量与负载系数表示，即得工程上最常用的效率公式：
\begin{equation}
\eta=\frac{S_N\beta\cos\varphi_2}{S_N\beta\cos\varphi_2+p_0+\beta^{2}p_{kN}}\times100\% .
\end{equation}

\keyformula{\eta=\frac{\beta S_N\cos\varphi_2}{\beta S_N\cos\varphi_2+p_0+\beta^{2}p_{kN}}\times100\%}

保持负载功率因数不变，效率随负载电流 $I_2$ 的变化曲线为一条先升后降的曲线，在某一负载电流时效率达到最大。将上式对 $I_2$（或 $\beta$）求导并令其等于零，得效率最大的条件：
\begin{equation}
p_{Fe}=mI_2'^{2}R_k'=p_{Cu}
\quad\Longleftrightarrow\quad
p_0=\beta^{2}p_{kN}
\quad\Longleftrightarrow\quad
\beta=\sqrt{\frac{p_0}{p_{kN}}} .
\end{equation}

\keyformula{p_{Fe}=p_{Cu}\ \Longleftrightarrow\ \beta_m=\sqrt{\frac{p_0}{p_{kN}}}}

即\textbf{当铜耗等于铁耗（可变损耗等于不变损耗）时，变压器效率达到最大值}。实际电力变压器中最大效率发生在 $\beta=0.5\sim0.6$ 之间——这是因为 $p_0<p_{kN}$，$\beta_m<1$，故意设计成"常在半载附近运行最经济"。

变压器的效率虽然可以直接测量，但电力变压器效率很高（通常在 95\% 以上，大容量可达 99\% 以上），$P_1$ 与 $P_2$ 差距很小，直接法误差大，且大容量变压器难以找到相适应的负载。故工程上常通过\textbf{空载和短路试验}求出铜耗和铁耗，再用\textbf{间接法}（损耗分析法）计算效率。

\begin{center}
\begin{tabular}{@{}p{1.9cm}p{3.2cm}p{4.0cm}p{3.9cm}@{}}
\toprule
项目 & 定义/公式 & 性质 & 测定方法 \\
\midrule
电压变化率 & $\Delta U^{*}\approx I_2^{*}(R_k^{*}\cos\varphi_2+X_k^{*}\sin\varphi_2)$ & 感性为正、容性可为负 & 由短路参数计算 \\
铁耗 $p_{Fe}$ & $\approx p_0$，$\propto U_1^2$ & 不变损耗 & 空载试验 \\
铜耗 $p_{Cu}$ & $\approx\beta^2p_{kN}$，$\propto I^2$ & 可变损耗 & 短路试验 \\
效率 $\eta$ & $\dfrac{\beta S_N\cos\varphi_2}{\beta S_N\cos\varphi_2+p_0+\beta^2p_{kN}}$ & 最大时 $p_{Fe}=p_{Cu}$ & 间接法（损耗分析） \\
\bottomrule
\end{tabular}
\end{center}

\corepoint{三个高频错点：① 把 $\Delta U$ 的分子写成 $U_{1N}-U_2$（应是与 $U_{2N}$ 比较，且是相电压之差）；② 容性负载时忘记 $\sin\varphi_2$ 取负号，得出"容性负载端电压下降"的错误结论；③ 求最大效率时用 $\beta=1$ 代入——最大效率条件是 $p_0=\beta^2p_{kN}$，不是满载。另外注意 $p_0$、$p_{kN}$ 都是"三相总损耗"，用标幺值时对应 $S_N$，不要漏乘相数 $m$。}

\section{变压器的并联运行}

所谓变压器的并联运行，是指将两台或多台变压器的一次绕组和二次绕组分别接到各自的公共母线上，共同向负载供电。并联运行可以提高供电可靠性，减小总的备用容量，并可根据负载变化调整投入运行的变压器台数，以提高运行效率。

\subsection{理想并联运行的条件}
变压器并联运行时，总希望损耗小、效率高、每台变压器都能得到充分利用，即达到理想状态：
\begin{enumerate}
  \item \textbf{空载时并联的变压器之间无环流}。环流会增加损耗，严重时使变压器过热甚至烧坏。
  \item \textbf{负载时各变压器有相同的负载率}。这样容量大的变压器承担负载电流大、容量小的承担小，各变压器容量得到最充分利用。
  \item \textbf{各变压器输出电流同相位}。这样当总负载电流一定时，各变压器分担的电流最小，总铜耗最小。
\end{enumerate}

要达到上述理想要求，并联运行的各台变压器应满足如下条件：
\begin{enumerate}
  \item 各台变压器的\textbf{额定电压与变比相等}；
  \item 各台变压器的\textbf{联结组别相同}；
  \item 各台变压器的\textbf{短路阻抗标幺值相等，阻抗角相同}。
\end{enumerate}

\subsection{变比不相等时的空载环流}
以两台变压器并联为例，设联结组相同。由于两台变压器接到公共母线上，一次绕组具有相同的电网电压 $\dot U_1$；当变比不等时，两台变压器二次绕组的感应电动势不相等。设第一台为 $\dot E_{2I}$、第二台为 $\dot E_{2II}$，公共母线把两台变压器的二次绕组构成回路。因此，即使二次母线不接负载，在电动势差 $\Delta\dot E_2=\dot E_{2I}-\dot E_{2II}$ 作用下，二次绕组也会产生\textbf{环流}。由磁平衡关系，二次绕组产生环流的同时一次绕组也产生环流；由于变比不等，两台变压器一次绕组的环流不相同。

应用简化等效电路，可求出变比不等时的空载环流：
\begin{equation}
\dot I_c=\frac{\Delta\dot E_2}{Z_{kI}'+Z_{kII}'} .
\end{equation}
由于变压器的短路阻抗很小，\textbf{即使变比相差很小也会产生较大的环流}，因此对并联变压器之间的变比差必须严格限制（工程上一般要求变比差不超过 $\pm0.5\%$）。

\subsection{联结组别不同时的空载环流}
若并联运行的两台变压器联结组别不同，后果更为严重：联结组别不同的两台变压器，线电动势的相位至少相差 $30^\circ$。设两台变压器线电动势大小相等，则二次绕组电动势差为
\begin{equation}
\Delta E=2E_{2}\sin\frac{30^\circ}{2}=0.52E_2 .
\end{equation}
这么大的电动势差作用在二次绕组回路中，不但会产生很大的空载环流，甚至会烧坏变压器。故\textbf{联结组不同的变压器绝对不能并联运行}。

\subsection{短路阻抗不等时的负载分配}
分析负载分配时假设两台变压器的变比和联结组别都相同，则变压器内不产生环流，两台变压器具有相同的输入电压 $\dot U_1$ 和输出电压 $\dot U_2$。由简化等效电路，两台变压器并联时短路阻抗上的压降应相等：
\begin{equation}
\dot I_I Z_{kI}=\dot I_{II}Z_{kII}
\quad\Longrightarrow\quad
\frac{\dot I_I}{\dot I_{II}}=\frac{Z_{kII}}{Z_{kI}} .
\end{equation}
即\textbf{并联运行的各台变压器所分担的负载电流与其短路阻抗成反比}。

两边同乘以额定电流并考虑两台变压器额定电压相同，可导出用标幺值表示的负载电流分配：
\begin{equation}
\frac{I_I^{*}}{I_{II}^{*}}=\frac{Z_{kII}^{*}}{Z_{kI}^{*}} .
\end{equation}

\keyformula{\frac{I_I^{*}}{I_{II}^{*}}=\frac{Z_{kII}^{*}}{Z_{kI}^{*}}}

由于两台变压器的额定电压和负载电压均相等，所以有
\begin{equation}
\frac{S_I}{S_{II}}=\frac{S_{NI}}{S_{NII}}\cdot\frac{Z_{kII}^{*}}{Z_{kI}^{*}} .
\end{equation}
可见：\textbf{短路阻抗标幺值不相等时，并联各变压器所分担的负载不能按额定容量成正比例分配}。由标幺值式还可看出，为满足理想并联运行的第三条（输出电流同相位），各台变压器短路阻抗的\textbf{阻抗角也应相等}。

实际并联运行时：联结组必须相同；变比差要严格控制；漏阻抗标幺值在数值上相差不要太大（通常相对误差小于 10\%），相角可有一定差别，通常不作严格要求。

\subsection{容量利用率}
若两台变压器短路阻抗标幺值不等，则阻抗小的一台先达到满载，另一台仍欠载，于是总设备容量不能充分利用。设负载按总容量 $S_L$ 分配，为使任何一台都不过载，应取"先满载者"为约束，可承担的最大负载与设备总容量之比即\textbf{容量利用率}。

\begin{center}
\begin{tabular}{@{}p{2.5cm}p{3.6cm}p{4.4cm}p{2.6cm}@{}}
\toprule
条件不满足 & 后果 & 定量关系 & 严重程度 \\
\midrule
变比不等 & 空载环流 & $I_c=\Delta E_2/(Z_{kI}'+Z_{kII}')$ & 须限制变比差 \\
联结组不同 & 巨大空载环流，可烧毁 & $\Delta E=0.52E_2$（差 $30^\circ$） & \textbf{绝对禁止}并联 \\
短路阻抗标幺值不等 & 负载不能按容量比例分配 & $I_I^{*}/I_{II}^{*}=Z_{kII}^{*}/Z_{kI}^{*}$ & 影响经济性 \\
阻抗角不等 & 电流不同相位，总电流偏大 & 相量和大于代数和 & 一般不作严格要求 \\
\bottomrule
\end{tabular}
\end{center}

\corepoint{把三条条件分成"能不能"和"划不划算"两类，记忆最省力：\textbf{变比相等 + 组别相同}是"能不能并联"的前提（否则空载就有环流）；\textbf{短路阻抗标幺值相等 + 阻抗角相同}是"并得好不好"的条件（否则负载分不匀、总铜耗偏大）。考试若问"两台变压器能否并联"，先答组别，再答变比，最后答阻抗。}

% ============================================================
% 第 7 章 其他常用变压器（教材第二篇，印刷页 68--75）
% 节名：第一节 自耦变压器；第二节 三绕组变压器；第三节 互感器
% ============================================================
\chapter{其他常用变压器}

前面各章讨论的变压器，每相只有一个一次绕组和一个二次绕组，称为\textbf{双绕组变压器}——它是最普通、最常用的变压器。除此之外还有许多型式和用途的变压器，本章介绍较常用的自耦变压器、三绕组变压器和互感器。它们的基本原理仍是双绕组变压器那一套，特殊之处在于\textbf{绕组的连接方式}或\textbf{绕组的数目}改变了功率传递的路径。

\section{自耦变压器}

\subsection{结构特点和变比}
普通双绕组变压器的一、二次绕组之间相互绝缘，只有磁的耦合，没有直接电的联系。如果把普通双绕组变压器的一次绕组和二次绕组\textbf{串联}起来作为新的一次绕组，而原二次绕组还同时作为新的二次绕组，便得到一台降压自耦变压器。

设接成普通双绕组时一、二次绕组匝数分别为 $N_1$、$N_2$；改为自耦连接时，一次绕组匝数 $N_{1a}=N_1+N_2$，二次绕组匝数 $N_{2a}=N_2$。当绕组 $Ax$ 两端外加交变正弦电压 $U_{1a}$ 时，铁芯内产生交变磁通，在 $Ax$ 间感应电动势
\begin{equation}
E_{1a}=4.44f\left(N_1+N_2\right)\Phi_m=E_1+E_2 ,
\end{equation}
在二次绕组 $ax$ 间感应电动势 $E_{2a}=4.44fN_2\Phi_m=E_2$。忽略漏阻抗压降，自耦变压器的变比为
\begin{equation}
k_a=\frac{N_{1a}}{N_{2a}}=\frac{N_1+N_2}{N_2}=1+\frac{N_1}{N_2}=1+\frac{1}{k},
\end{equation}
其中 $k=N_1/N_2$ 为接成双绕组变压器时的变比。

可见自耦变压器的变压原理与双绕组变压器相同，只是由于一、二次绕组有直接的电联系、一次绕组包含了二次绕组，因此自耦变压器的变比 $k_a$ 比接成双绕组时的变比 $k$ 大。用自耦变压器升压时，可把电源加在匝数为 $N_2$ 的绕组上，匝数为 $N_{1a}$ 的绕组感应电动势就比外加电压高。

\subsection{基本电磁关系}
\subsubsection{电压、电流关系}
对降压自耦变压器，二次绕组（$ax$ 段）的电流为
\begin{equation}
\dot I_{2}=\dot I_1+\dot I_{2a},
\end{equation}
即在忽略励磁电流的情况下，由磁平衡关系 $N_1\dot I_1+N_2\dot I_2=0$（对串联段而言），可得公共绕组中的电流为两电流之差。故自耦变压器二次电流为"一次电流与公共段电流的相量和"，公共绕组段的电流较小——这正是自耦变压器省材料的根本原因。

\subsubsection{功率关系}
作为普通双绕组变压器时，额定容量为
\begin{equation}
S_N=U_{1N}I_{1N}=U_{2N}I_{2N}.
\end{equation}
接成自耦变压器时，额定容量为
\begin{equation}
S_{aN}=U_{1aN}I_{1aN}=S_N+\frac{U_{2N}}{U_{1N}}S_N=S_N+S_Q ,
\end{equation}
即
\begin{equation}
S_{aN}=S_N+S_Q .
\end{equation}

\keyformula{S_{aN}=S_N+S_Q,\qquad S_N=\left(1-\frac{1}{k_a}\right)S_{aN}}

自耦变压器的视在功率由两部分组成：
\begin{itemize}
  \item \textbf{感应功率} $S_N$：借助于 $Aa$ 段绕组和 $ax$ 段绕组之间的电磁感应关系，从一次侧传递到二次侧的功率。这部分功率与普通双绕组变压器相同，变压器的\textbf{主要尺寸和材料消耗由这部分功率决定}，所以 $S_N$ 称为\textbf{计算容量}。
  \item \textbf{传导功率} $S_Q$：通过电路的直接联系，从一次绕组直接传递到负载侧的功率。它\textbf{不需要增加变压器的计算容量}。
\end{itemize}

\subsection{自耦变压器的优缺点}
\begin{enumerate}
  \item \textbf{体积小、质量轻、价格低}，便于运输和安装。自耦变压器的计算容量 $S_N=(1-1/k_a)S_{aN}<S_{aN}$，故在同样额定容量下，主要尺寸和所用铜线、硅钢片及其他结构材料都比较少；铜线和硅钢片减少，铜耗和铁耗也相应减少，效率提高。$k_a$ 越接近 1，计算容量占额定容量的比例越小，这些优点越显著。因此自耦变压器适合于一、二次电压相差不大的场合，一般 $k_a<2$。
  \item 将低压侧短路、从高压侧看，其\textbf{短路阻抗与普通双绕组变压器相同}，但由于自耦变压器的电压基值较高，故\textbf{短路阻抗标幺值比普通双绕组变压器小，短路电流较大}。
  \item 高、低压绕组之间具有\textbf{直接的电联系}：一方面要求低压绕组具有与高压绕组相同的绝缘水平，另一方面使\textbf{过电压保护变得复杂}。
\end{enumerate}

\begin{center}
\begin{tabular}{@{}p{2.6cm}p{5.0cm}p{5.0cm}@{}}
\toprule
项目 & 双绕组变压器 & 自耦变压器 \\
\midrule
电的联系 & 只有磁耦合，无直接电联系 & 有直接电联系（一次含二次） \\
变比 & $k=N_1/N_2$ & $k_a=1+1/k=N_{1a}/N_{2a}$ \\
额定容量 & $S_N=U_{1N}I_{1N}$ & $S_{aN}=S_N+S_Q$ \\
计算容量 & 等于额定容量 & $(1-1/k_a)S_{aN}<S_{aN}$ \\
短路电流 & 较小 & 标幺短路阻抗小，短路电流较大 \\
绝缘与保护 & 简单 & 低压侧需按高压绝缘；过电压保护复杂 \\
适用场合 & 通用 & $k_a<2$，电压相近的场合 \\
\bottomrule
\end{tabular}
\end{center}

\corepoint{自耦变压器的全部"异常"都源于一句话：\textbf{一次绕组包含了二次绕组}。由此推出三件事——变比 $k_a=1+1/k$ 而不是 $k$；容量由"感应功率 + 传导功率"两部分构成，故计算容量小、省材料；高低压直接电联系，故低压侧要按高压绝缘、短路电流大、过电压保护复杂。答题时先写这句，后面自然出来。}

\section{三绕组变压器}

三绕组变压器有高压、中压和低压三个绕组，大多用于二次侧需要两种不同电压等级的电力系统。三绕组变压器相当于不同变比的两台双绕组变压器，但使用三绕组变压器比用两台更经济。

\subsection{结构特点和容量分配}
三个绕组一般均套装在同一铁芯柱上。为便于绝缘，\textbf{高压绕组通常放置于最外层}。对于\textbf{降压}变压器，中压绕组放在中间层，低压绕组放在最里层；对于\textbf{升压}变压器，为使漏磁场分布均匀、漏抗分配合理，不致因低压和高压绕组相距较远而造成漏磁通及其附加损耗增加，保证有较好的电压变化率和运行性能，通常将中压绕组放在最里层、低压绕组放在中间层。

GB 1094.1—1996 规定，三相三绕组变压器的标准联结组有 \textbf{YNyn0d11} 和 \textbf{YNyn0y0} 两种。

三个绕组的容量可以相同也可以不同。三绕组变压器的\textbf{额定容量是指三个绕组中容量最大那个绕组的容量}。若以额定容量为 100\%，则三个绕组容量的配合有 100/100/50、100/50/100 和 100/100/100 三种。

\subsection{基本方程式和等效电路}
设绕组 1 是一次绕组，另外两个是二次绕组，则三绕组变压器的变比分别为
\begin{equation}
k_{12}=\frac{N_1}{N_2},\qquad k_{13}=\frac{N_1}{N_3}.
\end{equation}

三绕组变压器在磁路方面有三个绕组相互耦合，它们产生的磁通可分为三种：
\begin{itemize}
  \item \textbf{自漏磁通}：只与一个绕组相交链；
  \item \textbf{互漏磁通}：与两个绕组交链而不与第三个绕组相交链；
  \item 以上两部分漏磁通通过油或空气闭合，\textbf{不受铁芯磁路饱和的影响}；
  \item \textbf{主磁通}：与三个绕组同时交链，通过铁芯磁路闭合，\textbf{受铁芯饱和影响}。
\end{itemize}

主磁通由三个绕组的合成磁动势产生，故励磁磁动势为
\begin{equation}
N_1\dot I_m=N_1\dot I_1+N_2\dot I_2+N_3\dot I_3
\quad\Longrightarrow\quad
\dot I_1+\dot I_2'+\dot I_3'=\dot I_m ,
\end{equation}
其中 $\dot I_2'$、$\dot I_3'$ 分别表示归算到一次侧的两个二次绕组电流。忽略励磁电流时
\begin{equation}
\dot I_1+\dot I_2'+\dot I_3'=0 .
\end{equation}

由于存在互漏磁通，每个绕组内的漏磁压降都受另外两个绕组的影响。若按双绕组变压器的分析方法建立等效电路，三绕组变压器每个绕组的电抗必须是一个既包含自感电抗、又包含互感电抗的\textbf{等效值}。用 $X_1$、$X_2'$、$X_3'$ 和 $Z_1$、$Z_2'$、$Z_3'$ 表示三个绕组的等效电抗和等效阻抗，可列出电压方程
\begin{equation}
\dot U_1=-\dot E_1+\dot I_1Z_1,\qquad
\dot U_2'=\dot E_1-\dot I_2'Z_2',\qquad
\dot U_3'=\dot E_1-\dot I_3'Z_3' .
\end{equation}
据此可得三绕组变压器的\textbf{简化等效电路}：三个等效阻抗呈星形连接，公共点接励磁支路。

\subsection{参数测定}
简化等效电路中的 $Z_1$、$Z_2'$、$Z_3'$ 可通过\textbf{三次短路试验}确定：
\begin{enumerate}
  \item 绕组 2 短路、绕组 3 开路，在绕组 1 上加低电压做短路试验，测得 1、2 之间的短路阻抗 $Z_{k12}=Z_1+Z_2'$；
  \item 绕组 3 短路、绕组 2 开路，测得 $Z_{k13}=Z_1+Z_3'$；
  \item 绕组 3 短路、绕组 1 开路，测得 $Z_{k23}=Z_2'+Z_3'$。
\end{enumerate}
联解三式得
\begin{equation}
Z_1=\tfrac12\left(Z_{k12}+Z_{k13}-Z_{k23}'\right),\quad
Z_2'=\tfrac12\left(Z_{k12}+Z_{k23}'-Z_{k13}\right),\quad
Z_3'=\tfrac12\left(Z_{k13}+Z_{k23}'-Z_{k12}\right).
\end{equation}
其中 $Z_1$、$Z_2'$、$Z_3'$ 均为等效阻抗，其大小与三个绕组之间的布置有关；在某种特定布置下，其中某个等效电抗可能变成负值，此时该电抗可当作\textbf{容抗}计算。

\corepoint{三绕组变压器的等效电路不是"三个绕组各串一个阻抗再并到励磁支路"，而是\textbf{三个等效阻抗星形连接后接励磁支路}——因为互漏磁通使每个绕组的漏抗不再是独立的物理量，只能由三次短路试验"解方程"求出。若算出某个 $X$ 为负值不要慌，那是互漏磁通造成的等效结果，按容抗处理即可。}

\section{互感器}

互感器是测量用设备，分\textbf{电压互感器}和\textbf{电流互感器}两种，主要在高电压、大电流测量以及各种继电保护中使用。互感器的工作原理与普通变压器相同，但运行状态与使用禁忌截然相反——这是本节的核心。

\subsection{电压互感器}
电压互感器一次绕组匝数 $N_1$ 较多，直接接于被测的高压电路上；二次绕组匝数较少，一般只有一匝或几匝，与电压表或其他仪表的电压线圈相接。由于电压表或仪表的电压线圈内阻抗很大，所以电压互感器工作时相当于变压器的\textbf{空载运行}。

电压互感器二次绕组的额定电压通常取 100\,V。忽略漏阻抗压降，则
\begin{equation}
\frac{U_1}{U_2}=\frac{N_1}{N_2}=k_u\quad\Longrightarrow\quad U_1=k_uU_2 ,
\end{equation}
因此利用一、二次绕组不同的匝数比可将电路上的高电压变为低电压来测量。

\subsubsection{误差与准确级}
电压互感器的准确度由\textbf{变比误差}和\textbf{相位误差}衡量：
\begin{itemize}
  \item 变比误差：一次电压 $U_1$ 与二次电压归算值 $k_uU_2$ 的算术差的相对值，即 $\left[(U_1-k_uU_2)/U_1\right]\times100\%$；
  \item 相角误差：$\dot U_1$ 与 $-k_u\dot U_2$ 之间的相位差。
\end{itemize}
从变压器相量图可以看出，\textbf{励磁电流和一、二次绕组的漏阻抗压降}是产生这两种误差的主要原因。因此在设计电压互感器时，大都采用高磁导率的硅钢片，尽量减小磁路中的气隙，并使铁芯不饱和，一般取磁密为 0.6～0.8\,T。按误差大小，电压互感器的准确度分为 0.1、0.2、0.5、1.0 和 3.0 五级。

\subsubsection{使用注意}
\begin{enumerate}
  \item \textbf{二次绕组不允许短路}，否则会产生很大的短路电流；
  \item 互感器的\textbf{铁芯和二次绕组的一端必须可靠接地}，以保证操作人员的安全。
\end{enumerate}

\subsection{电流互感器}
电流互感器一次绕组匝数 $N_1$ 较少，通常为一匝或几匝，串入需要测量电流的电路中；二次绕组匝数 $N_2$ 较多，与电流表或其他仪表的电流线圈相接。电流表或仪表中电流线圈的阻抗很小，所以电流互感器工作时相当于普通变压器的\textbf{短路运行}。

电流互感器二次绕组的额定电流一般为 5\,A 或 1\,A。忽略励磁电流，根据磁平衡方程式 $I_1N_1+I_2N_2=0$，在数值上可得
\begin{equation}
I_1=\frac{N_2}{N_1}I_2=k_iI_2 ,
\end{equation}
即用匝数比把大电流变成小电流来测量。

\subsubsection{误差与准确级}
与电压互感器类似，电流互感器的误差也分变比误差与相角误差，主要由励磁电流和漏阻抗压降引起。为减小误差，应采用高磁导率材料、减小磁路气隙、降低铁芯磁密（即增大铁芯截面）。按误差大小分为 0.1、0.2、0.5、1.0、3.0 等级。

\subsubsection{使用注意}
\begin{enumerate}
  \item \textbf{二次绕组绝对不允许开路}。二次开路时，一次电流全部变为励磁电流，铁芯中磁通和磁密急剧增大，铁耗剧增、铁芯严重过热；同时在二次绕组两端感应出很高的电压，危及人身与设备安全。
  \item \textbf{铁芯和二次绕组的一端必须可靠接地}。
\end{enumerate}

\begin{center}
\begin{tabular}{@{}p{2.3cm}p{3.3cm}p{3.6cm}p{4.0cm}@{}}
\toprule
项目 & 电压互感器 & 电流互感器 & 记忆要点 \\
\midrule
一次绕组 & 匝数多，并联在被测电路 & 匝数少（1～几匝），串入电路 & 并联/串联 \\
工作状态 & 相当于空载运行 & 相当于短路运行 & 与变压器对应 \\
二次额定值 & 100\,V & 5\,A 或 1\,A & 标准化量程 \\
主要误差源 & 励磁电流、漏阻抗压降 & 励磁电流、漏阻抗压降 & 都靠高磁导率、低磁密减小 \\
使用禁忌 & \textbf{二次不许短路} & \textbf{二次不许开路} & 一字之差，后果都严重 \\
安全措施 & 铁芯与二次一端接地 & 铁芯与二次一端接地 & 防止高压窜入低压 \\
\bottomrule
\end{tabular}
\end{center}

\corepoint{记住这一组"反着来"的关系：电压互感器 = \textbf{空载}运行，因此怕\textbf{短路}（短路即成为短路运行，电流剧增）；电流互感器 = \textbf{短路}运行，因此怕\textbf{开路}（开路即失去去磁作用，铁芯饱和、二次出现高压）。两条安全措施完全相同：铁芯与二次一端必须接地。}

% ============================================================
% 第 8 章 交流绕组（教材第三篇，印刷页 78--89）
% 节名：第一节 同步发电机和感应电动机工作原理概述；第二节 三相单层绕组；第三节 三相双层绕组
% ============================================================
\chapter{交流绕组}

工业用交流电机主要分为两大类：\textbf{同步电机}与\textbf{感应电机}。两者工作原理和运行特点不同，但有一批共性问题：定子绕组都是三相交流绕组，绕组中感应电动势与磁动势的分析计算方法相同，发热、冷却与温升限值的规律也相同。本篇（第 8—11 章）解决这些共性问题，第 12 章起再分别研究两类电机的运行性能。

\section{同步发电机和感应电动机工作原理概述}

\subsection{同步发电机的基本工作原理}
三相同步发电机由\textbf{定子}、\textbf{转子}和定转子之间的工作气隙组成。定子是一个圆筒形铁芯，靠近铁芯内表面均匀开槽，槽里嵌放导体；把这些导体按一定规律连接起来构成\textbf{定子绕组}（也称电枢绕组）。最简单的三相定子绕组每相仅有一个线圈，首末端分别用 $A$—$X$、$B$—$Y$、$C$—$Z$ 表示，它们在空间位置上彼此相差 $120^\circ$。

圆筒形铁芯的中间是可旋转的转子，转子上装有\textbf{主磁极}，磁极铁芯上套有励磁绕组。给励磁绕组通入直流电流，磁极就表现出极性（$N$、$S$）。在 $N$ 极处磁力线由转子磁极表面流出进入定子铁芯，在 $S$ 极下磁力线由定子铁芯流出进入转子磁极，形成一对极的磁场。

当转子由原动机拖动以恒定转速 $n_1$（\textbf{同步转速}）旋转时，定子绕组切割转子磁场感应交流电动势，感应电动势方向由右手定则确定。转子每旋转一周，感应电动势交变一次；转速单位为 r/min，故
\begin{equation}
f=\frac{pn_1}{60}\ \text{(Hz)} .
\end{equation}
又因三相绕组在空间上互差 $120^\circ$，转子磁场被 $A$ 相绕组切割后再转过 $120^\circ$ 被 $B$ 相切割，故 $B$ 相电动势滞后 $A$ 相 $120^\circ$，$C$ 相又滞后 $B$ 相 $120^\circ$，得到三相对称感应电动势。

同步发电机三相绕组通常接成星形，加上三相对称负载后，在三相对称电动势作用下向负载输出电能；与此同时，三相绕组中的电流与磁场相互作用，产生与原动机拖动转矩相反的\textbf{电磁转矩}，克服原动机的机械能，把机械能转换成电能，形成机电能量转换。

\subsection{感应电动机的基本工作原理}
感应电动机的定子结构与同步发电机相同，由铁芯和交流绕组构成。转子也由铁芯和绕组构成，但结构与同步发电机不同：以常用的三相笼型感应电动机为例，靠近转子铁芯外表面均匀开槽，槽里嵌放\textbf{导条}（铜条或铸铝），导条两端用\textbf{端环}连接，形成笼型绕组。

\subsubsection{三相对称电流产生旋转磁场}
约定：当 $A$ 相电流为正值时，电流从 $A$ 导体流入纸面、从 $X$ 导体流出纸面；$B$、$C$ 相类似。

某时刻 $A$ 相电流为最大值 $I_m$、$B$、$C$ 相电流均为 $-I_m/2$。此时 $A$ 相电流从 $A$ 端流入、$X$ 端流出；$B$、$C$ 相电流分别从 $Y$、$Z$ 端流入、从 $B$、$C$ 端流出。由右手定则可判断，三相绕组电流产生的合成磁动势处于\textbf{$A$ 相绕组的轴线}上。此时 $A$ 相电流最大，对合成磁动势贡献最大；$B$、$C$ 相瞬时电流只有 $I_m/2$，且其作用方向与合成磁动势方向成 $60^\circ$ 夹角，故两相对合成磁动势的贡献分别为 $\tfrac12F_{\phi1}\cos60^\circ$。于是合成磁动势
\begin{equation}
F=\frac32F_{\phi1},
\end{equation}
即合成磁动势为单相绕组产生磁动势最大值的 $3/2$ 倍。

再过 $30^\circ$（时间上），合成磁动势位置与前一时刻相比\textbf{旋转了 $30^\circ$}；此时合成磁动势由 $A$ 相与 $B$（$C$）相电流共同产生，两者作用方向与合成磁动势的夹角都是 $30^\circ$，合成磁动势仍为 $\tfrac32F_{\phi1}$。既然 $A$ 相电流最大时磁动势最大值在 $A$ 相轴线上，那么 $B$ 相电流达最大时合成磁动势必在 $B$ 相轴线上，$C$ 相电流最大时在 $C$ 相轴线上。以此类推：\textbf{电流变化一个周期（$2\pi$），磁动势在空间旋转一周，且幅值始终保持为 $\tfrac32F_{\phi1}$}。该磁动势产生两极磁场，其转速
\begin{equation}
n_1=\frac{60f}{p}\ \text{(r/min)} .
\end{equation}
这就是三相对称交流绕组通以三相对称交流电流产生旋转磁场的基本概念。

\subsubsection{转子转动的物理过程}
把定子旋转磁场用旋转磁极表示，设其逆时针同步旋转、转速 $n_1$。当转子静止（$n=0$）或转速 $n<n_1$ 时，转子导条与旋转磁场有相对运动，导条中感应电动势。因转子速度低于 $n_1$，导条切割磁力线的运动方向为顺时针，由右手定则可判断感应电动势方向；导条被端环短接，若不计导条电抗影响，导条中电流与感应电动势同相位。载流导体在磁场中受电磁力，由左手定则可判断方向，电磁力与转子半径之积即\textbf{电磁转矩}。转子上的电磁转矩为驱动力矩，若能克服负载和摩擦阻力矩，电机就以转速 $n$ 旋转。

因为转子电路中没有外接电源，完全依靠转子与旋转磁场之间的相对运动感应电动势、电流并产生电磁转矩，故称\textbf{感应电动机}；又因为带动机械负载时转速总要低于同步转速 $n_1$，故又称\textbf{异步电动机}。

反之，定子绕组 $A$—$X$、$B$—$Y$、$C$—$Z$ 三相对称（即三相绕组导体沿定子内表面按 $AZBXCY$ 规律排列），当转子以同步速旋转时，定子绕组中将产生三相对称感应电动势。

\subsection{多极电机}
实际交流电机制成多极。多极电机可视作若干个两极电机拼合而成：四极电机可看作两个两极电机拼合，各占半个圆周（机械角 $180^\circ$）。转子转过一对极（四极电机的半周），定子绕组感应电动势交变一次；转子转过一周，电动势交变两次。一般地，对 $p$ 对极电机，转子转过一周感应电动势交变 $p$ 次，故
\begin{equation}
f=\frac{pn}{60}.
\end{equation}

感应电动机中，两极电机电流交变一次、磁场旋转一周（转过一对极）；$p$ 对极电机可视作 $p$ 个两极电机拼合，电流交变一次磁场仍转一对极（$1/p$ 圆周），故 $p$ 对极电机旋转磁场转速为
\begin{equation}
n_1=\frac{60f}{p}.
\end{equation}

\keyformula{n_1=\frac{60f}{p}\ \text{(r/min)},\qquad f=\frac{pn}{60}}

两式实际上是同一公式的两种写法：对同步发电机，转子旋转一周定子电动势变化 $p$ 次；对感应电动机，定子电流交变一次磁场旋转 $1/p$ 周。即电气量随时间变化 $360^\circ$，则磁场或转子在空间旋转 $360^\circ/p$。若用电气量变化的角度（\textbf{电角度}）表示空间机械角度，则
\begin{equation}
\text{电角度}=p\times\text{空间机械角度} .
\end{equation}

引入空间电角度的好处：无论一对极还是 $p$ 对极电机，用电角度表示空间机械角度后，该角度与对应的感应电动势（或电流）随时间变化的角度一致。另外，相邻两相绕组空间位置上的机械角差为 $120^\circ/p$，空间电角差为 $120^\circ$，恰好等于相邻两绕组中感应电动势的时间相位差——\textbf{时间相位差与空间电角度差相一致}，这是交流绕组分析的总纲。

\corepoint{三句话背下来，第 8 章就通了一半：① $n_1=60f/p$ 描述"旋转磁场多快"；② $f=pn/60$ 描述"感应电动势多快"；③ 电角度 $=p\times$ 机械角度描述"角度怎么换算"。凡是遇到"绕组空间位置差多少度、电动势相位差多少度"，一律先换到电角度再算，千万不要把机械角度当电角度用。}

\section{三相单层绕组}

\subsection{从集中绕组到分布绕组}
图 8-1、8-5 中的两极、四极三相定子绕组都属于\textbf{集中绕组}，实际电机中的绕组比这复杂。集中绕组有很多缺点：电气性能差、散热不好，而且浪费铜线、铁芯等有效材料。实际电机大多采用\textbf{分布绕组}，它由集中绕组演变而来：把原来集中在一个大槽中的导体分散放置在 $q$ 个小槽中，$q$ 个小槽对应的圆周角度称为\textbf{相带}。所有大槽都这样处理，就得到 $Q=2pmq$ 个小槽构成的分布绕组；每小槽的导体都属于同一线圈，故称\textbf{单层绕组}。属于同一相的 $q$ 个槽分布在 $60^\circ$ 空间电角度范围内，故称 \textbf{$60^\circ$ 相带绕组}。

分布绕组各相带内导体的连接应与集中时相同，即每相带 $q$ 个导体依次串联，串联的先后次序可以改变。为节省铜线、制造方便和端部布置合理，按线圈形状和端部连接方法，单层绕组可分为\textbf{同心式、链式和交叉式}。

\subsection{几个基本量}
\begin{enumerate}
  \item \textbf{极距} $\tau$：对应于一个极（即 $180^\circ$ 空间电角度）的定子内表面圆弧长度，可用槽数或圆弧长度度量：
  \begin{equation}
  \tau=\frac{Q}{2p}\ \text{(槽)}\qquad\text{或}\qquad \tau=\frac{\pi D}{2p},
  \end{equation}
  其中 $Q$ 为定子槽数，$D$ 为定子铁芯内径。
  \item \textbf{节距} $y$：一个线圈的两个线圈边之间的圆弧长度，可用槽数或圆弧长度度量。$y=\tau$ 为整距，$y<\tau$ 为短距。
  \item \textbf{每极每相槽数} $q$：每个极下每相所占的槽数，即由集中绕组过渡到分布绕组时的小槽数：
  \begin{equation}
  q=\frac{Q}{2pm},
  \end{equation}
  其中 $m$ 为相数，三相电机 $m=3$。
  \item \textbf{槽距角} $\alpha$：相邻两槽之间的电角度，
  \begin{equation}
  \alpha=\frac{p\times360^\circ}{Q}.
  \end{equation}
\end{enumerate}

\subsection{同心式绕组}
同心式绕组由不等节距的同心线圈构成。例如三相二极 24 槽的定子，每极每相槽数
\begin{equation}
q=\frac{Q}{2pm}=\frac{24}{2\times1\times3}=4 .
\end{equation}
按 $60^\circ$ 相带原则划分为 6 个相带，各相带内的槽号按顺序排列（每相带 4 槽：$23,24,1,2$ / $3,4,5,6$ / …）。

由于相带 $A$ 与相带 $X$ 处于相异磁极下，其电流方向相反：设导体 $23,24,1,2$ 中电流流入纸面，则导体 $11,12,13,14$ 中电流流出纸面。把 $1$ 和 $12$ 相连构成一个大线圈，$2$ 和 $11$ 相连构成一个小线圈，这一大一小构成一个\textbf{同心式线圈组}；再把 $13$ 与 $24$ 相连、$14$ 与 $23$ 相连，构成另一个同心式线圈组；然后顺着电流流向把两个线圈组连接起来（尾端与尾端相接、首首引出），作为 $A$、$X$ 端，即构成 $A$ 相绕组。其余两相同理。

\subsection{链式与交叉式绕组}
\begin{itemize}
  \item \textbf{链式绕组}：线圈节距相同、端部如链环相扣，各线圈依次串联，端部较短、省铜线。
  \item \textbf{交叉式绕组}：由大小两种节距的线圈交叉组合，端部排列更紧凑，多用于 $q$ 为奇数的场合。
\end{itemize}

单层绕组的主要优点是：\textbf{槽内无须层间绝缘、槽利用率高、嵌线方便}，多用于散下线的 10\,kW 以下小功率交流电机；主要缺点是\textbf{电气性能不如双层绕组}——不能通过选择节距来削弱谐波电动势与磁动势。

\section{三相双层绕组}

双层绕组的每个槽分上、下两层，每槽内放置两个线圈边（一个上层边、一个下层边）；线圈数等于槽数，一个线圈的一个边放在某槽上层，另一个边放在相隔节距 $y$ 的另一槽下层。

\subsection{结构特点}
\begin{enumerate}
  \item \textbf{可选择节距}：双层绕组可以做成短距（通常 $y_1\approx\tfrac56\tau$ 或 $\tfrac{5}{6}\tau$ 附近），从而\textbf{同时削弱 5 次和 7 次谐波}电动势与磁动势，改善波形（第 9 章第五节将定量说明）。
  \item \textbf{端部排列整齐}：所有线圈形状相同，端部易于布置、便于制造。
  \item \textbf{分叠绕组与波绕组}：叠绕组是相邻线圈依次串联、后一线圈紧叠在前一线圈上；波绕组是线圈沿圆周"波浪式"前进，多用于极数较多、支路数要求特殊的场合。
\end{enumerate}

在绕组展开图中，上层边用实线表示、下层边用虚线表示，每一个线圈都是一半实线、一半虚线。按 $60^\circ$ 相带划分后，各相带槽号的分配与单层绕组相同，只是每个槽号对应上、下两个线圈边。

\subsection{两种绕组对比}
\begin{center}
\begin{tabular}{@{}p{2.2cm}p{3.6cm}p{3.9cm}p{3.0cm}@{}}
\toprule
项目 & 单层绕组 & 双层绕组 & 说明 \\
\midrule
每槽线圈边 & 1 个 & 2 个（上、下层） & 双层槽利用率略低 \\
线圈数 & $Q/2$ & $Q$ & 双层线圈数等于槽数 \\
节距 & 一般为整距 & 可选短距（$\approx\tfrac56\tau$） & 短距可削弱 5、7 次谐波 \\
端部 & 同心式/链式/交叉式 & 叠绕组/波绕组 & 双层端部整齐 \\
电气性能 & 较差 & 好 & 大、中型电机多用双层 \\
典型用途 & 10\,kW 以下小功率电机 & 大、中型交流电机 & — \\
\bottomrule
\end{tabular}
\end{center}

\corepoint{绕组这一章的所有计算都归结为四个量：$p$（极对数）、$Q$（槽数）、$m$（相数）、$y$（节距）。由它们派生出 $\tau=Q/(2p)$、$q=Q/(2pm)$、$\alpha=360^\circ p/Q$。做题第一步永远是先算这四个派生量并写清楚，再画槽号分配表；槽号分配错了，后面电动势、磁动势全错。另外牢记：单层绕组\textbf{不能}用短距削弱谐波，双层才可以——这是选择绕组型式的第一条理由。}

% ============================================================
% 第 9 章 交流绕组的感应电动势（教材第三篇，印刷页 90--103）
% 节名：第一节 导体电动势星形图；第二节 导体基波电动势；第三节 线圈基波电动势及节距因数；
%       第四节 线圈组基波电动势及分布因数；第五节 谐波电动势的计算；第六节 削弱谐波电动势的方法
% ============================================================
\chapter{交流绕组的感应电动势}

第 8 章定性地说明了交流绕组中产生感应电动势的原理。本章把它定量化：先求一根导体，再求一个线圈，再求一个线圈组（极相组），最后求一相绕组。这条"导体 $\rightarrow$ 线圈 $\rightarrow$ 线圈组 $\rightarrow$ 相绕组"的链条，每一步只多乘一个因数（节距因数、分布因数、绕组因数），把整章串起来记最省力。

\section{导体电动势星形图}

交流绕组内的感应电动势通常为正弦交流电动势，可用相量表示和运算。对于多极电机，由第 8 章可知：转子旋转一周，定子绕组感应电动势变化 $p$ 次。

把定子各槽中导体的基波电动势画成相量，由于相邻两槽在空间相差电角度 $\alpha=p\cdot360^\circ/Q$，其感应电动势在时间上也相差 $\alpha$，于是各导体电动势相量构成一个\textbf{辐射形相量图}——这就是\textbf{导体电动势星形图}。$p$ 对极电机中，导体电动势星形图有 $p$ 个重合的"层"，每层包含 $Q/p$ 个相量，相邻相量夹角为 $\alpha$。

导体电动势星形图的作用：
\begin{itemize}
  \item 决定绕组的连接方式（哪些导体属于同一相带、如何串联）；
  \item 计算每相绕组感应电动势的大小（把属于同一相的导体电动势按相量相加）。
\end{itemize}
对于复杂的绕组（尤其是分数槽绕组），更需要借助电动势星形图才能解决绕组连接和感应电动势的计算问题。

\section{导体基波电动势}

由电磁感应定律，导体在正弦分布的气隙磁场中切割磁力线，感应电动势的有效值为
\begin{equation}
E_{c1}=2.22f\Phi_1 ,
\end{equation}
其中 $f$ 为电动势频率（$f=pn/60$），$\Phi_1$ 为每极基波磁通。

\keyformula{E_{c1}=2.22f\Phi_1}

\subsubsection{系数 2.22 的来源}
设气隙基波磁密幅值为 $B_1$，导体有效长度为 $l$，极距为 $\tau$，则每极基波磁通 $\Phi_1=\dfrac{2}{\pi}B_1\tau l$。导体以同步速度切割该磁场，感应电动势幅值 $E_{cm}=B_1lv$，而 $v=2\tau f$，故
\begin{equation}
E_{cm}=B_1l\cdot2\tau f=\frac{\pi}{\sqrt2}\cdot\frac{\sqrt2}{\pi}B_1l\tau\omega
=\pi f\left(\frac{2}{\pi}B_1\tau l\right)\sqrt2 ,
\end{equation}
取有效值即得 $E_{c1}=\dfrac{\pi}{\sqrt2}f\Phi_1=2.22f\Phi_1$。\textbf{2.22 是"一根导体、一个极距、正弦磁场"的标志}；后面线圈、线圈组、相绕组乘上各自的因数，系数就变成 4.44。

\section{线圈基波电动势及节距因数}

单层或双层绕组均由线圈构成。线圈的节距 $y_1$ 通常为整距（$y_1=\tau$）或短距（$y_1<\tau$）。

\subsection{整距线圈}
整距线圈的两个边相距 $\tau$（即 $180^\circ$ 电角度），两边中的感应电动势大小相等、方向相反（沿回路相加），故线圈电动势为两个导体电动势的\textbf{代数和}：
\begin{equation}
E_{y1}=2E_{c1}=4.44fN_c\Phi_1 ,
\end{equation}
其中 $N_c$ 为线圈匝数。

\subsection{短距线圈与节距因数}
当线圈短距时，两个线圈边相距 $y_1<\tau$，对应的空间电角度为 $\beta=\dfrac{y_1}{\tau}\times180^\circ$。两个边的电动势相量不共线，而是相差 $\beta$，其相量和为
\begin{equation}
E_{y1}=2E_{c1}N_c\cos\frac{180^\circ-\beta}{2}=2E_{c1}N_c\sin\frac{\beta}{2}
=4.44fN_c\Phi_1k_{p1},
\end{equation}
其中
\begin{equation}
k_{p1}=\sin\frac{\beta}{2}=\sin\left(\frac{y_1}{\tau}\times90^\circ\right)
\end{equation}
称为基波\textbf{节距因数}（短距因数）。

\keyformula{k_{p1}=\sin\left(\frac{y_1}{\tau}\times90^\circ\right),\qquad E_{y1}=4.44fN_c\Phi_1k_{p1}}

$k_{p1}<1$，表示短距使基波电动势"打了折扣"。整距时 $y_1=\tau$，$k_{p1}=1$；短距时 $k_{p1}<1$。例如 $y_1=\tfrac56\tau$ 时 $k_{p1}=\sin75^\circ=0.966$，基波只损失 3.4\%，却能把 5、7 次谐波大幅削弱——这是双层绕组普遍采用 $\tfrac56$ 短距的原因。

\section{线圈组基波电动势及分布因数}

\subsection{分布绕组}
实际电机的每相绕组由 $q$ 个线圈（分布在 $q$ 个槽中）串联组成，这 $q$ 个线圈构成一个\textbf{线圈组}（极相组）。由于 $q$ 个槽在空间依次相差槽距角 $\alpha$，各线圈电动势相量也依次相差 $\alpha$，其相量和小于代数和。

\subsection{分布因数}
设 $q$ 个线圈电动势大小相等、依次相差 $\alpha$，其相量和为 $qE_{y1}k_{d1}$，其中
\begin{equation}
k_{d1}=\frac{\sin\left(q\dfrac{\alpha}{2}\right)}{q\sin\dfrac{\alpha}{2}}
\end{equation}
称为基波\textbf{分布因数}。

\keyformula{k_{d1}=\frac{\sin\left(q\alpha/2\right)}{q\sin\left(\alpha/2\right)}}

$k_{d1}<1$，表示分布使基波电动势减小。$q=1$（集中绕组）时 $k_{d1}=1$。

\subsection{线圈组电动势与绕组因数}
一个线圈组（极相组）的基波电动势为
\begin{equation}
E_{q1}=qE_{y1}k_{d1}=4.44fqN_c\Phi_1k_{p1}k_{d1}=4.44fqN_c\Phi_1k_{w1},
\end{equation}
其中
\begin{equation}
k_{w1}=k_{p1}k_{d1}
\end{equation}
称为基波\textbf{绕组因数}。

\subsection{一相绕组的电动势}
设每相绕组串联总匝数为 $N$（$N=$ 每相串联线圈数 $\times$ 每个线圈匝数 $N_c$），则一相绕组的基波电动势为
\begin{equation}
E_{\phi1}=4.44fN\Phi_1k_{w1}.
\end{equation}

\keyformula{E_{\phi1}=4.44fN\Phi_1k_{w1},\qquad k_{w1}=k_{p1}k_{d1}}

若绕组为双层、每相有 $a$ 条并联支路，则每相串联匝数 $N=\dfrac{2pqN_c\cdot(\text{每极相组数}/\text{支路})}{a}$，计算时注意\textbf{串联}二字——并联支路中的匝数不参与串联相加。

\begin{center}
\begin{tabular}{@{}llll@{}}
\toprule
层次 & 电动势表达式 & 新增因数 & 因数表达式 \\
\midrule
一根导体 & $E_{c1}=2.22f\Phi_1$ & — & — \\
一个线圈 & $E_{y1}=4.44fN_c\Phi_1k_{p1}$ & 节距因数 & $k_{p1}=\sin\left(\dfrac{y_1}{\tau}90^\circ\right)$ \\
一个线圈组 & $E_{q1}=4.44fqN_c\Phi_1k_{w1}$ & 分布因数 & $k_{d1}=\dfrac{\sin(q\alpha/2)}{q\sin(\alpha/2)}$ \\
一相绕组 & $E_{\phi1}=4.44fN\Phi_1k_{w1}$ & 绕组因数 & $k_{w1}=k_{p1}k_{d1}$ \\
\bottomrule
\end{tabular}
\end{center}

\corepoint{这一节的"递进结构"必须记牢：\textbf{2.22（导体）$\rightarrow$ 4.44（线圈，整距）$\rightarrow$ 乘 $k_{p1}$（短距）$\rightarrow$ 乘 $k_{d1}$（分布）}。做题时先判断问的是哪一层，再决定乘哪些因数；$k_{w1}$ 永远等于 $k_{p1}k_{d1}$，若题目给了 $k_{w1}$ 就不必分开算。另外注意 $N$ 是"每相串联匝数"，不是"每相总匝数"，并联支路要除以 $a$。}

\section{谐波电动势的计算}

前四节假定主磁极产生的气隙磁场沿定子内表面\textbf{正弦分布}，此时绕组中感应电动势是频率为 $f$ 的正弦波。实际气隙磁场并非正弦分布（因磁极形状、齿槽、饱和等），可分解为基波与一系列奇次谐波（3、5、7、…）。$v$ 次谐波磁场的极对数为 $vp$，故其感应电动势频率为
\begin{equation}
f_v=vf .
\end{equation}
仿照基波的推导，$v$ 次谐波电动势的有效值为
\begin{equation}
E_{\phi v}=4.44f_vN\Phi_vk_{wv}=4.44(vf)N\Phi_vk_{wv},
\end{equation}
其中
\begin{equation}
k_{pv}=\sin\left(\frac{vy_1}{\tau}\times90^\circ\right),\qquad
k_{dv}=\frac{\sin\left(vq\dfrac{\alpha}{2}\right)}{q\sin\left(v\dfrac{\alpha}{2}\right)},\qquad
k_{wv}=k_{pv}k_{dv}.
\end{equation}
相电动势的有效值为基波与各次谐波电动势的\textbf{平方和开方}：
\begin{equation}
E_\phi=\sqrt{E_{\phi1}^{2}+E_{\phi3}^{2}+E_{\phi5}^{2}+E_{\phi7}^{2}+\cdots}
\end{equation}

\subsection{三相绕组中的 3 次及 3 的倍数次谐波}
在三相绕组中，3 次谐波电动势在三相之间\textbf{同相位}（相位差 $3\times120^\circ=360^\circ$）。因此：
\begin{itemize}
  \item 若三相接成\textbf{星形}，3 次谐波电动势在相电动势中存在，但\textbf{线电动势中互相抵消为零}；
  \item 若三相接成\textbf{三角形}，3 次谐波电动势将在闭合三角形回路内产生 3 次谐波\textbf{环流}，引起附加损耗和发热。
\end{itemize}
因此同步发电机大多采用\textbf{星形接法}——既消除线电动势中的 3 次谐波，又避免三角形接法中的环流。

\section{削弱谐波电动势的方法}

实际交流电机绕组中总会存在谐波电动势。谐波电动势使电动势波形畸变，谐波电流引起附加损耗、发热、振动和通信干扰，因此必须削弱。常用方法有四条。

\subsection{采用短距绕组}
短距可以削弱或消除某一指定次数的谐波。要使 $v$ 次谐波电动势为零，须 $k_{pv}=0$，即
\begin{equation}
\frac{vy_1}{\tau}\times90^\circ=180^\circ
\quad\Longrightarrow\quad
y_1=\tau\left(1-\frac1v\right).
\end{equation}
例如要消除 5 次谐波，取 $y_1=\tfrac45\tau=0.8\tau$；要消除 7 次谐波，取 $y_1=\tfrac67\tau\approx0.857\tau$。工程上常取 $y_1\approx\tfrac56\tau$（$0.833\tau$）——此时 5 次和 7 次谐波都被削弱到很小（$k_{p5}\approx0.259$、$k_{p7}\approx-0.259$），而基波只损失约 3.4\%，是"一笔划算的买卖"。

\subsection{采用分布绕组}
采用分布绕组也可以削弱谐波电动势，削弱的程度取决于谐波电动势的\textbf{分布因数} $k_{dv}$。$q$ 越大，$k_{dv}$ 越小（对 $v>1$ 的谐波尤其显著），而基波分布因数 $k_{d1}$ 下降不多，因此分布绕组对谐波"削得很狠、对基波很温柔"。当然 $q$ 增大会使槽数增多、槽利用率下降，工程上一般取 $q=2\sim6$。

\subsection{采用斜槽}
把定子（或转子）槽沿轴向扭斜一个角度，使同一根导体沿轴向各点切割的磁场相位不同，从而削弱齿谐波电动势。斜槽因数
\begin{equation}
k_{skv}=\frac{\sin\left(v\dfrac{\beta}{2}\right)}{v\dfrac{\beta}{2}},
\end{equation}
其中 $\beta$ 为斜槽对应的空间电角度。斜槽主要用于削弱\textbf{齿谐波}（$v=Q/p\pm1$ 次），因为齿谐波的绕组因数与基波相同，靠短距和分布无法削弱，只能靠斜槽或增大气隙。

\subsection{采用星形接法}
如第 9.5 节所述，星形接法使线电动势中不含 3 次及 3 的倍数次谐波，代价为零，因此几乎所有同步发电机都采用星形接法。

\begin{center}
\begin{tabular}{@{}p{2.4cm}p{4.0cm}p{5.6cm}@{}}
\toprule
方法 & 削弱对象 & 代价/限制 \\
\midrule
短距绕组 & 指定次数谐波（如 5、7 次） & 基波电动势略降（$\tfrac56\tau$ 时降 3.4\%） \\
分布绕组 & 各次谐波（$q$ 越大越显著） & 槽数增多、槽利用率下降，$q=2\sim6$ \\
斜槽 & 齿谐波（$Q/p\pm1$ 次） & 制造工艺复杂，基波略有削弱 \\
星形接法 & 线电动势中的 3 次及倍数次谐波 & 无代价；避免三角形环流 \\
\bottomrule
\end{tabular}
\end{center}

\corepoint{四条措施的"分工"要分清：\textbf{短距}打指定次谐波（$v$ 次要 $y_1=\tau(1-1/v)$）；\textbf{分布}打所有高次谐波（$q$ 越大越狠）；\textbf{斜槽}专打齿谐波（因为齿谐波的 $k_w$ 与基波相同，短距、分布都无效）；\textbf{星形接法}清掉线电动势里的 3 次谐波。考试若问"为什么取 $y_1\approx5/6\tau$"，答"同时削弱 5、7 次且基波仅损失 3.4\%，综合最经济"。}

% ============================================================
% 第 10 章 交流绕组的磁动势（教材第三篇，印刷页 104--115）
% 节名：第一节 单相绕组的脉振磁动势；第二节 单相脉振磁动势的分解；第三节 三相绕组的合成磁动势——旋转磁动势
% ============================================================
\chapter{交流绕组的磁动势}

第 9 章解决了"磁通变 $\rightarrow$ 电动势"，本章解决反向的问题："电流变 $\rightarrow$ 磁动势"。这一章是全书最抽象的一章，也是后面感应电机、同步电机全部理论的落脚点：感应电机的转矩来自定、转子磁动势的合成，同步电机的电枢反应就是电枢磁动势对励磁磁动势的作用。看这一章要抓住一条主线：\textbf{一个脉振磁动势 = 两个反向旋转磁动势；三个脉振磁动势合成 = 一个圆形旋转磁动势}。

\section{单相绕组的脉振磁动势}

\subsection{整距集中线圈的磁动势}
先看最简单的两极电机、一个整距集中线圈。线圈中通入电流 $i$ 时，线圈产生的磁动势沿气隙圆周的分布是一个\textbf{矩形波}：在包含线圈边的范围内磁动势为 $N_ci/2$（每极安匝），在范围外为 $-N_ci/2$，跳变发生在线圈边处。

对 $p$ 对极电机，每极磁动势幅值为
\begin{equation}
F_c=\frac{N_ci}{2p}.
\end{equation}
矩形波可分解为基波与一系列奇次谐波（1、3、5、…）。用傅里叶分解，基波幅值为矩形波幅值的 $4/\pi$ 倍：
\begin{equation}
F_{c1}=\frac{4}{\pi}\cdot\frac{N_ci}{2p}=0.9\frac{N_ci}{p}.
\end{equation}
$v$ 次谐波幅值为基波的 $1/v$：
\begin{equation}
F_{cv}=\frac{1}{v}F_{c1}=\frac{0.9}{v}\cdot\frac{N_ci}{p}.
\end{equation}
可见谐波磁动势的幅值随次数迅速减小，但次数越高、极对数越多、转速越低，其影响不可一概忽略。

\subsection{短距与分布的影响}
实际绕组不是整距集中线圈，而是短距、分布的线圈组。与电动势的处理完全类似：
\begin{itemize}
  \item 线圈\textbf{短距}时，两个线圈边的磁动势波在空间错开 $\beta$ 角，合成后基波幅值乘以\textbf{节距因数} $k_{p1}$；
  \item $q$ 个线圈\textbf{分布}放置时，各线圈的脉振磁动势在空间依次错开槽距角 $\alpha$，合成后基波幅值乘以\textbf{分布因数} $k_{d1}$；
  \item 两者的乘积 $k_{w1}=k_{p1}k_{d1}$ 即\textbf{绕组因数}。
\end{itemize}

\subsection{单相绕组的基波脉振磁动势}
设一相绕组每相串联匝数为 $N$，通入正弦电流 $i=\sqrt2I\sin\omega t$，则单相绕组产生的基波磁动势幅值为
\begin{equation}
F_{\phi1}=0.9\frac{Nk_{w1}}{p}I ,
\end{equation}
其沿气隙圆周的分布与时间变化可写成
\begin{equation}
f_{\phi1}(x,t)=F_{\phi1}\cos\frac{\pi}{\tau}x\cdot\sin\omega t ,
\end{equation}
其中 $x$ 为沿气隙圆周的坐标，$\tau$ 为极距。

\keyformula{F_{\phi1}=0.9\frac{Nk_{w1}}{p}I}

\subsubsection{"脉振"的含义}
上式表明：在空间任一固定点（$x$ 固定），磁动势随时间按 $\sin\omega t$ 交变；在任一固定时刻（$t$ 固定），磁动势沿空间按余弦分布。\textbf{空间位置固定、幅值随时间交变}的磁动势称为\textbf{脉振磁动势}。它的轴线固定在该相绕组的轴线上，忽正忽负地"原地脉动"，不旋转。

\begin{center}
\begin{tabular}{@{}lll@{}}
\toprule
层次 & 基波幅值 & 说明 \\
\midrule
整距集中线圈（两极） & $F_{c1}=0.9N_ci$ & 矩形波的 $4/\pi$ 倍 \\
整距集中线圈（$p$ 对极） & $F_{c1}=0.9N_ci/p$ & 每极安匝 $N_ci/(2p)$ \\
短距线圈 & 乘 $k_{p1}$ & 两边磁动势波错开 \\
分布线圈组 & 再乘 $k_{d1}$ & $q$ 个波错开 $\alpha$ \\
一相绕组 & $F_{\phi1}=0.9Nk_{w1}I/p$ & $N$ 为每相串联匝数 \\
\bottomrule
\end{tabular}
\end{center}

\corepoint{磁动势这一章的所有幅值公式都长成"$0.9\times$ 每极安匝 $\times$ 因数"的样子：$0.9$ 来自矩形波分解的 $4/\pi$ 与 $\sqrt2$ 的合成，$1/p$ 来自"每极分摊"，因数来自短距与分布。记住这个结构，就不会把 $0.9$ 与电动势里的 $2.22$、$4.44$ 搞混——\textbf{电动势公式里是 2.22/4.44，磁动势公式里是 0.9}，这是两个不同的傅里叶系数。}

\section{单相脉振磁动势的分解}

单相绕组通以交流电流时，产生基波和一系列谐波脉振磁动势。基波磁动势产生的磁场幅值随时间脉振。利用三角恒等式
\begin{equation}
\cos\frac{\pi}{\tau}x\cdot\sin\omega t
=\frac12\sin\left(\omega t-\frac{\pi}{\tau}x\right)+\frac12\sin\left(\omega t+\frac{\pi}{\tau}x\right),
\end{equation}
即
\begin{equation}
f_{\phi1}(x,t)=\frac{F_{\phi1}}{2}\sin\left(\omega t-\frac{\pi}{\tau}x\right)
+\frac{F_{\phi1}}{2}\sin\left(\omega t+\frac{\pi}{\tau}x\right).
\end{equation}

\keyformula{f_{\phi1}=\underbrace{\frac{F_{\phi1}}{2}\sin\left(\omega t-\frac{\pi}{\tau}x\right)}_{\text{正向旋转波}}
+\underbrace{\frac{F_{\phi1}}{2}\sin\left(\omega t+\frac{\pi}{\tau}x\right)}_{\text{反向旋转波}}}

两个分量都是行波：
\begin{itemize}
  \item 第一项 $\sin\left(\omega t-\pi x/\tau\right)$ 代表一个\textbf{正向旋转磁动势}，其幅值为 $F_{\phi1}/2$，转速为同步转速 $n_1=60f/p$，转向为 $+x$ 方向；
  \item 第二项 $\sin\left(\omega t+\pi x/\tau\right)$ 代表一个\textbf{反向旋转磁动势}，幅值同为 $F_{\phi1}/2$，转速仍为 $n_1$，但转向相反。
\end{itemize}

\subsection{结论}
\textbf{一个脉振磁动势可以分解为两个幅值相等（各为脉振幅值的一半）、转速相同（均为同步转速）、转向相反的两个旋转磁动势。}反过来，两个幅值相等、转向相反的旋转磁动势合成，就是一个脉振磁动势。

这个结论是理解单相感应电动机（第 15 章）的钥匙：单相绕组产生的磁场可看作正、反两个旋转磁场叠加，转子对正转磁场转差率为 $s$、对反转磁场转差率为 $2-s$，两个转矩方向相反，故单相感应电动机\textbf{没有起动转矩}。

\section{三相绕组的合成磁动势——旋转磁动势}

第 8 章已定性地说明：三相对称绕组通以三相对称电流，产生旋转磁场的磁动势。本节给出定量的结果。

\subsection{合成磁动势的表达式}
设三相绕组在空间互差 $120^\circ$ 电角度，三相电流对称且互差 $120^\circ$ 时间电角度。取 $A$ 相绕组轴线为坐标原点，则三相脉振磁动势分别为
\begin{equation}
\begin{aligned}
f_A&=F_{\phi1}\cos\frac{\pi}{\tau}x\sin\omega t,\\
f_B&=F_{\phi1}\cos\left(\frac{\pi}{\tau}x-120^\circ\right)\sin\left(\omega t-120^\circ\right),\\
f_C&=F_{\phi1}\cos\left(\frac{\pi}{\tau}x-240^\circ\right)\sin\left(\omega t-240^\circ\right).
\end{aligned}
\end{equation}
把每一个脉振磁动势分解为两个旋转磁动势，再三相相加。结果：\textbf{反转的两个分量互相抵消，正转的三个分量同相位相加}，合成的基波磁动势为
\begin{equation}
f_1(x,t)=\frac32F_{\phi1}\sin\left(\omega t-\frac{\pi}{\tau}x\right).
\end{equation}

\subsection{旋转磁动势的三个特征}
\begin{enumerate}
  \item \textbf{幅值恒定}：
  \begin{equation}
  F_1=\frac32F_{\phi1}=1.35\frac{Nk_{w1}}{p}I .
  \end{equation}
  幅值不随时间变化，故称\textbf{圆形旋转磁动势}（若三相电流不对称，则幅值随时间变化，轨迹为椭圆，称椭圆形旋转磁动势）。
  \item \textbf{转速为同步转速}：
  \begin{equation}
  n_1=\frac{60f}{p}.
  \end{equation}
  \item \textbf{转向由电流相序决定}：合成磁动势的轴线总是\textbf{从电流超前相绕组轴线转向电流滞后相绕组轴线}。当某相电流达到最大值时，合成磁动势的幅值恰好位于该相绕组的轴线上（第 8 章已用图解法验证过这一点）。
\end{enumerate}

\keyformula{F_1=\frac32F_{\phi1}=1.35\frac{Nk_{w1}}{p}I,\qquad n_1=\frac{60f}{p}}

\subsection{改变转向的方法}
由特征 3 立即可得：\textbf{把三相电源线中的任意两根对调}（即改变相序），旋转磁场即反向旋转。这就是三相感应电动机"改变相序实现反转"的原理。

\subsection{谐波磁动势的合成}
三相绕组产生的谐波磁动势中：
\begin{itemize}
  \item \textbf{3 次及 3 的倍数次谐波}：三相在空间互差 $3\times120^\circ=360^\circ$（同相位），在时间上也同相位，三相合成结果\textbf{为零}——即三相对称绕组中不存在 3 次谐波旋转磁动势（这与第 9 章"线电动势中无 3 次谐波"是同一件事的两面）。
  \item \textbf{5 次谐波}：三相合成后为一个\textbf{反向}旋转磁动势，幅值为基波的 $1/5$，转速为 $n_1/5$（极对数为 $5p$）。
  \item \textbf{7 次谐波}：三相合成后为\textbf{正向}旋转磁动势，幅值为基波的 $1/7$，转速为 $n_1/7$。
  \item 一般地，$v=6k\pm1$ 次谐波存在，$v=6k\pm3$ 次谐波为零；$v=6k+1$ 次正转、$v=6k-1$ 次反转（$k=1,2,\dots$）。
\end{itemize}
谐波磁动势会产生附加转矩、附加损耗和振动噪声，故用短距、分布、斜槽加以削弱（同第 9 章第六节）。

\begin{center}
\begin{tabular}{@{}p{1.6cm}p{2.6cm}p{2.6cm}p{3.4cm}@{}}
\toprule
磁动势类型 & 幅值 & 转速/转向 & 特征 \\
\midrule
单相基波脉振 & $F_{\phi1}=0.9Nk_{w1}I/p$ & 原地脉振（不转） & 轴线固定于该相绕组轴线 \\
两个旋转分量 & 各 $F_{\phi1}/2$ & $\pm n_1$，转向相反 & 合成即脉振波 \\
三相基波合成 & $F_1=\tfrac32F_{\phi1}=1.35Nk_{w1}I/p$ & $n_1$，由相序决定 & 幅值恒定（圆形旋转） \\
三相 3 次谐波 & 0 & — & 三相合成互相抵消 \\
三相 5 次谐波 & $F_1/5$ & $n_1/5$，反向 & 引起附加转矩与损耗 \\
三相 7 次谐波 & $F_1/7$ & $n_1/7$，正向 & 同上 \\
\bottomrule
\end{tabular}
\end{center}

\corepoint{本章三条"必须能默写"的结论：① $F_{\phi1}=0.9Nk_{w1}I/p$（单相基波脉振幅值）；② 脉振 $=$ 两个幅值各半、转向相反的旋转波（这条解释单相电机为何无起动转矩）；③ $F_1=\tfrac32F_{\phi1}=1.35Nk_{w1}I/p$，转速 $n_1=60f/p$，转向由相序决定（这条是感应电机、同步电机全部理论的起点）。注意 $\tfrac32$ 这个系数：三相合成磁动势是单相脉振幅值的 1.5 倍，不是 3 倍，也不是 $\sqrt3$ 倍。}

% ============================================================
% 第 11 章 电机的发热与温升限值（教材第三篇，印刷页 116--121）
% 节名：第一节 电机的温升和温升限值；第二节 电机的发热与冷却过程；第三节 电机的工作制；第四节 电机的冷却方式
% ============================================================
\chapter{电机的发热与温升限值}

前面各章只关心电磁关系，本章转向\textbf{热}的问题。电机的容量并不完全由电磁参数决定，往往受限于绝缘材料能承受的温度：损耗产生热量，热量使绕组温度升高，温度超过绝缘的耐热极限，绝缘就会加速老化甚至击穿。因此"温升限值"是电机设计和使用中一条不可逾越的红线。

\section{电机的温升和温升限值}

\subsection{温升的定义}
电机运行中产生的各种损耗（铜耗、铁耗、机械损耗）最终都变成热能，一部分使电机温度升高，一部分散入周围介质。电机某部分的温度 $\theta$ 与周围冷却介质温度 $\theta_0$ 之差，称为该部分的\textbf{温升}：
\begin{equation}
\tau=\theta-\theta_0 .
\end{equation}
温升的单位是 K（开尔文）或 ℃（作为温差时两者数值相同）。工程上关心温升而不是绝对温度，因为散热能力取决于温差，而绝缘老化速度也主要取决于温度高出环境多少。

\subsection{绝缘材料的耐热等级}
电机绕组绝缘按耐热能力分为若干等级，各级有其"最高允许温度"。常用等级如下（按 GB 规定的耐热分级）。

\begin{center}
\begin{tabular}{@{}llll@{}}
\toprule
耐热等级 & 最高允许温度 & 主要材料 & 典型应用 \\
\midrule
A & 105\,℃ & 浸渍过的棉纱、纸、丝绸等 & 老式小电机 \\
E & 120\,℃ & 环氧树脂、聚酯薄膜等 & 中小型电机 \\
B & 130\,℃ & 云母、玻璃丝、石棉与漆结合 & 常用等级 \\
F & 155\,℃ & 云母、玻璃丝与耐热漆 & 现代电机主流 \\
H & 180\,℃ & 硅有机树脂、云母、玻璃丝 & 高温、起重、冶金 \\
C & $>180$\,℃ & 云母、玻璃、陶瓷等无机材料 & 特殊场合 \\
\bottomrule
\end{tabular}
\end{center}

\subsection{温升限值}
\textbf{温升限值}是标准规定的、电机各部分温升不得超过的最大值。它由绝缘的极限温度扣除规定的冷却介质最高温度（我国规定环境空气温度最高为 40\,℃）以及绕组内部的\textbf{热点温差}确定。

例如 B 级绝缘、最高允许温度 130\,℃，若环境温度按 40\,℃ 计，则平均温升限值约为 80\,K，再考虑绕组最热点比平均温度高若干度（热点温差），标准给出的温升限值略低于这个差。这就是为什么铭牌上写"温升 80\,K"而不是"温度 130\,℃"。

\begin{center}
\begin{tabular}{@{}lll@{}}
\toprule
项目 & 含义 & 数值举例（B 级） \\
\midrule
最高允许温度 & 绝缘长期工作不加速老化的上限 & 130\,℃ \\
环境温度上限 & 标准规定的冷却介质最高温度 & 40\,℃ \\
温升限值 & 标准规定的允许温升 & 约 80\,K（扣除热点温差后） \\
\bottomrule
\end{tabular}
\end{center}

\corepoint{三条要分清：\textbf{温度}是绝对量，\textbf{温升}是温差，\textbf{温升限值}是标准规定值。考试常问"环境温度 40\,℃ 时 B 级绝缘电机绕组允许的最高温度是多少"，答案不是简单的 130\,℃ 也不是 120\,℃，而要按"温升限值 + 环境温度"算；反过来问"温升限值"，要用"极限温度 $-$ 环境温度 $-$ 热点温差"。}

\section{电机的发热与冷却过程}

\subsection{发热过程}
设电机某部分在单位时间内产生的热量为 $Q$，其热容量为 $C$（使该部分温度升高 1\,K 所需热量），散热系数为 $A$（单位温升、单位时间散出的热量）。由能量守恒，发热与散热之差等于热容量的增温：
\begin{equation}
Q\,\mathrm{d}t=C\,\mathrm{d}\tau+A\tau\,\mathrm{d}t ,
\end{equation}
解得温升随时间的变化规律为
\begin{equation}
\tau=\tau_\infty\left(1-\mathrm{e}^{-t/T}\right),\qquad
\tau_\infty=\frac{Q}{A},\qquad
T=\frac{C}{A}.
\end{equation}

\keyformula{\tau=\tau_\infty\left(1-\mathrm{e}^{-t/T}\right),\qquad T=\frac{C}{A}}

其中 $\tau_\infty$ 为\textbf{稳态温升}（发热与散热达到平衡时的温升），$T$ 为\textbf{发热时间常数}。温升按指数曲线上升：起始上升快，逐渐变慢，理论上 $t\to\infty$ 才达到 $\tau_\infty$，工程上认为 $t=(3\sim4)T$ 时已基本稳定（达 95\%～98\%）。

\subsection{冷却过程}
电机从运行状态切除电源（或负载减小）后，$Q=0$，温升按指数曲线下降：
\begin{equation}
\tau=\tau_0\mathrm{e}^{-t/T},
\end{equation}
其中 $\tau_0$ 为冷却开始时的温升。冷却时散热条件往往变差（如风扇停转），故冷却时间常数一般大于发热时间常数。

\subsection{时间常数的工程意义}
\begin{itemize}
  \item 小型电机发热时间常数约为十几分钟到一小时；
  \item 中大型电机可达数小时。
\end{itemize}
时间常数越长，说明电机"热惯性"越大：短时过载不会立即造成危险，但过载持续时间一长就危险。这正是"短时工作制"与"连续工作制"容量可以不同的原因。

\section{电机的工作制}

电机的工作制（工作方式）指电机承受负载的情况，它决定了电机的发热与冷却循环，也决定了额定功率的标定条件。按 GB 规定，主要有三类。

\subsection{连续工作制（S1）}
电机在恒定负载下连续运行，温升达到稳态值。绝大多数工业用电动机（水泵、风机、压缩机）属此类。铭牌上的额定功率就是在连续工作制、规定的环境温度和冷却条件下，能长期输出而不超过温升限值的功率。

\subsection{短时工作制（S2）}
电机在恒定负载下运行一段有限时间（尚未达到稳态温升），随后断电停转一段时间，温升降到接近环境温度。如闸门启闭机、某些机床辅助传动。

\subsection{断续周期工作制（S3）}
电机按"运行—停歇—运行"周期性循环，每次运行时间不足以达到稳态温升，停歇时间也不足以冷却到环境温度。用\textbf{负载持续率}（暂载率）$FC$ 表征：
\begin{equation}
FC=\frac{t_r}{t_r+t_0}\times100\%,
\end{equation}
其中 $t_r$ 为运行时间、$t_0$ 为停歇时间，一个周期一般不超过 10\,min。标准负载持续率有 15\%、25\%、40\%、60\% 等。起重机、电梯、轧机等属此类。

\corepoint{同一台电机，工作制不同，允许输出的功率不同。短时或断续工作的电机，因温升来不及达到稳态，可在短时间内输出比连续额定功率更大的功率；反之，把连续工作制的电机用于短时工作制是"浪费"，把短时工作制的电机用于连续运行则会过热烧毁。看到铭牌功率，先问"是哪种工作制下的额定值"。}

\section{电机的冷却方式}

\subsection{按冷却介质分}
\begin{itemize}
  \item \textbf{空气冷却}：最常见，分自扇冷（转子自带风扇）与他扇冷（独立风机）。
  \item \textbf{氢气冷却}：用于大型汽轮发电机。氢气的导热系数与比热容比空气大，且密度小，可显著提高散热能力并降低通风损耗；但必须严格密封、防止爆炸。
  \item \textbf{水冷却}：水的比热容和导热系数远大于空气，用于大型发电机定子绕组（水内冷）及大功率变频电机。
\end{itemize}

\subsection{按防护与通风结构分}
\begin{itemize}
  \item \textbf{开启式}：与外界空气直接相通，散热好但易进尘、受潮。
  \item \textbf{防护式}：有网罩或挡板，防止固体异物进入。
  \item \textbf{封闭式}：机壳封闭，靠机壳表面散热或内部循环冷却，防尘防潮好，但散热能力较差，同容量体积较大。
\end{itemize}

\subsection{通风路径}
按冷却气体在电机内部的流动路径，可分为\textbf{径向通风}（气体沿径向穿过铁芯与绕组）、\textbf{轴向通风}（沿轴向流过）和\textbf{混合通风}。大型电机常在铁芯中留出径向通风沟，把铁芯沿轴向分成若干段，以增大散热面积。

\begin{center}
\begin{tabular}{@{}p{2.4cm}p{4.2cm}p{5.2cm}@{}}
\toprule
项目 & 内容 & 关键结论 \\
\midrule
温升 & $\tau=\theta-\theta_0$ & 绝缘老化取决于温度，故标准规定温升限值 \\
绝缘等级 & A/E/B/F/H/C & 最高允许温度 105/120/130/155/180/$>180$\,℃ \\
发热过程 & $\tau=\tau_\infty(1-\mathrm{e}^{-t/T})$ & $3\sim4$ 个时间常数后基本稳定 \\
冷却过程 & $\tau=\tau_0\mathrm{e}^{-t/T}$ & 冷却时间常数通常大于发热时间常数 \\
工作制 & S1 连续 / S2 短时 / S3 断续 & 断续用负载持续率 $FC$ 表征 \\
冷却方式 & 空气、氢、水；开启、防护、封闭 & 容量越大越依赖强化冷却 \\
\bottomrule
\end{tabular}
\end{center}

% ============================================================
% 第 12 章 三相感应电机概述（教材第四篇，印刷页 125--130）
% 节名：第一节 三相感应电机的运行状态；第二节 三相感应电机的结构；第三节 三相感应电动机的额定值
% ============================================================
\chapter{三相感应电机概述}

感应电机（异步电机）多用作电动机，在少数场合也作发电机用。它是产量最大、应用最广的一种电机。本章只讲"是什么、怎么转、铭牌怎么看"，为第 13、14 章的理论分析做准备。

\section{三相感应电机的运行状态}

感应电机的定子接三相电源产生旋转磁场，转子电路自成闭合回路（笼型）或经外接电阻闭合（绕线型）。转子转速 $n$ 与旋转磁场同步转速 $n_1$ 之间的差异，决定了电机的运行状态。定义\textbf{转差率}
\begin{equation}
s=\frac{n_1-n}{n_1},
\end{equation}
则 $n=(1-s)n_1$。

\keyformula{s=\frac{n_1-n}{n_1},\qquad n=(1-s)n_1}

按转差率的正负与大小，感应电机有三种运行状态。

\subsection{电动机状态（$0<s<1$）}
$0<n<n_1$。转子切割旋转磁场的方向与磁场转向相同，转子感应电动势与电流产生的电磁转矩方向与旋转磁场转向相同，是\textbf{驱动}转矩；电机从电网吸收电功率，扣除损耗后以机械功率输出。这是感应电机最主要的运行状态。

\subsection{发电机状态（$s<0$）}
$n>n_1$。转子被原动机拖动超过同步转速，转子导体相对旋转磁场反向切割，感应电动势、电流反向，电磁转矩变为\textbf{制动}转矩，即吸收机械功率、输出电功率。感应发电机需要电网提供励磁（吸收无功），不能单独建压，故应用不如同步发电机广泛，但在风力发电中大量使用（双馈、笼型异步）。

\subsection{电磁制动状态（$s>1$）}
$n<0$（转子逆磁场方向旋转）。此时转子导体切割磁场的方向虽仍使电流产生制动转矩，但电机既吸收电网电功率、又吸收机械功率，两者都变成转子回路的损耗热——这是"反接制动"状态，用于快速停车，不能长期运行。

\begin{center}
\begin{tabular}{@{}llll@{}}
\toprule
运行状态 & 转速范围 & 转差率 & 功率流向 \\
\midrule
电动机 & $0<n<n_1$ & $0<s<1$ & 电网 $\rightarrow$ 机械（输出） \\
发电机 & $n>n_1$ & $s<0$ & 机械 $\rightarrow$ 电网（输出） \\
电磁制动 & $n<0$ & $s>1$ & 电网 $+$ 机械 $\rightarrow$ 转子损耗（发热） \\
\bottomrule
\end{tabular}
\end{center}

\corepoint{转差率是感应电机唯一的"状态变量"：$s$ 一变，转子频率、转子电动势、电流、转矩、功率分配全变。牢记三句话：$s=0$（同步转速）时转子无感应电动势、无转矩；$s=1$（堵转/起动瞬间）时转子电流最大、起动转矩有限；$s$ 略大于 0（额定运行）时效率最高。凡是"转速变化引起什么变化"的问题，一律先写出 $s$，再由 $s$ 推其余量。}

\section{三相感应电机的结构}

感应电机由\textbf{定子}、\textbf{转子}和定转子之间的\textbf{气隙}三部分组成。

\subsection{定子}
定子由定子铁芯、定子绕组和机座构成。
\begin{itemize}
  \item \textbf{定子铁芯}：由 0.5\,mm 厚的硅钢片叠成，内圆均匀开槽，用以嵌放定子绕组，同时构成主磁路。
  \item \textbf{定子绕组}：三相对称绕组，可接成星形或三角形，其六个出线端引至接线盒，便于改变接法（这是 Y—$\Delta$ 降压起动的物理基础）。
  \item \textbf{机座}：铸铁或钢板焊接，起支撑与防护作用；封闭式电机机座外表面带散热筋。
\end{itemize}

\subsection{转子}
按转子绕组形式分为\textbf{笼型}和\textbf{绕线型}两种。

\subsubsection{笼型转子}
转子铁芯外圆开槽，槽内嵌放\textbf{导条}（铜条或铸铝），导条两端用\textbf{端环}短接，形成自行闭合的多相绕组，形似鼠笼。笼型转子结构简单、坚固耐用、价格低廉、无需维护，是应用最广的形式；缺点是起动转矩较小、起动电流较大，且转子回路无法外接电阻。

\subsubsection{绕线型转子}
转子铁芯槽内嵌放三相对称绕组，一般接成星形，三个出线端通过\textbf{滑环}和电刷引至外电路，可串入外加电阻。绕线型转子的优点是可改善起动性能（串电阻提高起动转矩、减小起动电流）并可实现调速（串电阻调速、串级调速）；缺点是结构复杂、成本高、需维护滑环电刷。

\subsection{气隙}
感应电机的气隙很小（中小型电机一般为 0.2～1.5\,mm），因为气隙越大，励磁电流越大、功率因数越低。这是感应电机与同步电机结构上的一个重要差别（同步电机气隙可较大）。

\begin{center}
\begin{tabular}{@{}p{2.4cm}p{4.6cm}p{4.6cm}@{}}
\toprule
项目 & 笼型转子 & 绕线型转子 \\
\midrule
转子绕组 & 导条 + 端环，自行闭合 & 三相绕组，经滑环引出 \\
外接电阻 & 不能 & 可以 \\
起动性能 & 起动电流大、起动转矩小 & 串电阻可提高 $T_{st}$、减小 $I_{st}$ \\
调速 & 主要靠变频、变极 & 可转子串电阻、串级调速 \\
结构与维护 & 简单、坚固、免维护 & 复杂，滑环电刷需维护 \\
典型用途 & 风机、水泵、机床（绝大多数） & 起重、轧机、需大起动转矩场合 \\
\bottomrule
\end{tabular}
\end{center}

\section{三相感应电动机的额定值}

感应电动机铭牌上标明额定值和一些与运行有关的技术数据。按铭牌指定的条件和额定值运行，就叫额定运行。

\subsection{额定功率 $P_N$}
额定运行时电动机轴上输出的\textbf{机械功率}，单位 kW。注意这是\textbf{输出}功率（轴功率），不是输入电功率。

\subsection{额定电压 $U_N$}
额定运行时定子绕组应加的\textbf{线电压}，单位 V 或 kV。

\subsection{额定电流 $I_N$}
额定运行时定子绕组的\textbf{线电流}，单位 A。

\subsection{额定转速 $n_N$}
额定运行时的转子转速，单位 r/min。由 $n_N$ 可算出额定转差率
\begin{equation}
s_N=\frac{n_1-n_N}{n_1},
\end{equation}
其中同步转速 $n_1=60f_N/p$ 由额定频率和极对数确定。

\subsection{额定频率 $f_N$}
我国为 50\,Hz。

\subsection{额定功率因数 $\cos\varphi_N$}
额定运行时定子侧的功率因数。

\subsection{额定效率 $\eta_N$}
额定运行时输出功率与输入功率之比，故输入功率
\begin{equation}
P_{1N}=\frac{P_N}{\eta_N}=\sqrt3U_NI_N\cos\varphi_N .
\end{equation}

\keyformula{P_{1N}=\sqrt3U_NI_N\cos\varphi_N=\frac{P_N}{\eta_N}}

\subsection{接法与绝缘等级}
铭牌还标明定子绕组的接法（Y 或 $\Delta$）与绝缘等级。接法必须与电源线电压相匹配：例如铭牌"380\,V/$\Delta$ 接"，表示电源线电压 380\,V 时应接成三角形。

\begin{center}
\begin{tabular}{@{}lll@{}}
\toprule
额定值 & 含义 & 注意 \\
\midrule
$P_N$ & 轴上输出的机械功率 & 输出，不是输入 \\
$U_N$ & 定子线电压 & 线值 \\
$I_N$ & 定子线电流 & 线值 \\
$n_N$ & 额定转速 & 略低于 $n_1$，$s_N$ 约 1\%～5\% \\
$f_N$ & 额定频率 & 我国 50\,Hz \\
$\cos\varphi_N$ & 定子功率因数 & 中小型约 0.8～0.9 \\
$\eta_N$ & 额定效率 & 中小型约 85\%～93\% \\
接法 & Y 或 $\Delta$ & 必须与电源电压匹配 \\
\bottomrule
\end{tabular}
\end{center}

\corepoint{额定值计算最容易错的两处：① 把 $P_N$ 当成输入功率——由 $P_N$、$\eta_N$ 求 $P_1$ 时才乘/除效率，方向别反；② 用 $U_N$、$I_N$ 算功率时别忘了 $\sqrt3$（线值），而在等效电路里算每相损耗时又要把线值换成相值。记住一条链：$P_1=\sqrt3U_NI_N\cos\varphi_N$，$P_1=P_N+p_{Cu1}+p_{Fe}+p_{Cu2}+p_{mec}+p_{ad}$，$P_N=P_1-\sum p$。}

% ============================================================
% 第 13 章 三相感应电动机的基本理论（教材第四篇，印刷页 131--149）
% 节名：第一节 三相感应电动机的空载运行；第二节 三相感应电动机负载运行；
%       第三节 三相感应电动机的基本方程和等效电路；第四节 感应电动机的功率方程和转矩方程；
%       第五节 笼型转子参数；第六节 三相感应电动机参数测定
% ============================================================
\chapter{三相感应电动机的基本理论}

第 8—10 章说明了三相绕组如何产生旋转磁动势、如何在旋转磁场中感应电动势。本章把这些结论用到感应电动机上：定子产生旋转磁场，旋转磁场在转子中感应电动势与电流，转子电流产生转子磁动势，定、转子磁动势合成建立气隙磁场——如此循环，最终达到稳态。分析套路与变压器几乎完全相同：\textbf{写方程 $\rightarrow$ 归算 $\rightarrow$ 等效电路 $\rightarrow$ 求参数 $\rightarrow$ 算性能}。区别只有一处：感应电机的转子在转，转子量的频率随转差率变化。

\section{三相感应电动机的空载运行}

\subsection{空载磁场与空载电流}
定子接三相电源、轴上不带负载（$T_L=0$）时，电动机以接近同步转速 $n\approx n_1$ 空转，转差率 $s\approx0$。

定子三相电流产生的合成磁动势 $F_1$ 以同步转速 $n_1$ 旋转，建立气隙磁场。与变压器一样，磁通分为两部分：
\begin{itemize}
  \item \textbf{主磁通} $\Phi_m$：同时交链定子绕组与转子绕组，通过气隙与铁芯闭合，是机电能量转换的媒介；其路径受铁芯饱和影响。
  \item \textbf{漏磁通}：只交链定子（或只交链转子）绕组，主要经空气闭合，基本不受饱和影响，用常数漏抗 $X_{1\sigma}$、$X_{2\sigma}$ 表征。
\end{itemize}

空载时转子转速接近同步转速，转子导体与旋转磁场的相对速度极小，转子感应电动势、电流近似为零。因此\textbf{空载电流几乎全部是励磁电流}：
\begin{equation}
\dot I_0=\dot I_m ,
\end{equation}
它由磁化分量 $I_\mu$（建立主磁通，无功）和铁耗分量 $I_{Fe}$（供给铁耗，有功）组成。感应电机气隙小、磁路磁阻小，但励磁电流仍占额定电流的 20\%～50\%（远大于变压器的百分之几），这是感应电机功率因数偏低（空载时 $\cos\varphi_0$ 仅 0.1～0.2）的根本原因。

\subsection{空载时的电动势平衡}
主磁通在定子绕组中感应电动势，有效值为
\begin{equation}
E_1=4.44f_1N_1k_{w1}\Phi_m ,
\end{equation}
其中 $N_1$ 为定子每相串联匝数，$k_{w1}$ 为定子基波绕组因数。定子电压方程为
\begin{equation}
\dot U_1=-\dot E_1+\dot I_0\left(R_1+\mathrm{j}X_{1\sigma}\right)=-\dot E_1+\dot I_0Z_{1\sigma},
\end{equation}
且 $\dot E_1=-\dot I_mZ_m$，$Z_m=R_m+\mathrm{j}X_m$ 为励磁阻抗（$R_m$ 表征铁耗、$X_m$ 表征磁化性能）。

\section{三相感应电动机负载运行}

\subsection{转子频率、电动势与电流}
负载后转速下降为 $n<n_1$，转差率 $s>0$。旋转磁场以相对转速 $n_1-n=sn_1$ 切割转子导体，故转子感应电动势的频率为
\begin{equation}
f_2=\frac{p\left(n_1-n\right)}{60}=\frac{pn_1}{60}\cdot\frac{n_1-n}{n_1}=sf_1 .
\end{equation}

\keyformula{f_2=sf_1}

设转子静止（$s=1$）时转子每相感应电动势为 $E_2$、漏抗为 $X_{2\sigma}$，则转速为 $n$ 时
\begin{equation}
E_{2s}=sE_2,\qquad X_{2\sigma s}=sX_{2\sigma}.
\end{equation}
转子每相电流为
\begin{equation}
I_{2s}=\frac{sE_2}{\sqrt{R_2^{2}+\left(sX_{2\sigma}\right)^{2}}},
\end{equation}
其功率因数角为
\begin{equation}
\varphi_2=\arctan\frac{sX_{2\sigma}}{R_2},
\end{equation}
即\textbf{转子电流与功率因数都随转差率变化}——这一点与变压器完全不同（变压器二次频率恒等于一次频率）。

\subsection{转子磁动势}
转子绕组为三相对称（笼型可等效为多相），转子电流产生的磁动势 $F_2$ 也是一个旋转磁动势，其相对转子的转速为
\begin{equation}
n_{2}=\frac{60f_2}{p}=\frac{60sf_1}{p}=sn_1 .
\end{equation}
而转子本身以转速 $n=(1-s)n_1$ 旋转，转向与旋转磁场相同，故 $F_2$ 在\textbf{空间}的转速为
\begin{equation}
n_{2}+n=sn_1+(1-s)n_1=n_1 .
\end{equation}
即\textbf{转子磁动势与定子磁动势同速、同向旋转，二者相对静止}——这是感应电机能够产生恒定平均转矩的前提。正因如此，才能像变压器一样把定、转子磁动势合成为一个合成磁动势，建立共同的气隙磁场。

\subsection{磁动势平衡}
与变压器同理，负载时定子磁动势与转子磁动势合成，建立与空载时基本相同的主磁通（因为 $U_1$ 不变、$E_1\approx U_1$ 不变），故
\begin{equation}
\dot F_1+\dot F_2=\dot F_0
\quad\Longrightarrow\quad
\dot I_1=\dot I_m+\left(-\dot I_2'\right),
\end{equation}
其中 $\dot I_2'$ 为归算后的转子电流。物理意义：定子电流的负载分量恰好抵消转子磁动势的去磁作用，维持主磁通不变。

\corepoint{这一节最重要的一句话是"\textbf{转子磁动势在空间也是以同步转速旋转}"。它一举解决三个疑问：① 为什么定、转子磁场能相对静止、产生恒定转矩；② 为什么转子频率随 $s$ 变，而气隙磁场转速恒为 $n_1$；③ 为什么可以用变压器的方法做归算。凡是"转速变了磁场还同步吗"的疑问，都用 $n_2+n=n_1$ 回答。}

\section{三相感应电动机的基本方程和等效电路}

\subsection{电压方程}
定子（与变压器一次侧完全相同）：
\begin{equation}
\dot U_1=-\dot E_1+\dot I_1\left(R_1+\mathrm{j}X_{1\sigma}\right).
\end{equation}
转子（注意频率为 $f_2=sf_1$，电动势为 $sE_2$、漏抗为 $sX_{2\sigma}$）：
\begin{equation}
\dot E_{2s}=\dot I_{2s}\left(R_2+\mathrm{j}sX_{2\sigma}\right),
\end{equation}
其中 $\dot E_{2s}=s\dot E_2$。两边同除以 $s$，得
\begin{equation}
\dot E_2=\dot I_{2s}\left(\frac{R_2}{s}+\mathrm{j}X_{2\sigma}\right).
\end{equation}
这一步非常关键：\textbf{把"转子在转"的问题转化成了"转子静止"的问题}，代价是转子电阻由 $R_2$ 变为 $R_2/s$。

\subsection{归算}
与变压器相同，把转子各量归算到定子侧。由于感应电机定、转子相数可能不同（$m_1$、$m_2$）、绕组因数不同，归算系数为
\begin{equation}
k_e=\frac{N_1k_{w1}}{N_2k_{w2}},\qquad
k_i=\frac{m_1N_1k_{w1}}{m_2N_2k_{w2}},\qquad
k_z=k_ek_i .
\end{equation}
归算规律与变压器一致：电动势、电压乘 $k_e$；电流除 $k_i$；阻抗乘 $k_ek_i$。

笼型转子的处理见第 13.5 节。

\subsection{T 型等效电路}
归算后可得与变压器结构完全相同的 T 型等效电路：定子漏阻抗 $Z_{1\sigma}=R_1+\mathrm{j}X_{1\sigma}$ 串联，励磁支路 $Z_m=R_m+\mathrm{j}X_m$ 并在中间，转子支路为
\begin{equation}
Z_2'=\frac{R_2'}{s}+\mathrm{j}X_{2\sigma}' .
\end{equation}

\keyformula{Z_2'=\frac{R_2'}{s}+\mathrm{j}X_{2\sigma}'=\underbrace{R_2'}_{\text{转子铜耗}}+\underbrace{\frac{1-s}{s}R_2'}_{\text{附加电阻（代表机械功率）}}+\mathrm{j}X_{2\sigma}'}

把 $R_2'/s$ 拆成两项，是感应电机等效电路最有物理意义的一步：
\begin{itemize}
  \item $R_2'$ 上的损耗对应\textbf{转子铜耗} $p_{Cu2}$（真实存在的热损耗）；
  \item 附加电阻 $\dfrac{1-s}{s}R_2'$ 上的"损耗"并不发热，它代表\textbf{转换成机械功率的那部分电磁功率}（等效静止转子的机械负载）。
\end{itemize}

\subsection{基本方程组}
\begin{equation}
\left\{
\begin{aligned}
\dot U_1&=-\dot E_1+\dot I_1\left(R_1+\mathrm{j}X_{1\sigma}\right),\\
\dot E_1&=-\dot I_mZ_m,\\
\dot I_1&=\dot I_m+\left(-\dot I_2'\right),\\
\dot E_1&=\dot E_2'=\dot I_2'\left(\frac{R_2'}{s}+\mathrm{j}X_{2\sigma}'\right).
\end{aligned}
\right.
\end{equation}

\section{感应电动机的功率方程和转矩方程}

\subsection{功率流程}
输入电功率
\begin{equation}
P_1=m_1U_1I_1\cos\varphi_1 .
\end{equation}
扣除定子铜耗 $p_{Cu1}=m_1I_1^2R_1$ 与铁耗 $p_{Fe}=m_1I_m^2R_m$ 后，余下的是通过气隙传递到转子的\textbf{电磁功率}
\begin{equation}
P_{em}=P_1-p_{Cu1}-p_{Fe}=m_1I_2'^{2}\frac{R_2'}{s} .
\end{equation}
电磁功率中，一部分是转子铜耗
\begin{equation}
p_{Cu2}=m_1I_2'^{2}R_2'=sP_{em},
\end{equation}
余下的是\textbf{总机械功率}
\begin{equation}
P_{mec}=P_{em}-p_{Cu2}=\left(1-s\right)P_{em} .
\end{equation}
再扣除机械损耗 $p_{mec}$ 与杂散损耗 $p_{ad}$，即得轴上输出功率
\begin{equation}
P_2=P_{mec}-p_{mec}-p_{ad} .
\end{equation}

\keyformula{P_{em}:p_{Cu2}:P_{mec}=1:s:(1-s)}

这个比例关系是感应电机最漂亮的结论之一：\textbf{转差率 $s$ 就是"转子铜耗占电磁功率的比例"}。$s$ 越大，损耗比例越大、效率越低，所以感应电机额定转差率只有 1\%～5\%。

\subsection{转矩方程}
由 $P_{em}=T\Omega_1$（$\Omega_1=2\pi n_1/60$ 为同步机械角速度）与 $P_{mec}=T\Omega$（$\Omega=2\pi n/60$），得
\begin{equation}
T=\frac{P_{em}}{\Omega_1}=\frac{P_{mec}}{\Omega}.
\end{equation}
稳态时 $T=T_2+T_0$（输出转矩 $T_2=P_2/\Omega$ 加空载制动转矩 $T_0$）。

\subsection{电磁转矩的实用表达式}
由 $P_{em}=m_1I_2'^2R_2'/s$ 与等效电路解得 $I_2'$，代入 $T=P_{em}/\Omega_1$，得
\begin{equation}
T=\frac{m_1pU_1^{2}\dfrac{R_2'}{s}}{\Omega_1\left[\left(R_1+\dfrac{R_2'}{s}\right)^{2}+\left(X_{1\sigma}+X_{2\sigma}'\right)^{2}\right]},
\qquad \Omega_1=\frac{2\pi f_1}{p}.
\end{equation}

\keyformula{T=\frac{m_1pU_1^{2}R_2'/s}{\Omega_1\left[\left(R_1+R_2'/s\right)^{2}+\left(X_{1\sigma}+X_{2\sigma}'\right)^{2}\right]}}

其中 $U_1$ 为定子\textbf{相电压}，$m_1$ 为定子相数。这个式子是第 14 章分析转矩—转差率曲线、最大转矩、起动转矩的出发点，必须会推导也要会背。

\section{笼型转子参数}

笼型转子没有明显的"相"，需等效为多相绕组：
\begin{itemize}
  \item 导条数 $Q_2$ 就是等效相数，即 $m_2=Q_2$；
  \item 每相只有一根导条，故每相串联匝数 $N_2=\tfrac12$（一根导条构成"一去一回"的半匝）；
  \item 绕组因数 $k_{w2}=1$。
\end{itemize}
代入归算系数，得笼型转子的电压比、电流比与阻抗归算系数：
\begin{equation}
k_e=\frac{N_1k_{w1}}{N_2k_{w2}}=2N_1k_{w1},\qquad
k_i=\frac{m_1N_1k_{w1}}{Q_2\cdot\frac12\cdot1}=\frac{2m_1N_1k_{w1}}{Q_2},\qquad
k_z=k_ek_i .
\end{equation}
这表明：\textbf{笼型转子不需要"制造"相数和匝数，它们由导条数自然确定}——这是笼型电机设计上的便利之处，也使笼型转子可以用铸铝一次浇成。

\section{三相感应电动机参数测定}

与三相变压器相似，用\textbf{空载试验}和\textbf{堵转试验}测定参数。

\subsection{空载试验}
\textbf{目的}：求励磁参数 $Z_m$、$R_m$、$X_m$，并分离铁耗与机械损耗。

\textbf{方法}：电动机轴上不带负载，定子加对称三相电压，从（1.1～1.3）$U_N$ 逐步降压到电机转速明显下降为止，测量 $U_1$、$I_0$、$p_0$。

\textbf{数据处理}：
\begin{enumerate}
  \item 空载时 $s\approx0$，转子支路近似开路，故测得的是励磁支路参数：
  \begin{equation}
  Z_0=\frac{U_{1\phi}}{I_{0\phi}},\qquad R_0=\frac{p_0'}{3I_{0\phi}^{2}},\qquad X_0=\sqrt{Z_0^{2}-R_0^{2}},
  \end{equation}
  再扣除定子漏阻抗得 $X_m\approx X_0-X_{1\sigma}$、$R_m\approx R_0-R_1$。
  \item \textbf{分离铁耗与机械损耗}：空载输入功率 $p_0$ 包含定子铜耗 $3I_0^2R_1$、铁耗 $p_{Fe}$ 和机械损耗 $p_{mec}$。令
  \begin{equation}
  p_0'=p_0-3I_0^{2}R_1=p_{Fe}+p_{mec} .
  \end{equation}
  铁耗 $\propto U_1^2$，机械损耗与电压无关（转速基本不变）。作 $p_0'=f(U_1^2)$ 曲线，其\textbf{延长线与纵轴的交点}即机械损耗 $p_{mec}$，从而分离出额定电压下的 $p_{Fe}$。
\end{enumerate}

\subsection{堵转试验}
\textbf{目的}：求短路参数 $R_k=R_1+R_2'$、$X_k=X_{1\sigma}+X_{2\sigma}'$。

\textbf{方法}：卡住转子（$n=0$，$s=1$），定子加降低的对称三相电压，使电流约为额定值，测量 $U_k$、$I_k$、$p_k$。

\textbf{数据处理}：
\begin{equation}
Z_k=\frac{U_{k\phi}}{I_{k\phi}},\qquad
R_k=\frac{p_k}{3I_{k\phi}^{2}},\qquad
X_k=\sqrt{Z_k^{2}-R_k^{2}} .
\end{equation}
因 $R_1$ 可由直流电阻测得，故 $R_2'=R_k-R_1$；而 $X_{1\sigma}$ 与 $X_{2\sigma}'$ 难以分离，工程上按
\begin{equation}
X_{1\sigma}\approx X_{2\sigma}'\approx\frac{X_k}{2}
\end{equation}
处理（对大型电机可取 $X_{1\sigma}=0.4X_k$、$X_{2\sigma}'=0.6X_k$ 之类的经验分配）。

堵转时电流大、磁路饱和，漏抗比正常运行值小，故测得的 $X_k$ 需按额定电流下的饱和程度修正；此外堵转试验应在尽可能低的电压、短时间内完成，以免绕组过热。

\begin{center}
\begin{tabular}{@{}p{2.0cm}p{3.6cm}p{3.8cm}p{3.6cm}@{}}
\toprule
试验 & 条件 & 主要求得 & 关键判据 \\
\midrule
空载试验 & 轴上无负载，$s\approx0$ & $Z_m$、$R_m$、$X_m$；分离 $p_{Fe}$ 与 $p_{mec}$ & 转子支路近似开路 \\
堵转试验 & 转子卡住，$n=0$、$s=1$ & $R_k$、$X_k$，进而 $R_2'$ & 转子支路阻抗最小 \\
\bottomrule
\end{tabular}
\end{center}

\corepoint{三个易错点：① 空载试验得到的 $p_0$ 不能全算作铁耗，必须扣除定子铜耗 $3I_0^2R_1$ 并分离机械损耗——这是与变压器空载试验最大的区别（变压器没有机械损耗）；② 堵转试验对应 $s=1$，不是"短路"，切勿与变压器短路试验的"二次短路"混淆，物理上是同一个意思但要清楚 $R_2'/s=R_2'$；③ 所有公式里的 $U_1$、$I_1$ 若用铭牌数据必须\textbf{先换算成相值}（Y 接 $U_\phi=U_N/\sqrt3$，$\Delta$ 接 $U_\phi=U_N$、$I_\phi=I_N/\sqrt3$），这一步每年都有人栽跟头。}

% ============================================================
% 第 14 章 三相感应电动机的运行特性（教材第四篇，印刷页 150--166）
% 节名：第一节 三相感应电动机转矩—转差率曲线；第二节 三相感应电动机的机械特性和工作特性；
%       第三节 深槽和双笼感应电动机；第四节 三相感应电动机的启动及电磁制动；第五节 三相感应电动机的调速
% ============================================================
\chapter{三相感应电动机的运行特性}

第 13 章给出了电磁转矩的表达式。本章把它画成曲线（转矩—转差率曲线、机械特性），再由曲线解释三个工程问题：\textbf{能不能稳定运行}、\textbf{怎么起动}、\textbf{怎么调速}。这三个问题也是感应电动机全部工程应用的落脚点。

\section{三相感应电动机转矩—转差率曲线}

在额定电压和额定频率下，电磁转矩与转差率的关系 $T=f(s)$ 称为\textbf{转矩—转差率曲线}。由第 13 章
\begin{equation}
T=\frac{m_1pU_1^{2}\dfrac{R_2'}{s}}{\Omega_1\left[\left(R_1+\dfrac{R_2'}{s}\right)^{2}+\left(X_{1\sigma}+X_{2\sigma}'\right)^{2}\right]} .
\end{equation}

\subsection{曲线的三个特征点}
\begin{enumerate}
  \item \textbf{同步点}（$s=0$，$n=n_1$）：$T=0$。转子与旋转磁场无相对运动，无感应电动势、无电流、无转矩。
  \item \textbf{最大转矩点}（$s=s_m$）：$T=T_{\max}$。对 $s$ 求导并令其为零，得\textbf{临界转差率}
  \begin{equation}
  s_m=\frac{R_2'}{\sqrt{R_1^{2}+\left(X_{1\sigma}+X_{2\sigma}'\right)^{2}}} ,
  \end{equation}
  代回转矩式得\textbf{最大转矩}
  \begin{equation}
  T_{\max}=\frac{m_1pU_1^{2}}{2\Omega_1\left[R_1+\sqrt{R_1^{2}+\left(X_{1\sigma}+X_{2\sigma}'\right)^{2}}\right]} .
  \end{equation}
  \item \textbf{起动点}（$s=1$，$n=0$）：$T=T_{st}$ 称为\textbf{起动转矩}，是曲线与 $s=1$ 的交点。
\end{enumerate}

\keyformula{s_m=\frac{R_2'}{\sqrt{R_1^{2}+X_k^{2}}},\qquad
T_{\max}=\frac{m_1pU_1^{2}}{2\Omega_1\left[R_1+\sqrt{R_1^{2}+X_k^{2}}\right]},\qquad X_k=X_{1\sigma}+X_{2\sigma}'}

\subsection{两个重要比例关系}
\begin{itemize}
  \item $T_{\max}\propto U_1^{2}$：\textbf{最大转矩与电压平方成正比}。电网电压降低 10\%，最大转矩降低约 19\%，这是电动机"低电压时容易堵转烧毁"的原因。
  \item $s_m\propto R_2'$：\textbf{临界转差率与转子电阻成正比}，而 $T_{\max}$ \textbf{与 $R_2'$ 无关}（忽略 $R_1$ 时）。这一条是绕线型电动机"转子串电阻改善起动性能"的理论依据：增大 $R_2'$ 可使 $s_m$ 向 1 靠近，从而提高起动转矩，而最大转矩不变。
\end{itemize}

\subsection{稳定运行区}
把曲线分为两段：
\begin{itemize}
  \item $0<s<s_m$（转速较高段）：$T$ 随 $s$ 增大而增大，$\dfrac{\mathrm{d}T}{\mathrm{d}n}<0$，为\textbf{稳定运行区}。负载增加时转速略降、转矩增大，自动达到新的平衡。
  \item $s_m<s<1$（转速较低段）：$T$ 随 $s$ 增大反而减小，为\textbf{不稳定区}。若负载转矩超过 $T_{\max}$，电机将堵转。
\end{itemize}
稳定运行的判据是：负载扰动消失后能回到原平衡点，即 $\dfrac{\mathrm{d}T}{\mathrm{d}n}<\dfrac{\mathrm{d}T_L}{\mathrm{d}n}$。额定运行点一般在 $s_N=0.01\sim0.05$，距 $s_m$ 尚有裕量，故过载能力
\begin{equation}
\lambda=\frac{T_{\max}}{T_N}=1.6\sim2.5 .
\end{equation}

\section{三相感应电动机的机械特性和工作特性}

\subsection{机械特性}
把 $T=f(s)$ 换成 $n=f(T)$ 即得\textbf{机械特性}。它近似为一条略微下倾的直线（工作段），在 $T_{\max}$ 处出现拐点，之后转速迅速下降到零（堵转）。

机械特性是分析电动机与负载配合的工具：把电动机的 $n=f(T)$ 与负载的 $n=f(T_L)$ 画在同一图上，交点即稳态工作点。

\subsection{工作特性}
工作特性指 $U_1=U_N$、$f_1=f_N$ 时，转速 $n$、电磁转矩 $T$、定子电流 $I_1$、功率因数 $\cos\varphi_1$、效率 $\eta$ 随输出功率 $P_2$ 变化的曲线。

\begin{itemize}
  \item \textbf{转速特性} $n=f(P_2)$：从空载到额定负载，转速下降很少（$s$ 由 0 增到 $s_N$），曲线略向下倾斜，属"硬特性"。
  \item \textbf{转矩特性} $T=f(P_2)$：$T=T_0+P_2/\Omega$，近似为一条略上翘的直线。
  \item \textbf{定子电流特性} $I_1=f(P_2)$：空载时 $I_1=I_0$（较大，占额定 20\%～50\%）；随负载增加，$I_1$ 近似按线性增长。
  \item \textbf{功率因数特性} $\cos\varphi_1=f(P_2)$：空载时很低（0.1～0.2，因励磁电流大），随负载增加迅速上升，在额定负载附近达最大（0.8～0.9），过载后又下降。
  \item \textbf{效率特性} $\eta=f(P_2)$：空载时为零，随负载增加而上升，在可变损耗（铜耗）等于不变损耗（铁耗+机械损耗）时达到最大，一般在 $(0.7\sim0.8)P_N$ 附近。
\end{itemize}

\begin{center}
\begin{tabular}{@{}lll@{}}
\toprule
特性 & 变化规律 & 原因 \\
\midrule
$n=f(P_2)$ & 略下降（硬特性） & $s_N$ 仅 1\%～5\% \\
$T=f(P_2)$ & 近似直线略上翘 & $T=T_0+P_2/\Omega$ \\
$I_1=f(P_2)$ & 空载较大，随负载近线性增大 & 定子电流含励磁分量与负载分量 \\
$\cos\varphi_1=f(P_2)$ & 空载低、额定附近最高 & 励磁电流占比随负载下降 \\
$\eta=f(P_2)$ & 先升后降，$(0.7\sim0.8)P_N$ 最高 & 可变损耗 $=$ 不变损耗 \\
\bottomrule
\end{tabular}
\end{center}

\corepoint{感应电动机"轻载运行不经济"的根子在功率因数：空载时 $I_1$ 中绝大部分是无功的励磁电流，$\cos\varphi_1$ 只有 0.1～0.2；负载增加后有功分量上升，$\cos\varphi_1$ 才提高。所以"大马拉小车"不仅效率低，还使电网功率因数变差——这是工程上按负载选配电机容量的理由。}

\section{深槽和双笼感应电动机}

普通笼型电动机起动时 $s=1$，$X_{2\sigma}$ 相对较大、转子功率因数低，起动转矩不大而起动电流很大（$I_{st}=(4\sim7)I_N$）。改善起动性能的关键是\textbf{起动时增大转子电阻、运行时恢复小电阻}——利用\textbf{集肤效应}可以自动实现这一点。

\subsection{集肤效应}
当转子导条中电流频率较高时，交流电流因电感作用趋向于导条\textbf{上部（槽口侧）}流动，有效截面减小、等效电阻增大、漏抗减小，这种现象称为\textbf{集肤效应}（挤流效应）。集肤效应的强弱与频率有关：频率越高越显著，频率越低越不明显。

而转子频率 $f_2=sf_1$：\textbf{起动瞬间} $s=1$，$f_2=f_1=50$\,Hz，集肤效应强，转子电阻大 $\Rightarrow$ 起动转矩大、起动电流小；\textbf{正常运行} $s$ 很小，$f_2$ 仅 1～3\,Hz，集肤效应几乎消失，转子电阻回到小值 $\Rightarrow$ 效率高、转差率小。这就是深槽与双笼转子的巧妙之处：\textbf{不需要外部控制，靠频率自动切换}。

\subsection{深槽转子}
转子槽形做得又深又窄（高宽比可达 10～12），导条为高而窄的矩形。起动时集肤效应使电流集中在上部，等效电阻大；运行时电流均匀分布，电阻小。

\subsection{双笼转子}
转子上有两套笼：\textbf{上笼}（起动笼）用黄铜或铝青铜等电阻率高的材料，截面小、靠近槽口，漏抗小；\textbf{下笼}（运行笼）用紫铜，截面大、埋在槽底，漏抗大。
\begin{itemize}
  \item 起动时 $f_2$ 高，下笼漏抗大、电流被挤到上笼，上笼电阻大 $\Rightarrow$ \textbf{起动转矩大}；
  \item 运行时 $f_2$ 低，两笼电流按电阻反比分配，下笼电阻小、承担大部分电流 $\Rightarrow$ \textbf{运行效率高}。
\end{itemize}
双笼电动机的转矩—转差率曲线呈现"两段上凸"的形状，兼顾了起动与运行性能。

\section{三相感应电动机的启动及电磁制动}

\subsection{对起动性能的要求}
\begin{enumerate}
  \item \textbf{起动转矩足够大}（$T_{st}/T_N$ 一般要求 1.0～2.0），能带动负载起转；
  \item \textbf{起动电流不要太大}（普通笼型 $I_{st}/I_N=4\sim7$），以免电网电压跌落、绕组过热；
  \item 起动设备简单、经济、可靠。
\end{enumerate}
起动瞬间 $s=1$、转子支路阻抗最小，故电流最大而转矩并不最大——这是感应电动机起动问题的根源。

\subsection{笼型电动机的起动方法}

\subsubsection{直接起动（全压起动）}
把定子直接接入额定电压。方法最简单、起动转矩最大，但起动电流大。适用于小容量电动机（一般满足"电动机容量不超过电源变压器容量的 20\%～30\%"这类经验条件）。

\subsubsection{星—三角（Y—$\Delta$）降压起动}
运行时为 $\Delta$ 接的电动机，起动时先接成 Y，起动完毕再切换为 $\Delta$。由于 Y 接时每相电压降为 $1/\sqrt3$，故
\begin{equation}
\frac{I_{stY}}{I_{st\Delta}}=\frac13,\qquad
\frac{T_{stY}}{T_{st\Delta}}=\left(\frac{1}{\sqrt3}\right)^{2}=\frac13 .
\end{equation}
即\textbf{起动电流和起动转矩都降为直接起动的 1/3}。设备简单、成本低，但只适用于正常运行时为 $\Delta$ 接的电机，且起动转矩降低较多，只宜轻载起动。

\subsubsection{自耦变压器降压起动}
用自耦变压器把电压降低后加到定子，起动完毕再切换到全压。若抽头电压比为 $k_A<1$，则
\begin{equation}
\frac{I_{st}}{I_{st,\text{直接}}}=k_A^{2},\qquad
\frac{T_{st}}{T_{st,\text{直接}}}=k_A^{2} .
\end{equation}
有 55\%、64\%、73\% 等抽头可选，比 Y—$\Delta$ 灵活，可用于 Y 接或 $\Delta$ 接的电机；缺点是设备体积大、价格高。

\subsubsection{软起动与变频起动}
用晶闸管调压（软起动器）使电压平滑上升，或用变频器从低频低压起动，可把起动电流限制在 1.5～2 倍额定值以内，且起动平滑、无冲击，是现代主流方案。

\subsection{绕线型电动机的起动}
绕线型转子可外串电阻。由 $s_m\propto R_2'$、$T_{\max}$ 不变可知：串入适当电阻既\textbf{减小起动电流}，又\textbf{提高起动转矩}（可做到 $T_{st}=T_{\max}$）。起动过程中逐段切除电阻，可使转矩保持在较大范围内，是起重、轧钢等重载起动的理想方案。

\subsection{电磁制动}
\begin{enumerate}
  \item \textbf{反接制动}：把定子任意两相对调，旋转磁场反向，转差率 $s>1$，转矩变为制动转矩，转速迅速下降；到零附近必须及时切断电源，否则电机会反转。制动强烈但能量全部消耗在转子回路。
  \item \textbf{能耗制动}：切断三相交流电源，立即在定子绕组中通入直流电流，产生静止磁场，转子切割该磁场产生制动转矩。制动平稳、无冲击，能量也变为热。
  \item \textbf{回馈制动}：电机转速超过同步转速（$s<0$），进入发电机状态，把机械能回馈电网。多用于起重机下放重物、变频调速减速过程，最经济。
\end{enumerate}

\section{三相感应电动机的调速}

由 $n=(1-s)n_1=(1-s)\dfrac{60f_1}{p}$ 可知，调速只有三条路径：\textbf{改变极对数 $p$}、\textbf{改变频率 $f_1$}、\textbf{改变转差率 $s$}。

\subsection{变极调速}
改变定子绕组的接法（如 Y—YY、$\Delta$—YY）使极对数成倍变化，同步转速随之成倍变化，故为\textbf{有级调速}。优点是不需额外设备、效率高；缺点是级差大、不能连续调节，多用于机床、风机等。

\subsection{变频调速}
改变供电频率 $f_1$ 即可连续改变 $n_1$，是目前性能最好的调速方式。关键约束是\textbf{保持磁通恒定}：由 $U_1\approx E_1=4.44f_1N_1k_{w1}\Phi_m$ 可知，若只降频不降压，$\Phi_m$ 将增大导致磁路饱和、励磁电流剧增，故需按
\begin{equation}
\frac{U_1}{f_1}=\text{常数}
\end{equation}
进行"恒磁通"控制（基频以下恒转矩调速）；基频以上因电压不能超过额定值，转为"恒功率"调速（弱磁升速）。变频调速范围宽、平滑、效率高，是现代交流调速的主流。

\subsection{改变转差率调速}
\begin{itemize}
  \item \textbf{转子串电阻调速}（绕线型）：串电阻使 $s_m$ 增大，同一转矩下转差率增大、转速下降。设备简单，但转差功率消耗在电阻上，效率随转速降低而下降。
  \item \textbf{串级调速}（绕线型）：把转子回路的转差功率经整流、逆变后回馈电网或用于驱动，效率高，但设备复杂。
  \item \textbf{改变定子电压调速}：降低 $U_1$ 使 $T_{\max}\propto U_1^2$ 下降，工作点沿负载特性下移，转速降低。方法简单但调速范围小、低速时损耗大，多用于风机类平方转矩负载。
\end{itemize}

\begin{center}
\begin{tabular}{@{}p{2.3cm}p{3.4cm}p{3.8cm}p{3.4cm}@{}}
\toprule
调速方式 & 依据 & 特点 & 适用 \\
\midrule
变极 & 改变 $p$，$n_1\propto1/p$ & 有级、高效、便宜 & 机床、风机 \\
变频 & 改变 $f_1$，$n_1\propto f_1$ & 无级、范围宽、效率高 & 现代主流，需变频器 \\
转子串电阻 & 改变 $s$（$s_m\propto R_2'$） & 简单，转差功率损耗大 & 绕线型，短时调速 \\
串级 & 转差功率回馈 & 效率高、设备复杂 & 绕线型大容量 \\
改变定子电压 & $T_{\max}\propto U_1^2$ & 范围小、低速损耗大 & 风机、泵类 \\
\bottomrule
\end{tabular}
\end{center}

\corepoint{本章的两条主线要抓牢：① \textbf{转矩公式的三个比例}——$T\propto U_1^2$、$T_{\max}\propto U_1^2$ 且与 $R_2'$ 无关、$s_m\propto R_2'$，几乎所有起动与调速问题都由这三条推出；② \textbf{起动与调速都在"动 $s$ 或动 $n_1$"}——降压、串电阻是动 $s$（并牺牲转矩或效率），变极、变频是动 $n_1$（不改转差率，故效率高）。答题时先说清"我动的是 $n_1$ 还是 $s$"，思路就不会乱。}

% ============================================================
% 第 15 章 其他常用感应电机（教材第四篇，印刷页 167--174）
% 节名：第一节 单相感应电动机；第二节 直线感应电机；第三节 感应发电机
% ============================================================
\chapter{其他常用感应电机}

三相感应电动机之外，还有单相感应电动机（家用电器、小功率场合）、直线感应电机（直线运动场合）和感应发电机（风力发电等）。它们的基本原理仍由第 10 章"脉振磁动势 = 两个反向旋转磁动势"这一条推出。

\section{单相感应电动机}

\subsection{单相绕组的磁场}
单相感应电动机定子上只有一相工作绕组（单相），通入单相交流电流后产生的是\textbf{脉振磁动势}，其轴线固定、幅值随时间脉振，而不是旋转磁场。由第 10 章第二节，
\begin{equation}
f_{\phi1}(x,t)=\frac{F_{\phi1}}{2}\sin\left(\omega t-\frac{\pi}{\tau}x\right)
+\frac{F_{\phi1}}{2}\sin\left(\omega t+\frac{\pi}{\tau}x\right),
\end{equation}
即一个脉振磁场可分解为\textbf{两个幅值相等（各为 $F_{\phi1}/2$）、转速相同（均为 $n_1$）、转向相反}的旋转磁场。

\subsection{为什么没有起动转矩}
设转子转速为 $n$，则转子对正转磁场的转差率为
\begin{equation}
s_+=\frac{n_1-n}{n_1}=s,
\end{equation}
对反转磁场的转差率为
\begin{equation}
s_-=\frac{n_1+n}{n_1}=2-s .
\end{equation}
正转磁场产生正向转矩 $T_+$，反转磁场产生反向转矩 $T_-$。在起动瞬间 $n=0$，$s_+=s_-=1$，两磁场对称，故 $T_+=T_-$，合成转矩
\begin{equation}
T=T_++T_-=0 .
\end{equation}
\textbf{单相感应电动机没有起动转矩}，不能自行起动——这是它区别于三相感应电动机的根本特点。

若用外力使转子转起来（$n>0$），则 $s_+<1$、$s_->1$，正转转矩大于反转转矩，合成转矩为正，电机可继续加速到接近同步转速运行；转向取决于初始转动方向，故单相电动机"往哪边拨就往哪边转"。

\subsection{起动方法}
解决起动问题的思路只有一个：\textbf{人为制造一个在空间上不同相位的第二套绕组电流，使气隙磁场"转"起来}。

\subsubsection{分相式（电阻分相）}
在定子上另加一套\textbf{起动绕组}，其在空间上与工作绕组相差 $90^\circ$ 电角度；起动绕组用细导线（电阻大、电抗小），使其电流与工作绕组电流在时间上接近 $90^\circ$ 相位差，从而产生椭圆形旋转磁场、产生起动转矩。起动后由离心开关切除起动绕组。

\subsubsection{电容分相式}
在起动绕组中串入电容器，使起动绕组电流\textbf{超前}工作绕组电流接近 $90^\circ$，可获得接近圆形的旋转磁场，起动转矩较大、起动电流较小。若电容长期接入，则称\textbf{电容运转式}，可提高运行功率因数与效率。

\subsubsection{罩极式}
在凸极磁极的一部分套上短路铜环（罩极环），环中感应电流使被罩部分的磁通滞后，形成"移行磁场"而产生起动转矩。结构最简单、最可靠，但起动转矩小、效率低，多用于小风扇。

\begin{center}
\begin{tabular}{@{}p{2.2cm}p{3.6cm}p{3.2cm}p{3.4cm}@{}}
\toprule
起动方式 & 原理 & 起动转矩 & 特点 \\
\midrule
电阻分相 & 起动绕组电阻大，电流相位差 $<90^\circ$ & 中等 & 简单便宜，离心开关切除 \\
电容分相 & 串电容，相位差接近 $90^\circ$ & 较大 & 可做成电容运转式提高功率因数 \\
罩极式 & 短路环使磁通移相 & 小 & 结构最简单，效率低 \\
\bottomrule
\end{tabular}
\end{center}

\corepoint{单相感应电动机是"脉振磁动势"最典型的应用：正、反两个旋转磁场叠加 $\Rightarrow$ 起动时两转矩相等 $\Rightarrow$ 无起动转矩；转动起来后 $s_+\ne s_-$ $\Rightarrow$ 有净转矩。记住 $s_+=s$、$s_-=2-s$ 这一对式子，几乎所有单相电机的定性题都能推出答案。}

\section{直线感应电机}

直线感应电机的定子被"剖开铺平"，成为\textbf{初级}（原边）；转子相应成为\textbf{次级}（副边），可以是金属板（如铝板、钢板）。初级通三相电流产生\textbf{行波磁场}（而非旋转磁场），在次级中感应电流，产生直线方向的电磁推力，使次级（或初级）做直线运动。

主要特点：
\begin{itemize}
  \item 直接产生直线运动，不需旋转—直线转换机构（如丝杠、齿条），结构简单、无磨损、可靠性高；
  \item 存在\textbf{端部效应}（初级两端磁场畸变）和\textbf{边缘效应}（横向磁场不均），使效率与功率因数低于旋转电机；
  \item 气隙较大（机械上必须留出间隙），励磁电流大、功率因数低。
\end{itemize}
应用于地铁与磁浮列车牵引、传送带、电磁弹射、电梯等。

\section{感应发电机}

感应发电机由原动机拖动到\textbf{超过同步转速}（$s<0$）运行，定子接电网（或接电容器自励），向电网输出有功功率。特点：
\begin{itemize}
  \item \textbf{需要无功励磁}：感应发电机没有独立的励磁绕组，靠电网（或并联电容器）提供励磁电流，故输出有功的同时吸收无功，功率因数滞后；
  \item \textbf{不能独立建压}：除非并联电容器自励（自励感应发电机），否则必须并网运行；
  \item \textbf{结构简单、免维护、无需同步装置}：这是它在风力发电（尤其是早期定桨距风机）中被大量采用的原因；
  \item 转速变化范围小（$s$ 绝对值很小），效率略低于同步发电机。
\end{itemize}

双馈感应发电机（绕线型转子经变频器接电网）可在变速范围内运行并独立调节有功、无功，是现代变速风电机组的主流机型。

% ============================================================
% 第 16 章 同步电机概述（教材第五篇，印刷页 176--182）
% 节名：第一节 同步电机的基本结构；第二节 同步电机的运行状态；第三节 同步电机的励磁方式；第四节 同步电机的型号和额定值
% ============================================================
\chapter{同步电机概述}

同步电机是电力系统的核心设备：几乎全部交流电能都由同步发电机产生。它的特点是\textbf{转速恒等于同步转速} $n_1=60f/p$，转子转速与旋转磁场"同步"。本章介绍结构与额定值，第 17 章讲基本理论，第 18 章讲并网运行，第 19 章讲电动机与补偿机。

\section{同步电机的基本结构}

同步电机由\textbf{定子}（电枢）、\textbf{转子}（磁极）和气隙组成。定子结构与感应电机基本相同：铁芯由硅钢片叠成，内圆开槽嵌放三相对称电枢绕组。区别主要在转子：同步电机转子装有\textbf{直流励磁的主磁极}。

\subsection{隐极式转子}
转子为圆柱形整体（由整块高强度合金钢锻成），励磁绕组嵌放在转子表面的槽内，槽的分布使气隙磁密接近正弦。特点：
\begin{itemize}
  \item 气隙均匀、转子直径小、轴向长度大；
  \item 机械强度高，可高速旋转（3000\,r/min、1500\,r/min）；
  \item 用于\textbf{汽轮发电机}（火力发电、核电）。
\end{itemize}

\subsection{凸极式转子}
转子磁极是凸出的（磁极铁芯由薄钢板叠成，套上励磁绕组），磁极与磁极之间有较大的极间空隙。特点：
\begin{itemize}
  \item 气隙不均匀（极下小、极间大），$d$、$q$ 轴磁阻不同（这是凸极同步电机分析的出发点）；
  \item 转子直径大、轴向长度短，转速低（几十到几百 r/min）；
  \item 用于\textbf{水轮发电机}及柴油发电机。
\end{itemize}

\begin{center}
\begin{tabular}{@{}p{2.2cm}p{4.2cm}p{4.2cm}@{}}
\toprule
项目 & 隐极式 & 凸极式 \\
\midrule
转子形状 & 圆柱形整体锻钢，槽内嵌励磁绕组 & 凸出磁极，极间空隙大 \\
气隙 & 均匀 & 不均匀（$d$、$q$ 轴磁阻不同） \\
转速 & 高（1500、3000\,r/min） & 低（几十～几百 r/min） \\
极数 & 少（2、4 极） & 多（可达数十极） \\
典型应用 & 汽轮发电机、核电机组 & 水轮发电机、柴油发电机 \\
分析方法 & 隐极机相量图（同步电抗 $X_t$） & 双反应理论（$X_d$、$X_q$） \\
\bottomrule
\end{tabular}
\end{center}

\section{同步电机的运行状态}

同步电机的运行状态由\textbf{功率角}（功角）$\delta$ 的符号与大小决定：$\delta$ 是励磁电动势 $\dot E_0$ 与端电压 $\dot U$ 之间的夹角。

\subsection{发电机状态}
转子由原动机拖动，转速为同步转速 $n_1$，$\dot E_0$ 超前 $\dot U$ 一个正角 $\delta$，电机向电网输出有功功率。这是同步电机最主要的运行状态。

\subsection{电动机状态}
转子拖着负载旋转，从电网吸收有功功率，$\dot E_0$ 滞后 $\dot U$（$\delta<0$），电磁转矩为驱动转矩。同步电动机的转速严格等于 $n_1$，与负载无关（除非失步），故用于要求恒速的场合（如大型压缩机、磨机）。

\subsection{调相机（补偿机）状态}
电机轴上不带机械负载，专门向电网输出或吸收\textbf{无功功率}以改善功率因数。过励时输出无功（相当于电容器），欠励时吸收无功（相当于电抗器）。这是同步电机"第三种用途"。

\begin{center}
\begin{tabular}{@{}llll@{}}
\toprule
运行状态 & 有功流向 & 功角 $\delta$ & 主要用途 \\
\midrule
发电机 & 电机 $\rightarrow$ 电网 & $\delta>0$ & 发电（主要用途） \\
电动机 & 电网 $\rightarrow$ 电机 & $\delta<0$ & 恒速拖动 \\
调相机 & 只交换无功 & $\delta\approx0$ & 改善功率因数、调压 \\
\bottomrule
\end{tabular}
\end{center}

\section{同步电机的励磁方式}

励磁系统为转子提供直流励磁电流，直接决定电机的无功功率与端电压，是同步电机最重要的辅助系统之一。

\begin{enumerate}
  \item \textbf{直流励磁机励磁}：用同轴的直流发电机供电。结构简单、独立性强，但容量受换向限制、维护量大，已逐渐淘汰。
  \item \textbf{静止半导体励磁}：由励磁变压器从机端取电，经可控整流后供给转子。响应快、维护简单，是大中型机组的主流。
  \item \textbf{旋转半导体（无刷）励磁}：交流励磁机与主发电机同轴，其电枢（旋转）输出的交流经装在转轴上的旋转整流器整流后直接供主发电机励磁，无需滑环电刷。可靠性高，用于大型机组。
  \item \textbf{三次谐波励磁}：利用气隙磁场中的三次谐波在附加绕组中感应电动势，经整流后励磁，具有自励恒压特性，多用于小型机组。
\end{enumerate}

\section{同步电机的型号和额定值}

\subsection{额定容量 $S_N$ 与额定功率 $P_N$}
同步发电机铭牌上给出\textbf{额定视在功率} $S_N$（kV·A）和\textbf{额定有功功率} $P_N$（kW），二者关系为
\begin{equation}
P_N=S_N\cos\varphi_N .
\end{equation}
对于同步电动机，额定功率指轴上输出的机械功率；对于调相机，只给额定视在功率。

\subsection{额定电压 $U_N$}
额定运行时定子绕组的\textbf{线电压}，单位 V 或 kV。

\subsection{额定电流 $I_N$}
额定运行时定子绕组的\textbf{线电流}，单位 A。由
\begin{equation}
S_N=\sqrt3U_NI_N
\end{equation}
确定。

\subsection{额定功率因数 $\cos\varphi_N$}
额定运行时的功率因数，一般发电机为 0.8（滞后），电动机为 0.8～0.9（超前或滞后按铭牌）。

\subsection{额定频率 $f_N$}
我国为 50\,Hz。极对数 $p$ 与转速严格对应：$n_N=60f_N/p$，故同步电机转速是"分档"的（3000、1500、1000、750、600…r/min）。

\subsection{额定励磁电流 $I_{fN}$ 与额定励磁电压 $U_{fN}$}
额定运行（额定电压、额定电流、额定功率因数）时的转子励磁电流与电压。

\keyformula{S_N=\sqrt3U_NI_N,\qquad P_N=S_N\cos\varphi_N,\qquad n_N=\frac{60f_N}{p}}

\corepoint{同步电机额定值有两处与感应电机不同：① 发电机铭牌给的是\textbf{视在功率} $S_N$（因为要同时保证电压和功率因数），电动机才给轴上机械功率 $P_N$；② 转速不是"额定"出来的，而是由极对数和频率\textbf{唯一确定}（$n=f/p$ 的整数分档），所以同步电机没有"额定转差率"这一说。}

% ============================================================
% 第 17 章 同步发电机的基本理论（教材第五篇，印刷页 183--203）
% 节名：第一节 同步发电机的运行方式；第二节 隐极同步发电机的分析方法；
%       第三节 凸极同步发电机的分析方法；第四节 同步发电机功率和转矩方程；第五节 同步发电机的运行特性和电抗测定
% ============================================================
\chapter{同步发电机的基本理论}

同步发电机与变压器、感应电机的根本区别在于：它的主磁场由\textbf{转子上的直流励磁绕组}独立建立，而定子（电枢）三相电流产生的磁动势会反过来影响主磁场——这种现象称为\textbf{电枢反应}。整章的逻辑是：先确定电枢反应的性质（去磁、助磁还是交磁），再把电枢反应"折算"成电抗，于是就能像变压器那样写出电压方程、画相量图、算功率。

\section{同步发电机的运行方式}

同步发电机转子由原动机拖动，以同步转速 $n_1=60f/p$ 旋转。转子励磁绕组通入直流电流 $I_f$，建立\textbf{励磁磁动势} $F_f$ 与主磁场；主磁场切割定子三相对称绕组，感应三相对称电动势。定子接上负载后，三相电流产生\textbf{电枢磁动势} $F_a$，它与 $F_f$ 共同建立气隙磁场。

稳态对称运行时，只研究一相（通常取 $A$ 相），并用相量表示。各量正方向按发电机惯例规定：电枢电流 $\dot I$ 的正方向与感应电动势 $\dot E$ 相同（由中性点流向出线端），即电机\textbf{输出}电流为正。

\section{隐极同步发电机的分析方法}

\subsection{电枢反应的性质}
电枢磁动势 $F_a$ 对主磁场的作用称为\textbf{电枢反应}，其性质取决于负载的性质（即 $\dot I$ 与 $\dot E_0$ 的相位关系）。设 $\psi$ 为 $\dot E_0$ 与 $\dot I$ 的夹角（内功率因数角）：
\begin{itemize}
  \item $\psi=0$（纯电阻负载，$\dot I$ 与 $\dot E_0$ 同相）：$F_a$ 的轴线与主磁极轴线相差 $90^\circ$ 电角度，称为\textbf{交轴（交磁）电枢反应}，使气隙磁场轴线扭斜，磁场半边增强、半边削弱，不改变总磁通量；
  \item $\psi=90^\circ$（纯电感负载，$\dot I$ 滞后 $\dot E_0$ $90^\circ$）：$F_a$ 的轴线与主磁极轴线\textbf{相反}，称为\textbf{直轴去磁电枢反应}，使主磁通削弱、端电压下降（所以感性负载需要更大的励磁电流）；
  \item $\psi=-90^\circ$（纯电容负载，$\dot I$ 超前 $\dot E_0$ $90^\circ$）：$F_a$ 轴线与主磁极轴线\textbf{相同}，称为\textbf{直轴助磁电枢反应}，使主磁通增强、端电压升高。
\end{itemize}

\begin{center}
\begin{tabular}{@{}llll@{}}
\toprule
负载性质 & $\psi$ & 电枢反应性质 & 对主磁通/端电压的影响 \\
\midrule
纯电阻 & $0$ & 交轴（交磁） & 磁场扭斜，总量基本不变 \\
纯电感（感性） & $90^\circ$ & 直轴去磁 & 削弱主磁通，端电压下降 \\
纯电容（容性） & $-90^\circ$ & 直轴助磁 & 增强主磁通，端电压升高 \\
\bottomrule
\end{tabular}
\end{center}

\subsection{电枢反应电抗与同步电抗}
把电枢反应磁动势 $F_a$ 的作用等效为一个电抗压降：$F_a$ 产生的磁通在电枢绕组中感应电动势 $\dot E_a$，且 $\dot E_a=-\mathrm{j}\dot I X_a$（$X_a$ 为\textbf{电枢反应电抗}）。再加上漏磁通产生的漏抗压降 $-\mathrm{j}\dot I X_\sigma$，两者合并为
\begin{equation}
X_t=X_a+X_\sigma ,
\end{equation}
称为\textbf{同步电抗}。

\keyformula{X_t=X_a+X_\sigma}

$X_t$ 综合反映了电枢反应与漏磁两种效应。\textbf{电枢反应是同步电机区别于其他电机的核心}：它把"磁场互相作用"这个物理过程，压缩成一个电抗压降，从而使同步电机可以用电路方法来分析。这是"路的分析方法"最漂亮的一次应用。

\subsection{电压方程与相量图}
计及电枢电阻 $R_a$，隐极同步发电机的电压方程为
\begin{equation}
\dot E_0=\dot U+\dot I\left(R_a+\mathrm{j}X_t\right).
\end{equation}

\keyformula{\dot E_0=\dot U+\dot I\left(R_a+\mathrm{j}X_t\right)}

通常 $R_a\ll X_t$，可忽略，得
\begin{equation}
\dot E_0\approx\dot U+\mathrm{j}\dot I X_t .
\end{equation}

\subsubsection{相量图画法}
以端电压 $\dot U$ 为参考，按负载功率因数角 $\varphi$ 画出 $\dot I$；在 $\dot U$ 上加 $\dot I R_a$（平行 $\dot I$）与 $\mathrm{j}\dot I X_t$（超前 $\dot I$ $90^\circ$），其相量和即 $\dot E_0$；$\dot E_0$ 与 $\dot U$ 的夹角即\textbf{功率角}（功角）$\delta$，$\dot E_0$ 与 $\dot I$ 的夹角为 $\psi$，且有
\begin{equation}
\psi=\delta+\varphi .
\end{equation}

\section{凸极同步发电机的分析方法}

凸极电机气隙不均匀（极下小、极间大），同样大小的电枢磁动势作用在不同位置时产生的磁通不同，因此不能用一个同步电抗描述。工程上采用\textbf{双反应理论}：把电枢磁动势分解为两个分量，分别沿\textbf{直轴}（磁极轴线，$d$ 轴）与\textbf{交轴}（极间轴线，$q$ 轴），各自用相应的电抗计算。

\subsection{双反应理论}
把电枢电流分解为
\begin{equation}
\dot I=\dot I_d+\dot I_q,\qquad
I_d=I\sin\psi,\qquad I_q=I\cos\psi ,
\end{equation}
其中 $\dot I_d$ 产生直轴电枢反应（$X_{ad}$），$\dot I_q$ 产生交轴电枢反应（$X_{aq}$）。相应的同步电抗为
\begin{equation}
X_d=X_{ad}+X_\sigma,\qquad X_q=X_{aq}+X_\sigma .
\end{equation}
对凸极机通常 $X_d>X_q$（直轴气隙小、磁阻小，故电抗大）。

\subsection{电压方程}
\begin{equation}
\dot E_0=\dot U+\dot I R_a+\mathrm{j}\dot I_dX_d+\mathrm{j}\dot I_qX_q .
\end{equation}

\keyformula{\dot E_0=\dot U+\dot I R_a+\mathrm{j}\dot I_dX_d+\mathrm{j}\dot I_qX_q}

忽略 $R_a$ 时
\begin{equation}
\dot E_0\approx\dot U+\mathrm{j}\dot I_dX_d+\mathrm{j}\dot I_qX_q .
\end{equation}

\subsubsection{相量图的特点}
凸极机的相量图中，$\dot E_0$ 的位置不能像隐极机那样直接由 $\dot U+\mathrm{j}\dot I X_t$ 一次求出，因为 $X_d\ne X_q$、$I_d$、$I_q$ 又依赖 $\psi$。实用的作图法是：先由 $\dot U$ 与 $\dot I$ 求出 $\dot E_Q=\dot U+\dot I R_a+\mathrm{j}\dot I X_q$，其方向即 $q$ 轴方向（$\dot E_0$ 的方向），据此分解出 $I_d$、$I_q$，再求 $\dot E_0$。这一"辅助电动势 $\dot E_Q$"的技巧是解凸极机题目的关键。

\begin{center}
\begin{tabular}{@{}p{2.2cm}p{3.4cm}p{3.6cm}p{3.2cm}@{}}
\toprule
项目 & 隐极机 & 凸极机 & 说明 \\
\midrule
气隙 & 均匀 & 不均匀 & 决定能否用单一电抗 \\
分析方法 & 同步电抗 $X_t$ & 双反应理论（$X_d$、$X_q$） & 分解 $I_d$、$I_q$ \\
电压方程 & $\dot E_0=\dot U+\dot I(R_a+\mathrm{j}X_t)$ & $\dot E_0=\dot U+\dot IR_a+\mathrm{j}\dot I_dX_d+\mathrm{j}\dot I_qX_q$ & — \\
电抗关系 & $X_t=X_a+X_\sigma$ & $X_d=X_{ad}+X_\sigma$，$X_q=X_{aq}+X_\sigma$ & 通常 $X_d>X_q$ \\
相量图 & 一次求出 $\dot E_0$ & 需先求 $\dot E_Q$ 定 $q$ 轴 & 凸极机作图较繁 \\
\bottomrule
\end{tabular}
\end{center}

\corepoint{隐极与凸极的分界只有一句话：\textbf{气隙是否均匀}。均匀 $\Rightarrow$ 一个 $X_t$ 够了；不均匀 $\Rightarrow$ 必须分解 $d$、$q$ 两轴。考试若给了 $X_d$、$X_q$ 两个电抗，一定是凸极机；只给一个 $X_t$（或 $X_s$），一定是隐极机——先看参数再动笔，可省大量时间。}

\section{同步发电机功率和转矩方程}

\subsection{电磁功率}
发电机输出的有功功率为
\begin{equation}
P_2=mUI\cos\varphi ,
\end{equation}
加上电枢铜耗 $p_{Cu}=mI^2R_a$ 得电磁功率 $P_{em}=P_2+p_{Cu}$。忽略 $R_a$ 时 $P_{em}\approx P_2$。

\subsubsection{隐极机的功角特性}
由相量图的几何关系可得
\begin{equation}
P_{em}=\frac{mE_0U}{X_t}\sin\delta .
\end{equation}

\keyformula{P_{em}=\frac{mE_0U}{X_t}\sin\delta}

\subsubsection{凸极机的功角特性}
\begin{equation}
P_{em}=\frac{mE_0U}{X_d}\sin\delta+\frac{mU^{2}}{2}\left(\frac{1}{X_q}-\frac{1}{X_d}\right)\sin2\delta .
\end{equation}

\keyformula{P_{em}=\frac{mE_0U}{X_d}\sin\delta+\frac{mU^{2}}{2}\left(\frac1{X_q}-\frac1{X_d}\right)\sin2\delta}

两项的物理意义：
\begin{itemize}
  \item 第一项（$\sin\delta$）是\textbf{励磁电磁功率}，由励磁磁动势与电枢磁动势的相互作用产生，与 $E_0$（即励磁电流）成正比，是发电机传递有功的主要成分；
  \item 第二项（$\sin2\delta$）是\textbf{磁阻功率}（凸极效应功率），由 $d$、$q$ 轴磁阻不等产生，与励磁无关，即使 $E_0=0$ 也存在——这是\textbf{磁阻同步电动机}（无励磁转子）能工作的原理。
\end{itemize}

\subsection{转矩方程}
\begin{equation}
T=\frac{P_{em}}{\Omega_1},\qquad \Omega_1=\frac{2\pi n_1}{60}=\frac{2\pi f}{p}.
\end{equation}
稳态时 $T=T_1=T_2+T_0$（原动机拖动转矩 $=$ 电磁转矩 $=$ 输出转矩 $+$ 空载制动转矩）。

\section{同步发电机的运行特性和电抗测定}

\subsection{空载特性}
$n=n_1$、$I=0$ 时，端电压（空载电动势）$E_0$ 与励磁电流 $I_f$ 的关系 $E_0=f(I_f)$。由于铁芯饱和，曲线起始段近似直线（气隙线），随后弯曲并趋于饱和。空载特性本质上就是电机的磁化曲线。

\subsection{短路特性}
$n=n_1$、三相稳态短路（$U=0$）时，短路电流 $I_k$ 与励磁电流 $I_f$ 的关系 $I_k=f(I_f)$。由于短路时电枢反应为纯去磁、磁路不饱和，故\textbf{短路特性是一条直线}。

\subsection{外特性与调整特性}
\begin{itemize}
  \item \textbf{外特性}：$n=n_1$、$I_f=$ 常数、$\cos\varphi=$ 常数时，$U=f(I)$。感性负载时随 $I$ 增大 $U$ 下降；容性负载时可能上升。
  \item \textbf{调整特性}：$n=n_1$、$U=U_N$、$\cos\varphi=$ 常数时，$I_f=f(I)$。为维持端电压不变，感性负载需增加励磁。
\end{itemize}

\subsection{同步电抗与短路比}
由空载特性与短路特性可求同步电抗：在短路特性上取额定电流对应的励磁电流 $I_{fk}$，再在空载特性上查同一励磁电流下的空载电动势 $E_0$，则
\begin{equation}
X_d=\frac{E_0}{I_k}\quad\left(I_k=I_N\right).
\end{equation}
\textbf{短路比}定义为空载时产生额定电压所需的励磁电流 $I_{f0}$ 与短路时产生额定电流所需的励磁电流 $I_{fk}$ 之比：
\begin{equation}
K_c=\frac{I_{f0}}{I_{fk}}\approx\frac{1}{X_d^{*}} .
\end{equation}
短路比大 $\Rightarrow$ $X_d$ 小 $\Rightarrow$ 电压调整率小、并列运行稳定性好，但气隙大、材料多、造价高。汽轮发电机 $K_c\approx0.5\sim0.7$，水轮发电机 $K_c\approx1.0\sim1.4$。

\begin{center}
\begin{tabular}{@{}p{2.2cm}p{3.6cm}p{4.2cm}p{3.0cm}@{}}
\toprule
特性/参数 & 条件 & 形状或关系 & 用途 \\
\midrule
空载特性 & $I=0$，$n=n_1$ & $E_0=f(I_f)$，饱和曲线 & 求 $X_d$、判断饱和程度 \\
短路特性 & $U=0$，$n=n_1$ & $I_k=f(I_f)$，直线 & 求 $X_d$ 与短路比 \\
外特性 & $I_f$、$\cos\varphi$ 恒定 & $U=f(I)$，感性下降 & 电压调整 \\
调整特性 & $U=U_N$、$\cos\varphi$ 恒定 & $I_f=f(I)$ & 励磁调节依据 \\
短路比 & — & $K_c\approx1/X_d^{*}$ & 稳定性与经济性权衡 \\
\bottomrule
\end{tabular}
\end{center}

\corepoint{本章最容易混的两组概念：① \textbf{$X_t$ 与 $X_d$}——隐极机用 $X_t$，凸极机用 $X_d$、$X_q$，两者不能互换代入；② \textbf{空载电动势 $E_0$ 与端电压 $U$}——$E_0$ 是励磁决定的"内部电动势"，$U$ 是带负载后的端电压，两者之差就是电枢反应与漏抗压降。所有同步电机的计算题都围绕 $\dot E_0=\dot U+\mathrm{j}\dot I X$（或其凸极形式）与 $P_{em}=\frac{mE_0U}{X}\sin\delta$ 这两组式子转。}

% ============================================================
% 第 18 章 同步发电机的并联运行（教材第五篇，印刷页 204--214）
% 节名：第一节 同步发电机的并网条件和方法；第二节 同步发电机电磁功率与功角特性；
%       第三节 同步发电机并联运行时的功率调节
% ============================================================
\chapter{同步发电机的并联运行}

电力系统中的同步发电机都是并联运行的。并网要解决三个问题：\textbf{怎么安全地并上去}（并网条件与方法）、\textbf{并上以后能承担多大负载}（功角特性与静态稳定）、\textbf{怎么分配有功与无功}（功率调节）。这一章是同步电机理论中最"工程"的一章，也是电力系统课程的接口。

\section{同步发电机的并网条件和方法的}

\subsection{理想并网条件}
把同步发电机投入电网并联运行，必须满足下列条件（否则并网瞬间会产生巨大的冲击电流与转矩，甚至损坏机组）：
\begin{enumerate}
  \item 发电机的\textbf{相序}与电网相序相同；
  \item 发电机的\textbf{频率}与电网频率相同；
  \item 发电机的\textbf{端电压大小}与电网电压相同；
  \item 发电机端电压的\textbf{相位}与电网电压相位一致（即电压相量重合）。
\end{enumerate}
条件 1—4 全部满足时，发电机与电网之间没有环流，合闸瞬间冲击最小，称为\textbf{理想并网}。

\begin{center}
\begin{tabular}{@{}p{2.2cm}p{4.4cm}p{5.2cm}@{}}
\toprule
条件 & 不满足时的后果 & 调整手段 \\
\midrule
相序不同 & 相当于两相短路，冲击电流极大 & 并网前核对接线（唯一必须先满足的） \\
频率不同 & 产生拍频，功率振荡、电流冲击 & 调原动机转速 \\
电压大小不等 & 合闸瞬间产生无功冲击电流 & 调励磁电流 \\
相位不一致 & 产生很大的有功冲击与转矩冲击 & 调相位（准同期最关键的"合闸时机"） \\
\bottomrule
\end{tabular}
\end{center}

\subsection{准同期法（准同步法）}
发电机先由原动机拖到接近同步转速，加励磁使端电压接近电网电压，然后\textbf{在相位一致的瞬间合闸}。
\begin{itemize}
  \item \textbf{灯光熄灭法}：把三组灯接在发电机与电网的同名相之间，灯全暗时表示两侧电压同相位，此时合闸。
  \item \textbf{灯光旋转法}：把一组灯接在一侧的同名相与另一侧的异名相之间，灯依次明暗旋转，表示频率有微小差；当直接接同相的那组灯熄灭、另两组灯亮度相同时合闸。
  \item 现代机组普遍采用\textbf{自动准同期装置}，自动调节频率与电压并捕捉合闸时机。
\end{itemize}
准同期法的优点是合闸冲击小，缺点是操作要求高、并网较慢。

\subsection{自同期法}
发电机不加励磁、转速接近同步转速时先合闸，随后立即投入励磁，靠定子磁场把转子\textbf{牵入同步}。优点：操作简单、并网迅速，可在系统故障时快速投入；缺点：合闸瞬间从电网吸取很大的励磁涌流（可达额定电流数倍）、引起电网电压短时下降。适用于事故备用机组的紧急并网。

\corepoint{并网条件的"代价排序"要记住：\textbf{相序错}是灾难（相当于短路），必须事先核对；\textbf{相位差}直接决定冲击大小（$\Delta U=2U\sin(\Delta\delta/2)$，差 $30^\circ$ 就有 $0.52U$ 的电压差），所以准同期法把"捕捉相位"当成核心；频率差与电压差相对好调整。考试若问"并网时为什么先调频后调压"，答：调频决定相位漂移速度，必须先把频差压到很小，才可能捕捉到相位一致的瞬间。}

\section{同步发电机电磁功率与功角特性}

第 17 章已得隐极机的功角特性
\begin{equation}
P_{em}=\frac{mE_0U}{X_t}\sin\delta .
\end{equation}
这一节讨论它的工程含义。

\subsection{功角的物理意义}
$\delta$ 是励磁电动势 $\dot E_0$ 与端电压 $\dot U$ 之间的相位差，对应着\textbf{转子磁极轴线}与\textbf{气隙合成磁场轴线}之间的空间夹角。发电机输出有功功率越大，转子磁极必须越"超前"于气隙磁场——就像两人拉一根弹性绳，拉的力越大，绳子的伸长（角度差）越大。所以 $\delta$ 本质上是"电磁弹簧"的变形量。

\subsection{静态稳定与整步功率}
当 $\delta<90^\circ$ 时，$\dfrac{\mathrm{d}P_{em}}{\mathrm{d}\delta}>0$：负载增加时电机能自动增加电磁功率以平衡，运行是\textbf{稳定}的。
当 $\delta>90^\circ$ 时，$\dfrac{\mathrm{d}P_{em}}{\mathrm{d}\delta}<0$：负载再增加，电磁功率反而减小，转子继续加速、$\delta$ 继续增大——电机\textbf{失步}（失去同步）。

因此 $\delta=90^\circ$ 是隐极机的\textbf{静态稳定极限}，此时
\begin{equation}
P_{\max}=\frac{mE_0U}{X_t}
\end{equation}
称为\textbf{功率极限}。整步功率系数（比整步功率）定义为
\begin{equation}
P_{sy}=\frac{\mathrm{d}P_{em}}{\mathrm{d}\delta}=\frac{mE_0U}{X_t}\cos\delta ,
\end{equation}
它表征电机维持同步的能力：$\delta$ 越大 $P_{sy}$ 越小，越接近失步。

\keyformula{P_{\max}=\frac{mE_0U}{X_t},\qquad P_{sy}=\frac{mE_0U}{X_t}\cos\delta}

\subsection{过载能力}
定义过载能力（静态过载倍数）
\begin{equation}
k_m=\frac{P_{\max}}{P_N}=\frac{1}{X_t^{*}\sin\delta_N},
\end{equation}
一般要求 $k_m=1.7\sim2.0$，即额定运行时的功角 $\delta_N\approx25^\circ\sim35^\circ$，留出足够的稳定裕量。

\subsection{提高静态稳定的措施}
由 $P_{\max}=mE_0U/X_t$ 可见：
\begin{itemize}
  \item \textbf{增大励磁}（提高 $E_0$）可提高功率极限——这是励磁调节（自动电压调节器 AVR）稳定电网的基本手段；
  \item \textbf{减小同步电抗}（设计上增大气隙、增大短路比）可提高功率极限，但造价上升；
  \item 采用\textbf{快速励磁}、\textbf{电力系统稳定器（PSS）}可进一步改善动态稳定。
\end{itemize}

\section{同步发电机并联运行时的功率调节}

并网后的发电机与电网"锁"在一起，转速由电网频率决定，不能随意改变。此时调节的是\textbf{有功}与\textbf{无功}两个通道，且两者互相独立（这是同步电机最优雅的性质之一）。

\subsection{有功功率的调节——调原动机}
增大原动机（汽轮机/水轮机）的进汽量或进水量，输入机械功率增加，转子加速、$\delta$ 增大，电磁功率随之增大，直到与输入功率重新平衡。故
\begin{equation}
P_{em}\uparrow\ \Longleftrightarrow\ \delta\uparrow\ \Longleftrightarrow\ \text{原动机输入转矩}\uparrow .
\end{equation}
\textbf{有功功率只由原动机的输入功率决定，与励磁电流无关}。这就是电力系统"调频靠调有功"（调机组出力）的依据。

\subsection{无功功率的调节——调励磁}
在保持原动机输入不变（有功不变）的条件下增大励磁电流 $I_f$：$E_0$ 增大，功率因数角改变，电枢电流的无功分量增大，发电机向电网输出\textbf{感性无功}（过励状态）；反之减小励磁，发电机吸收感性无功（欠励状态）。

\textbf{无功功率只由励磁电流决定，与输入机械功率基本无关}。这是电力系统"调压靠调无功"（调励磁）的依据。

\begin{center}
\begin{tabular}{@{}p{2.4cm}p{3.4cm}p{3.8cm}p{3.2cm}@{}}
\toprule
调节量 & 调节对象 & 主要变化 & 电网侧效果 \\
\midrule
有功功率 & 原动机进汽/进水量 & $\delta$ 变化，$E_0$ 基本不变 & 调频（有功平衡） \\
无功功率 & 励磁电流 $I_f$ & $E_0$ 变化，$\delta$ 基本不变 & 调压（无功平衡） \\
\bottomrule
\end{tabular}
\end{center}

\subsection{并联运行时的无功分配}
多台发电机并联时，无功功率按各机的\textbf{励磁状态}分配，而非按容量自动分配：励磁电流偏大的机组承担更多无功。若励磁分配不当，会出现"一台过励、一台欠励"的环流，故需配置无功分配装置（无功电流补偿）。

\subsection{调相运行}
若把发电机轴上负载去掉、只带励磁运行，则电机成为\textbf{调相机}：过励时向电网输出感性无功（相当于并联电容器），欠励时吸收感性无功（相当于并联电抗器）。调相机用于长距离输电线路的无功补偿与电压支撑。

\corepoint{本章的两条"独立性"是全章的灵魂：\textbf{有功只跟原动机走，无功只跟励磁走}。由此可推出所有定性结论：增励磁 $\Rightarrow$ 无功增大、功率因数下降（过励）、端电压升高、功角略减；增原动机输入 $\Rightarrow$ 有功增大、$\delta$ 增大、无功基本不变。凡是"调某个量会怎样"的问题，先判断它属于有功通道还是无功通道，再作答，几乎不会错。}

% ============================================================
% 第 19 章 同步电动机和同步补偿机（教材第五篇，印刷页 215--225）
% 节名：第一节 同步电动机电压方程和相量图；第二节 同步电动机的功角特性与 V 形曲线；
%       第三节 同步电动机的启动与调速；第四节 同步补偿机；第五节 永磁同步电动机
% ============================================================
\chapter{同步电动机和同步补偿机}

同步电机是可逆的：同一台机器，输入机械功率、输出电功率时为发电机；输入电功率、输出机械功率时为电动机。本章把第 17 章的方程改写成\textbf{电动机惯例}，重点解决两个新问题：\textbf{怎么起动}（同步电动机没有起动转矩）和\textbf{V 形曲线}（电流随励磁变化）。

\section{同步电动机电压方程和相量图}

\subsection{正方向的改变}
改作电动机运行时，电枢电流的正方向改为\textbf{从电网流入电机}，与感应电动势 $\dot E_0$ 反方向（吸收电功率），即
\begin{equation}
\dot U=\dot E_0+\dot I\left(R_a+\mathrm{j}X_t\right),
\end{equation}
或写成
\begin{equation}
\dot E_0=\dot U-\dot I\left(R_a+\mathrm{j}X_t\right).
\end{equation}

\keyformula{\dot U=\dot E_0+\dot I\left(R_a+\mathrm{j}X_t\right)\quad\text{（电动机惯例）}}

与发电机方程相比，只是 $\dot I$ 的符号变了。这意味着：\textbf{同一台电机，发电机与电动机的方程是同一个，区别只在电流方向与功率流向}。

\subsection{相量图与功率角}
电动机的功率角 $\delta$ 为 $\dot E_0$ 滞后 $\dot U$ 的角度（取负值或按绝对值定义），表示转子磁极轴线\textbf{滞后}于气隙磁场轴线。物理图像：电动机拖着负载，转子"落后"于旋转磁场，靠电磁"弹簧"拉着走。

电枢电流与端电压的夹角即功率因数角 $\varphi$；$\dot E_0$ 与 $\dot I$ 的夹角为 $\psi$。电动机可以工作在\textbf{超前}功率因数（过励）状态——这是同步电动机相对感应电动机的最大优点：可以补偿电网的无功。

\section{同步电动机的功角特性与 V 形曲线}

\subsection{功角特性}
电动机的电磁功率（由电网输入、转换为机械功率的部分）为
\begin{equation}
P_{em}=\frac{mE_0U}{X_t}\sin\delta ,
\end{equation}
凸极机为
\begin{equation}
P_{em}=\frac{mE_0U}{X_d}\sin\delta+\frac{mU^{2}}{2}\left(\frac{1}{X_q}-\frac{1}{X_d}\right)\sin2\delta .
\end{equation}
形式上与发电机完全相同，只是 $\delta$ 取相反符号（功率流向相反）。最大电磁功率（过载能力）同样出现在 $\delta=90^\circ$ 附近。

\subsection{V 形曲线}
在\textbf{输出功率恒定}（即输入有功不变）的条件下，改变励磁电流 $I_f$，电枢电流 $I$ 随之变化，其曲线形如字母"V"，故称 \textbf{V 形曲线}。

\subsubsection{形成机理}
输出功率不变意味着
\begin{equation}
P=\sqrt3UI\cos\varphi=\text{常数}\quad\Longrightarrow\quad I\cos\varphi=\text{常数},
\end{equation}
即电枢电流的\textbf{有功分量不变}。而励磁改变的是 $E_0$，从而改变\textbf{无功分量} $I\sin\varphi$。于是
\begin{equation}
I=\sqrt{\left(I\cos\varphi\right)^{2}+\left(I\sin\varphi\right)^{2}},
\end{equation}
当无功分量为零（$\cos\varphi=1$）时 $I$ 最小，即为 V 形曲线的\textbf{最低点}（正常励磁）；励磁大于该值时过励（电流超前、输出容性无功），小于该值时欠励（电流滞后、吸收感性无功）。

\subsubsection{不同负载下的 V 形曲线}
对应不同的输出功率可画出一族 V 形曲线：负载越大，曲线整体上移、最低点电流越大；各曲线最低点的连线称为\textbf{临界励磁线}（$\cos\varphi=1$ 线），其右侧为过励区、左侧为欠励区。欠励过度时，电机会因励磁不足而\textbf{失步}，故存在\textbf{稳定极限}（对应 $\delta=90^\circ$），V 形曲线左侧的部分区域是不稳定区。

\begin{center}
\begin{tabular}{@{}p{2.2cm}p{3.0cm}p{4.0cm}p{3.4cm}@{}}
\toprule
励磁状态 & 电枢电流性质 & 对电网的作用 & 典型用途 \\
\midrule
正常励磁 & $\cos\varphi=1$，$I$ 最小 & 不交换无功 & V 形曲线最低点 \\
过励 & 电流超前（容性） & 向电网输出感性无功 & 改善功率因数（进相运行） \\
欠励 & 电流滞后（感性） & 从电网吸收感性无功 & 一般不采用，易失步 \\
\bottomrule
\end{tabular}
\end{center}

\corepoint{V 形曲线的关键前提是"\textbf{有功不变}"。若输出功率变了，曲线整体平移，最低点位置也变——所以问"增大励磁电流，电流如何变化"时，必须先确认"输出功率是否不变"。另外记住：\textbf{同步电动机过励时功率因数超前}，这正是它能当"无功补偿装置"用的原因，也是它与感应电动机最本质的差别。}

\subsection{同步电动机的优缺点}
\begin{itemize}
  \item \textbf{优点}：转速严格恒定（$n=n_1$，与负载无关）；功率因数可调，可过励运行补偿电网无功；过载能力强（气隙大）、效率高。
  \item \textbf{缺点}：需要直流励磁电源；\textbf{没有起动转矩}，必须采用专门的起动方法；结构复杂、造价高。
\end{itemize}

\section{同步电动机的启动与调速}

\subsection{为什么不能自起动}
稳态运行时转子以同步转速旋转，转子磁极与气隙磁场相对静止。若在转子静止时投入三相电源，气隙旋转磁场以 $n_1$ 旋转，转子受到的电磁转矩在一个周期内\textbf{平均值为零}（时而正、时而负），转子无法加速，故同步电动机不能自起动。

\subsection{起动方法}
\begin{enumerate}
  \item \textbf{异步起动法}（最常用）：在转子上加装\textbf{笼型起动绕组}（阻尼绕组）。起动时先不加励磁，定子接三相电源，电机作为笼型感应电动机起动加速；当转速接近同步转速（约 95\% $n_1$）时投入励磁，靠同步转矩把转子\textbf{牵入同步}。牵入同步的瞬间会产生振荡，阻尼绕组可抑制振荡。
  \item \textbf{辅助电动机起动法}：用一台小电动机把同步电动机拖到接近同步转速，再投入电网与励磁。设备多、操作繁，仅用于特殊场合。
  \item \textbf{变频起动法}：用变频装置从低频开始供电，使旋转磁场转速从零缓慢上升，转子始终同步跟随。起动平滑、无冲击，是现代大容量同步电动机（如轧机、磨机）的主流方案。
\end{enumerate}

\subsection{调速}
同步电动机的转速 $n=60f/p$，只能通过\textbf{变频}（改变 $f$）或\textbf{变极}（改变 $p$）调速，且为有级（变极）或无级（变频）。变频调速时同样要保持 $U/f$ 恒定以维持磁通。这比感应电动机的调速手段少（因为没有转差率可调），但转速精度最高。

\section{同步补偿机}

同步补偿机（调相机）是\textbf{轴上不带任何机械负载}的同步电机，专门用来向电网输出或吸收无功功率。

\begin{itemize}
  \item \textbf{过励运行}：向电网输出感性无功，相当于并联电容器，可提高电网功率因数、支撑受端电压；
  \item \textbf{欠励运行}：从电网吸收感性无功，相当于并联电抗器，可抑制长线路空载或轻载时的工频过电压。
\end{itemize}
调相机的优点是无功调节连续、平滑、容量大；缺点是旋转设备有损耗、维护量大、造价高。近年来大量被\textbf{静止无功补偿器（SVC）}与\textbf{静止同步补偿器（STATCOM）}取代，但在特高压长距离输电中仍有应用。

\section{永磁同步电动机}

用\textbf{永磁体}代替励磁绕组建立主磁场，取消励磁电源、滑环与电刷，是同步电机的重要发展方向。

\subsubsection{结构特点}
按永磁体在转子上的位置，分为\textbf{表贴式}（永磁体贴在转子表面，$L_d\approx L_q$）与\textbf{内置式}（永磁体嵌入转子内部，$L_d<L_q$，存在磁阻转矩）。永磁材料常用钕铁硼（NdFeB）、钐钴等。

\subsubsection{优点}
\begin{itemize}
  \item 无励磁损耗，效率高（比同容量感应电动机高 2\%～8\%）；
  \item 功率密度大、体积小、质量轻；
  \item 无电刷滑环，可靠性高、维护简单；
  \item 配合变频器可实现宽范围调速与精确的位置/转矩控制。
\end{itemize}

\subsubsection{问题}
\begin{itemize}
  \item 永磁体价格高，且高温下可能\textbf{失磁}（需控制温度）；
  \item 无法通过调节励磁来调压、调无功，弱磁调速范围受限；
  \item 停机检修时反电动势仍然存在，安全上需注意。
\end{itemize}
永磁同步电动机广泛用于新能源汽车驱动、伺服系统、电梯、压缩机、风电（直驱永磁）等领域，是当前电机技术最活跃的方向之一。

\begin{center}
\begin{tabular}{@{}p{2.2cm}p{3.6cm}p{3.8cm}p{3.4cm}@{}}
\toprule
项目 & 同步电动机 & 感应电动机 & 直流电动机 \\
\midrule
转速 & 严格 $n_1$，与负载无关 & 随负载略降（有转差） & 随负载下降较多 \\
功率因数 & 可调（可超前） & 滞后，轻载时很低 & 较高（有换向器） \\
起动 & 需异步/变频起动 & 可自起动 & 可自起动，需限流 \\
调速 & 变频/变极 & 变频/变极/变转差 & 调压/调磁，最方便 \\
维护 & 需励磁（永磁机免） & 笼型免维护 & 换向器电刷需维护 \\
\bottomrule
\end{tabular}
\end{center}

% ============================================================
% 第 20 章 直流电机概述（教材第六篇，印刷页 227--234）
% 节名：第一节 直流电机的工作原理；第二节 直流电机的结构；第三节 直流电机的额定值
% ============================================================
\chapter{直流电机概述}

直流电机是历史上最早实用的电机，至今仍在需要宽范围平滑调速、大起动转矩的场合（轧钢、矿山、电车、起重）中大量使用。它的核心是一个\textbf{机械换向器}：把电枢绕组内部的交流电动势"整流"成对外输出的直流，或把外部直流"逆变"成绕组内部的交流。本章介绍原理、结构与额定值。

\section{直流电机的工作原理}

\subsection{直流电动机}
电枢绕组中的导体位于定子主磁极产生的磁场中，通入电流后受电磁力 $f=Bli$ 作用产生转矩，使电枢旋转。若导体固定不动，转矩方向恒定；但导体随转子旋转到另一极下时，磁场方向变了，若不改变电流方向，转矩就会反向，电机无法连续旋转。

解决的办法是\textbf{换向器}：它由许多彼此绝缘、与电枢绕组线圈按一定规律连接的换向片组成，电刷固定不动压在换向器表面。线圈从一个极下转到另一个极下时，电刷使线圈中的电流方向同时改变，因此\textbf{转矩方向始终保持不变}。

因此，换向器的作用是：\textbf{对电动机，把外部直流"逆变"为电枢绕组内部的交流（不断改向的电流）；对发电机，把绕组内部的交流电动势"整流"为对外输出的直流。}换向器称为\textbf{机械式逆变器/整流器}。

\subsection{直流发电机}
原动机拖动电枢旋转，导体切割主磁通感应电动势 $e=Blv$，方向由右手定则确定；同样由于导体交替进入不同极性的磁极下，绕组内部电动势是交变的，经换向器与电刷才输出方向不变的直流电动势。

\subsection{可逆性}
同一台直流电机，通入直流电流产生转矩、拖动机械负载，即为电动机；由原动机拖动产生电动势、输出直流电能，即为发电机。故直流电机同样具有可逆性，与第 2 章"能量守恒与机电能量转换"的结论一致。

\subsection{换向器工作的分步描述}
以最简单的单匝线圈为例（线圈两端分别接两个换向片，电刷压在换向片上）：
\begin{enumerate}
  \item 线圈边 $ab$ 在 $N$ 极下、$cd$ 在 $S$ 极下，电流由 $a\to b\to c\to d$，两导体所受电磁力方向一致，产生同一方向的转矩；
  \item 电枢转过 $180^\circ$ 后，$ab$ 到了 $S$ 极下、$cd$ 到了 $N$ 极下；若电流方向不变，转矩就会反向；
  \item 恰在此刻，两换向片被同一电刷短接，线圈中电流换向；随后电刷接触的换向片互换，线圈中电流反向；
  \item 于是 $ab$ 在 $S$ 极下时电流方向已反过来，转矩方向仍与原来一致——电枢得以连续旋转。
\end{enumerate}
这个过程每转半周重复一次，故电枢绕组内部的电流是\textbf{交变}的，而对外（电刷端）的电流方向恒定。这就是"机械换向"的本质：\textbf{靠机械接触的位置变化，在电流过零的瞬间完成换向}。

\subsection{与交流电机的对比}
\begin{center}
\begin{tabular}{@{}p{2.4cm}p{4.4cm}p{4.4cm}@{}}
\toprule
项目 & 直流电机 & 交流电机 \\
\midrule
磁场 & 定子主磁极（直流励磁） & 定子绕组产生旋转磁场 \\
电枢 & 转子（旋转） & 定子（静止） \\
绕组内电流/电动势 & 交流（靠换向器整流/逆变） & 交流（本身就是交流） \\
对外电流 & 直流 & 交流 \\
薄弱环节 & 换向器与电刷 & 无（笼型） \\
\bottomrule
\end{tabular}
\end{center}

\section{直流电机的结构}

直流电机由定子（静止部分）和转子（电枢，旋转部分）组成，两者之间有气隙。

\subsection{定子部分}
\begin{itemize}
  \item \textbf{主磁极}：由磁极铁芯（薄钢板叠成）和套在其上的\textbf{励磁绕组}构成，用于建立主磁场。磁极下部的扩大部分称\textbf{极靴}，使气隙磁密分布合理。
  \item \textbf{换向极}：装在相邻主磁极之间（几何中性线上），绕组与电枢绕组串联，用于改善换向（见第 25 章）。
  \item \textbf{机座}：既作机械支撑，又作为主磁路的一部分（磁轭）。
  \item \textbf{电刷装置}：由电刷、刷握、刷杆和连线组成，把电枢绕组与外电路连接起来。
\end{itemize}

\subsection{转子（电枢）部分}
\begin{itemize}
  \item \textbf{电枢铁芯}：由硅钢片叠成，外圆冲有齿槽以嵌放电枢绕组，同时构成主磁路。
  \item \textbf{电枢绕组}：由绝缘铜线绕成线圈嵌入槽内，按一定规律与换向片连接（单叠或单波，见第 21 章）。
  \item \textbf{换向器}：由许多梯形铜换向片与片间云母绝缘组成，是直流电机的关键部件，也是最易出故障的部件。
  \item \textbf{转轴、风扇}：支撑与冷却。
\end{itemize}

\begin{center}
\begin{tabular}{@{}p{2.2cm}p{4.6cm}p{5.2cm}@{}}
\toprule
部件 & 作用 & 备注 \\
\midrule
主磁极 & 建立主磁场 & 励磁绕组通直流 \\
换向极 & 改善换向、抑制火花 & 绕组与电枢串联 \\
机座 & 机械支撑 + 磁轭 & 构成磁路的一部分 \\
电枢铁芯 & 嵌放绕组 + 主磁路 & 硅钢片叠压 \\
电枢绕组 & 感应电动势、产生转矩 & 经换向片连接 \\
换向器 & 机械整流/逆变 & 直流电机的核心与薄弱环节 \\
电刷 & 电枢与外电路的滑动连接 & 需定期更换 \\
\bottomrule
\end{tabular}
\end{center}

\section{直流电机的额定值}

直流电机的额定值都标注在铭牌上，是电机长期安全运行的规定值。

\subsection{额定功率 $P_N$}
\textbf{电动机}：额定运行时轴上输出的机械功率（kW）；\textbf{发电机}：额定运行时出线端输出的电功率（kW）。注意两者定义不同。

\subsection{额定电压 $U_N$}
额定运行时的端电压（V）。直流电机额定电压等级常见为 110、220、440、750\,V 等。

\subsection{额定电流 $I_N$}
额定运行时出线端的电流（A）。对电动机，$I_N$ 是输入电流；对发电机，$I_N$ 是输出电流。

\subsection{额定转速 $n_N$}
额定运行时的转速（r/min）。

\subsection{额定励磁电压 $U_{fN}$ 与额定励磁电流 $I_{fN}$}
额定运行时的励磁电压与励磁电流。

\subsection{额定效率 $\eta_N$}
额定运行时输出功率与输入功率之比。

\begin{center}
\begin{tabular}{@{}lll@{}}
\toprule
额定值 & 电动机 & 发电机 \\
\midrule
$P_N$ & 轴上输出机械功率 & 出线端输出电功率 \\
$U_N$ & 输入端电压 & 输出端电压 \\
$I_N$ & 输入电流 & 输出电流 \\
基本关系 & $P_N=U_NI_N\eta_N$ & $P_N=U_NI_N$ \\
\bottomrule
\end{tabular}
\end{center}

\keyformula{\text{电动机：}P_1=U_NI_N,\ P_N=U_NI_N\eta_N;\qquad
\text{发电机：}P_2=U_NI_N,\ P_1=\frac{U_NI_N}{\eta_N}}

\corepoint{直流电机额定值最容易混的是 $P_N$ 的定义：电动机的 $P_N$ 是\textbf{输出机械功率}，发电机的 $P_N$ 是\textbf{输出电功率}。因此同样铭牌"$P_N=10$\,kW、$U_N=220$\,V"，电动机的额定电流是 $10\,000/(220\eta_N)$，发电机的额定电流是 $10\,000/220$——一个要除效率、一个不除。做题先问"这是电动机还是发电机"。}

% ============================================================
% 第 21 章 直流电机的电枢绕组（教材第六篇，印刷页 235--243）
% 节名：第一节 直流电枢绕组的构成；第二节 单叠绕组；第三节 单波绕组
% ============================================================
\chapter{直流电机的电枢绕组}

直流电机的电枢绕组是产生电动势、通过电流产生转矩的核心部件。与交流绕组不同，它的线圈经换向片串联成一个\textbf{闭合回路}，再通过电刷分成若干并联支路。本章讨论两种基本型式：单叠绕组与单波绕组，重点掌握它们的\textbf{节距计算}与\textbf{并联支路数}。

\section{直流电枢绕组的构成}

\subsection{元件与换向片}
电枢绕组由许多\textbf{绕组元件}（线圈）组成，每个元件的两个端头分别接到两个相邻的换向片上。元件数 $S$、换向片数 $K$、虚槽数 $u$ 之间满足
\begin{equation}
S=K=uQ_a ,
\end{equation}
其中 $Q_a$ 为电枢槽数，$u$ 为每槽的虚槽数（每槽内并放的元件边对数）。为了分析方便，常把每个虚槽的上、下层元件边分别编号，用"虚槽"代替实际槽。

\subsection{节距}
直流电枢绕组有四个基本节距（均以\textbf{虚槽数}计）：
\begin{itemize}
  \item \textbf{第一节距} $y_1$：一个元件的两个元件边在电枢表面的距离，$y_1\approx\tau=\dfrac{Q_a}{2p}$（整距）；
  \item \textbf{第二节距} $y_2$：前一元件的下层边与后一元件的上层边之间的距离；
  \item \textbf{合成节距} $y$：相邻两元件对应边之间的距离，$y=y_1-y_2$（叠绕组）或 $y=y_1+y_2$（波绕组）；
  \item \textbf{换向器节距} $y_k$：相邻两元件所接换向片之间的距离，$y_k=y$。
\end{itemize}
为使 $y_1$ 取整数，$y_1$ 可整距、短距或长距，工程上多用整距或短距。

\subsection{并联支路数}
电枢绕组经电刷分成若干并联支路，设\textbf{支路对数}为 $a$，则电枢总电流 $I_a$ 与每条支路电流 $i_a$ 的关系为
\begin{equation}
I_a=2a\,i_a .
\end{equation}
$a$ 的大小由绕组型式决定，这是两种绕组最本质的区别。

\section{单叠绕组}

\subsection{连接规律}
单叠绕组的特点是\textbf{每个元件的末端接到相邻的下一个换向片上}，即
\begin{equation}
y=y_k=\pm1 .
\end{equation}
通常取 $y=+1$（右行绕组）。合成节距为
\begin{equation}
y=y_1-y_2=1\quad\Longrightarrow\quad y_2=y_1-1 .
\end{equation}
元件依次"叠"在一起前进，绕电枢一周后回到出发点附近，$p$ 对极电机需绕 $p$ 周才能把所有元件串完。

\subsection{并联支路数}
单叠绕组中，每个主磁极下的元件串联成一条支路，故
\begin{equation}
2a=2p\quad\Longrightarrow\quad a=p .
\end{equation}

\keyformula{\text{单叠绕组：}y=y_k=1,\qquad a=p,\qquad 2a=2p}

\subsection{特点}
\begin{itemize}
  \item 并联支路数多（$2a=2p$），每条支路电流小，适合\textbf{大电流、低电压}的电机；
  \item 电刷数目与极数相同，换向片间的电压较低；
  \item 常用于中小型直流电机及需要大电流的场合。
\end{itemize}

\section{单波绕组}

\subsection{连接规律}
单波绕组的特点是\textbf{每个元件的末端接到相隔较远的换向片上}，元件像波浪一样沿电枢前进，绕一周后回到出发点\textbf{相邻}的换向片，再继续第二周。其合成节距为
\begin{equation}
y=y_k=\frac{K\pm1}{p},
\end{equation}
取"$-$"号为左行绕组（常用），"$+$"号为右行绕组。合成节距与第一节距的关系为
\begin{equation}
y=y_1+y_2\quad\Longrightarrow\quad y_2=y-y_1 .
\end{equation}

\subsection{并联支路数}
单波绕组把所有同一极性下的元件都串联起来，只形成\textbf{两条}并联支路，与极数无关：
\begin{equation}
2a=2\quad\Longrightarrow\quad a=1 .
\end{equation}

\keyformula{\text{单波绕组：}y=y_k=\frac{K\pm1}{p},\qquad a=1,\qquad 2a=2}

\subsection{特点}
\begin{itemize}
  \item 并联支路数少（$2a=2$），每条支路串联元件多、电压高，适合\textbf{高电压、小电流}的电机；
  \item 同样容量下所需电刷数少（理论上只需两组电刷即可，实用上仍按极数布置以减小换向片间电压）；
  \item 常用于额定电压较高的直流电机及某些特殊场合（如电车牵引电机）。
\end{itemize}

\begin{center}
\begin{tabular}{@{}p{2.0cm}p{3.4cm}p{3.6cm}p{4.2cm}@{}}
\toprule
项目 & 单叠绕组 & 单波绕组 & 记忆要点 \\
\midrule
合成节距 & $y=y_k=\pm1$ & $y=y_k=\dfrac{K\pm1}{p}$ & 叠"1"、波"跨极" \\
$y_1$、$y_2$ 关系 & $y=y_1-y_2$ & $y=y_1+y_2$ & 叠减、波加 \\
支路对数 $a$ & $a=p$ & $a=1$ & 叠随极数、波恒为 1 \\
并联支路数 & $2a=2p$ & $2a=2$ & — \\
均压线 & 需要 & 一般不需要 & 由磁路不对称决定 \\
适用 & 大电流、低电压 & 高电压、小电流 & 由 $I_a=2ai_a$ 判断 \\
\bottomrule
\end{tabular}
\end{center}

\section{绕组连接步骤与均压线}

\subsection{连接的实际步骤}
以单叠绕组为例，画绕组展开图（或计算节距）的一般步骤如下：
\begin{enumerate}
  \item 由 $Q_a$、$u$、$p$ 求虚槽数 $Q=uQ_a$、元件数 $S=K=Q$、极距 $\tau=Q/(2p)$；
  \item 取第一节距 $y_1\le\tau$ 且为整数（整距或短距）；
  \item 由 $y=y_k=\pm1$ 定合成节距，再由 $y_2=y_1-y$ 求第二节距，检查 $y_2$ 是否为整数且合理；
  \item 从 1 号换向片（1 号元件上层边）出发，按 $+y_1$（上层边到下层边）、$-y_2$（下层边到下一元件上层边）逐元件编号，直到回到起点；
  \item 检查是否把所有元件都用上（单叠绕组需绕 $p$ 周），并确定电刷位置（放在几何中性线上，与主磁极轴线相差 $90^\circ$ 电角度处）。
\end{enumerate}
单波绕组的步骤相同，只是用 $y=(K\pm1)/p$、$y_2=y-y_1$，且绕行一周后回到起点\textbf{相邻}的换向片，继续第二周直到闭合。

\subsection{均压线}
单叠绕组中，若各极下磁路不对称（气隙不均、材料不均），同一极性各支路的感应电动势不相等，会在绕组内部产生\textbf{环流}，使损耗增大、发热不均。为此在电枢上装设\textbf{均压线}，把理论上等电位的换向片连接起来，让环流在均压线中流通而不经过电刷。

单波绕组的所有元件串联成两条支路，本身对磁路不对称不敏感，一般不需要均压线（这也是它在低速、多极电机中受欢迎的原因之一）。

\corepoint{两种绕组只需记两组式子：\textbf{单叠} $y=1$、$a=p$、$y_2=y_1-y$；\textbf{单波} $y=(K\pm1)/p$、$a=1$、$y_2=y-y_1$。判断"该选哪种"的标准是电流与电压：绕组并联支路越多，每条支路分担的电流越小、总电压越低——所以大电流选单叠、高电压选单波。另外注意 $y$、$y_1$、$y_2$、$y_k$ 都要取整数，$y_1$ 须满足 $y_1\le\tau$ 且为整数。}

% ============================================================
% 第 22 章 直流电机的磁场（教材第六篇，印刷页 244--253）
% 节名：第一节 直流电机的励磁方式；第二节 空载时直流电机的磁场；
%       第三节 负载时直流电机的磁场及电枢反应；第四节 直流电机的感应电动势和电磁转矩
% ============================================================
\chapter{直流电机的磁场}

直流电机的磁场由\textbf{主磁极的励磁绕组}建立（这部分称为主磁场），负载后电枢绕组电流又产生\textbf{电枢磁场}，两者叠加即为气隙磁场。电枢磁场对主磁场的影响称为\textbf{电枢反应}，它是直流电机一系列运行特性（转矩、转速、换向）的根源。

\section{直流电机的励磁方式}

按励磁绕组与电枢绕组的连接方式，直流电机分为四类。

\subsection{他励}
励磁绕组由\textbf{独立的直流电源}供电，与电枢回路无电的联系。特点是励磁电流不受电枢电压影响，便于调节，常用于需要宽范围调速的场合（如轧钢电动机）及发电机试验。

\subsection{并励}
励磁绕组与电枢绕组\textbf{并联}，励磁电流 $I_f=U/R_f$ 取自电枢端电压。并励电动机的励磁电流随端电压变化，具有"自励"特性；并励发电机是最常用的直流发电机。

\subsection{串励}
励磁绕组与电枢绕组\textbf{串联}，励磁电流即为电枢电流 $I_f=I_a$。串励电动机起动转矩大、过载能力强，但轻载时转速极高（"飞车"），绝不能空载运行。

\subsection{复励}
主磁极上有两套励磁绕组，一套与电枢并联（并励绕组），一套与电枢串联（串励绕组）。若两者磁动势方向相同称\textbf{积复励}（常用），相反称\textbf{差复励}。复励兼顾了并励的稳定性和串励的大转矩。

\begin{center}
\begin{tabular}{@{}p{1.8cm}p{3.4cm}p{3.0cm}p{4.4cm}@{}}
\toprule
励磁方式 & 连接 & 励磁电流 & 特点与用途 \\
\midrule
他励 & 独立电源 & $I_f$ 独立可调 & 调速范围宽，轧钢、试验 \\
并励 & 与电枢并联 & $I_f=U/R_f$ & 特性较硬，常用发电机、机床 \\
串励 & 与电枢串联 & $I_f=I_a$ & 起动转矩大，电车、起重机；忌空载 \\
复励 & 并励 + 串励 & 两者合成 & 兼顾稳定性与起动转矩 \\
\bottomrule
\end{tabular}
\end{center}

\section{空载时直流电机的磁场}

空载（电枢电流为零或很小）时，气隙磁场完全由主磁极励磁磁动势建立。

\subsection{主磁场分布}
主磁极的极靴做得比磁极本体宽，使极靴下气隙较均匀、极间气隙很大。因此气隙磁密沿圆周的分布接近\textbf{平顶波}：极靴下磁密基本恒定，极间迅速下降。工程上取\textbf{平均磁密} $B_{av}$ 进行计算，它与最大磁密的关系由\textbf{极弧系数} $\alpha$ 描述：$B_{av}=\alpha B_\delta$（$\alpha$ 一般 0.6～0.75）。

每极磁通
\begin{equation}
\Phi=B_{av}\,\tau\,l=\alpha B_\delta\tau l ,
\end{equation}
其中 $\tau=\pi D_a/(2p)$ 为极距，$l$ 为电枢有效长度。

\subsection{磁化曲线}
空载时改变励磁电流 $I_f$，端电压（或电枢电动势）随之变化，得\textbf{空载特性} $E_0=f(I_f)$。它实质上就是电机的磁化曲线：起始段近似直线（气隙线），随后因铁芯饱和而弯曲。磁路饱和程度直接影响电枢反应的大小与换向的好坏。

\subsection{漏磁}
主磁极的磁通除穿过气隙进入电枢的主磁通外，还有一部分不经电枢、在磁极之间直接闭合，称为\textbf{漏磁通}。漏磁使磁极与机座需要更大的截面，用\textbf{漏磁系数} $\sigma=(\Phi+\Phi_\sigma)/\Phi\approx1.1\sim1.2$ 描述。

\section{负载时直流电机的磁场及电枢反应}

\subsection{电枢磁场}
电枢绕组中流过电流 $I_a$ 时产生电枢磁动势。由于电刷通常放在\textbf{几何中性线}上，电枢磁动势的轴线与主磁极轴线相差 $90^\circ$ 电角度，故其性质为\textbf{交轴电枢反应}。

\subsection{交轴电枢反应的影响}
电枢磁场与主磁场叠加后：
\begin{itemize}
  \item 主磁极的\textbf{前极尖}（进入端）磁通增强、\textbf{后极尖}（退出端）磁通削弱，气隙磁密分布\textbf{畸变}，不再是平顶波；
  \item 由于铁芯饱和，增强的一侧饱和程度加深、增磁有限，而削弱的一侧去磁作用充分，故合成磁通\textbf{略有减小}（一般减小百分之几）；
  \item 磁密的零点（\textbf{物理中性线}）不再与几何中性线重合，而是发生偏移。对\textbf{发电机}，物理中性线\textbf{顺着电枢旋转方向}偏移；对\textbf{电动机}，\textbf{逆着旋转方向}偏移。
\end{itemize}

\subsection{直轴电枢反应}
若电刷不在几何中性线上（例如为改善换向而移刷），电枢磁动势将出现\textbf{直轴分量}：电刷顺转向移动时产生去磁作用，逆转向移动时产生助磁作用。此外，换向极绕组若设计不当也会引入直轴磁动势。

\subsection{电枢反应的不良后果}
\begin{enumerate}
  \item 气隙磁密畸变使\textbf{换向条件恶化}，电刷下产生火花（第 25 章）；
  \item 磁通略有减小，使电动势、转矩下降，需略微增大励磁补偿；
  \item 对大型电机，畸变严重时可能引起环火。
\end{enumerate}

\subsection{抑制措施}
\begin{itemize}
  \item 装\textbf{换向极}：产生直轴磁动势抵消交轴电枢反应，并改善换向（第 25 章）；
  \item 装\textbf{补偿绕组}：嵌在主磁极极靴的槽内，与电枢绕组串联，其磁动势直接抵消交轴电枢反应磁动势，用于负载重、过载大、换向困难的大型电机。
\end{itemize}

\begin{center}
\begin{tabular}{@{}p{2.2cm}p{4.2cm}p{5.0cm}@{}}
\toprule
项目 & 空载 & 负载 \\
\midrule
磁场来源 & 仅励磁磁动势 & 励磁磁动势 + 电枢磁动势 \\
气隙磁密 & 平顶波（极靴下均匀） & 畸变，前后极尖一边增一边减 \\
磁通量 & $\Phi_0$ & $\Phi\approx(0.95\sim0.98)\Phi_0$（略减） \\
物理中性线 & 与几何中性线重合 & 偏移（发电机顺转向、电动机逆转向） \\
后果 & — & 换向恶化、火花、转矩略降 \\
\bottomrule
\end{tabular}
\end{center}

\corepoint{直流电机电枢反应的三句话必须背熟：① 电刷在几何中性线上时是\textbf{交轴电枢反应}，只使磁场\textbf{畸变}，不改变磁通方向；② 畸变的具体表现是"前极尖增磁、后极尖去磁"，因饱和而\textbf{净磁通略减}，\textbf{物理中性线偏移}；③ 偏移方向：发电机\textbf{顺}转向、电动机\textbf{逆}转向。凡是问"电枢反应有什么影响"，按"畸变 $\rightarrow$ 磁通略减 $\rightarrow$ 中性线偏移 $\rightarrow$ 换向变坏"这四步答。}

\section{直流电机的感应电动势和电磁转矩}

\subsection{感应电动势}
电枢旋转时，绕组切割气隙磁通感应电动势。设电枢总导体数为 $N$、并联支路对数为 $a$、每极磁通为 $\Phi$、转速为 $n$，则电枢电动势为
\begin{equation}
E_a=C_e\Phi n,
\end{equation}
其中
\begin{equation}
C_e=\frac{pN}{60a}
\end{equation}
为\textbf{电动势常数}（由电机结构决定）。

\keyformula{E_a=C_e\Phi n,\qquad C_e=\frac{pN}{60a}}

\subsection{电磁转矩}
电枢电流 $I_a$ 在磁场中受力产生转矩：
\begin{equation}
T=C_T\Phi I_a,
\end{equation}
其中
\begin{equation}
C_T=\frac{pN}{2\pi a}
\end{equation}
为\textbf{转矩常数}。两者关系为
\begin{equation}
C_T=\frac{60}{2\pi}C_e=9.55\,C_e .
\end{equation}

\keyformula{T=C_T\Phi I_a,\qquad C_T=\frac{pN}{2\pi a}=9.55C_e}

\subsection{电磁功率}
由 $P_{em}=E_aI_a$ 与 $P_{em}=T\Omega$ 可得
\begin{equation}
E_aI_a=C_e\Phi n\cdot I_a=\frac{C_T\Phi I_a\cdot2\pi n}{60}=T\Omega ,
\end{equation}
即电磁功率既等于电枢吸收（或发出）的电功率，又等于转换为机械（或电）功率的那部分——这是第 2 章 $ei=fv$ 的直流电机版本。

\begin{center}
\begin{tabular}{@{}lll@{}}
\toprule
量 & 表达式 & 说明 \\
\midrule
电枢电动势 & $E_a=C_e\Phi n$，$C_e=pN/(60a)$ & 与转速、磁通成正比 \\
电磁转矩 & $T=C_T\Phi I_a$，$C_T=pN/(2\pi a)$ & 与电流、磁通成正比 \\
常数关系 & $C_T=9.55C_e$ & 由 $\Omega=2\pi n/60$ 导出 \\
电磁功率 & $P_{em}=E_aI_a=T\Omega$ & 机电能量转换的桥梁 \\
\bottomrule
\end{tabular}
\end{center}

\corepoint{$E_a=C_e\Phi n$ 与 $T=C_T\Phi I_a$ 是直流电机的"两把钥匙"，后面所有特性都由这两式推出。注意三点：① $\Phi$ 是\textbf{每极磁通}，不是总磁通；② $C_e$、$C_T$ 只与结构（$p$、$N$、$a$）有关，与运行状态无关；③ 由 $U=E_a+I_aR_a$（电动机）或 $U=E_a-I_aR_a$（发电机）可见，$E_a$ 与 $U$ 之间的差就是电枢电阻压降——这一条是分析机械特性与起动电流的根基。}

% ============================================================
% 第 23 章 直流发电机（教材第六篇，印刷页 254--258）
% 节名：第一节 直流发电机的基本电磁关系；第二节 他励直流发电机的运行特性；第三节 并励直流发电机的自励及运行特性
% ============================================================
\chapter{直流发电机}

直流发电机由原动机拖动，把机械能转换为直流电能。它已大量被整流电源取代，但其"自励建压"的物理过程仍是直流电机理论中最精彩的部分，也是理解并励电动机与励磁系统的钥匙。

\section{直流发电机的基本电磁关系}

\subsection{电动势与电压方程}
电枢旋转感应电动势 $E_a=C_e\Phi n$。作为发电机，电枢输出电流，故电枢回路有电阻压降，端电压为
\begin{equation}
U=E_a-I_aR_a ,
\end{equation}
其中 $R_a$ 为电枢回路总电阻（含电枢绕组电阻与电刷接触电阻）。

\subsection{电磁转矩}
$T=C_T\Phi I_a$。发电机中电磁转矩是\textbf{制动转矩}，与原动机拖动转矩方向相反；稳态时 $T_1=T+T_0$。

\subsection{功率关系}
\begin{equation}
P_1=P_{em}+p_{mec}+p_{Fe}+p_{ad},\qquad P_{em}=E_aI_a,\qquad P_2=UI_a=P_{em}-p_{Cu a},
\end{equation}
其中 $P_1$ 为原动机输入功率，$P_2$ 为电输出功率，$p_{mec}$ 为机械损耗，$p_{Fe}$ 为铁耗，$p_{ad}$ 为杂散损耗，$p_{Cua}=I_a^2R_a$ 为电枢铜耗（对并励机还应计入励磁损耗 $U I_f$）。

\subsection{功率流程与效率}
把上述关系画成流程：原动机输入 $P_1$ 扣除机械损耗、铁耗与杂散损耗后成为\textbf{电磁功率} $P_{em}=E_aI_a$；电磁功率扣除电枢铜耗后即为\textbf{输出电功率} $P_2=UI_a$。效率为
\begin{equation}
\eta=\frac{P_2}{P_1}\times100\%=\frac{UI_a}{UI_a+p_{Cua}+p_{mec}+p_{Fe}+p_{ad}}\times100\% .
\end{equation}
其中铁耗与机械损耗近似为\textbf{不变损耗}（与负载无关，只与转速、磁通有关），铜耗为\textbf{可变损耗}（$\propto I_a^2$）。与变压器、感应电机一样，当可变损耗等于不变损耗时效率最高。

对并励发电机，励磁功率由电枢自身供给，故输出电流 $I=I_a-I_f$，输出功率 $P_2=UI$；计算效率时既可用"输入机械功率"作分母，也可用"输出电功率 $+$ 全部损耗"作分母，两者必须一致（不要重复计入励磁损耗）。

\section{他励直流发电机的运行特性}

他励发电机的励磁由独立电源供给，$I_f$ 可独立调节，是最基本的分析对象。

\subsection{空载特性}
$n=n_N$、$I_a=0$ 时，$U_0=E_a=f(I_f)$。这就是电机的磁化曲线：起始段线性（气隙线），随后弯曲饱和。空载特性是求磁路饱和程度与电枢反应大小的基础。

\subsection{外特性}
$n=n_N$、$I_f=$ 常数时，$U=f(I_a)$。随负载电流增大，$U$ 下降，原因有二：① 电枢电阻压降 $I_aR_a$；② 电枢反应的去磁作用使 $\Phi$ 略减。外特性是一条略向下弯的曲线。

\subsection{调节特性}
$n=n_N$、$U=$ 常数时，$I_f=f(I_a)$。为维持端电压不变，负载增加时须增大励磁电流，故曲线略上翘。

\subsection{电压变化率}
从空载到额定负载，端电压的相对下降称为电压变化率：
\begin{equation}
\Delta U=\frac{U_0-U_N}{U_N}\times100\% ,
\end{equation}
他励发电机一般约 5\%～10\%。

\section{并励直流发电机的自励及运行特性}

\subsection{自励建压过程}
并励发电机的励磁绕组与电枢并联，励磁电流取自电枢端电压。它不需要外部直流电源即可建立电压，过程如下：
\begin{enumerate}
  \item 电机由原动机拖动到额定转速，电枢绕组切割主磁极的\textbf{剩磁}，感应一个不大的剩磁电动势 $E_r$（约为额定电压的 2\%～5\%）；
  \item $E_r$ 加在励磁绕组上，产生不大的励磁电流 $I_f$；
  \item 若该励磁电流产生的磁通与剩磁\textbf{同方向}，则磁通增强，$E_a$ 增大，$I_f$ 进一步增大……
  \item 如此循环，电压不断上升，直到励磁绕组的伏安特性 $U=I_fR_f$ 与电机的空载特性 $E_0=f(I_f)$ 相交于稳定点为止。
\end{enumerate}

\subsection{自励的三个条件}
\begin{enumerate}
  \item 主磁极必须\textbf{有剩磁}（否则须先用外电源充磁）；
  \item 励磁绕组与电枢的\textbf{接法必须正确}，使励磁电流产生的磁通与剩磁同向（接反则电压建不起来，甚至反转极性）；
  \item 励磁回路的\textbf{电阻必须小于临界电阻} $R_{f,cr}$。若 $R_f$ 过大，励磁特性线 $U=I_fR_f$ 与空载特性的交点过低，电压只能建立到很低的数值。
\end{enumerate}

\keyformula{R_f<R_{f,cr}\quad\text{（自励建压的必要条件）}}

\subsection{外特性}
并励发电机的外特性 $U=f(I_a)$ 比他励机\textbf{更陡}：负载增加时，除电枢电阻压降与电枢反应去磁外，端电压下降还使励磁电流减小，进一步削弱磁场。因此并励机存在\textbf{临界电流}：负载电阻减小时，电流增大到某一值后反而减小（因电压崩溃）。

\subsection{建压失败的排查}
并励发电机建不起电压是常见的故障现象，按自励三条件逐一排查：

\begin{center}
\begin{tabular}{@{}p{3.2cm}p{4.0cm}p{4.8cm}@{}}
\toprule
现象 & 可能原因 & 处理 \\
\midrule
电压表指示为零或极小 & 无剩磁 & 用外电源充磁 \\
电压建立但极性相反 & 励磁绕组接反 & 对调励磁绕组两端 \\
电压只能建立到很低值 & 励磁回路电阻过大（超过临界电阻） & 减小磁场变阻器阻值 \\
电压建立后不稳、振荡 & 转速过低或电刷接触不良 & 提高转速、研磨/更换电刷 \\
\bottomrule
\end{tabular}
\end{center}

\subsection{并励与他励的比较}
\begin{center}
\begin{tabular}{@{}p{2.2cm}p{4.0cm}p{4.4cm}@{}}
\toprule
项目 & 他励发电机 & 并励发电机 \\
\midrule
励磁电源 & 独立直流电源 & 取自电枢端电压 \\
建压 & 加励磁即有电压 & 需剩磁自励建压 \\
外特性 & 略下降 & 下降更陡，有临界电流 \\
适用 & 宽范围调压、试验 & 常用直流电源 \\
\bottomrule
\end{tabular}
\end{center}

\corepoint{自励建压这一节是"正反馈"的经典例子：剩磁 $\rightarrow$ 小电动势 $\rightarrow$ 小励磁电流 $\rightarrow$ 磁通增强 $\rightarrow$ 电动势增大 $\rightarrow$ …直到饱和与电阻压降把增长"刹住"。三个条件缺一不可：\textbf{有剩磁、接法对（助磁）、电阻小于临界值}。考试常问"并励发电机建不起电压的可能原因"，就按这三条逐一排查（另外还要检查转速是否过低、电刷接触是否良好）。}

% ============================================================
% 第 24 章 直流电动机（教材第六篇，印刷页 259--269）
% 节名：第一节 直流电动机的基本电磁关系；第二节 直流电动机的运行特性；
%       第三节 直流电动机的启动、调速及电磁制动
% ============================================================
\chapter{直流电动机}

直流电动机最大的优点是\textbf{调速范围宽、平滑、起动转矩大}，因此在轧钢、矿山、电车、起重等需要低速大转矩的场合仍不可替代。本章的核心只有一个式子：由 $U=E_a+I_aR_a$ 与 $E_a=C_e\Phi n$ 解出的\textbf{转速公式}——它同时解释了机械特性、起动限流和三种调速方式。

\section{直流电动机的基本电磁关系}

\subsection{电压方程与转速公式}
电动机接直流电源，电枢电流 $I_a$ 由电网流入（与 $E_a$ 反方向），故
\begin{equation}
U=E_a+I_aR_a\quad\Longrightarrow\quad
I_a=\frac{U-E_a}{R_a}.
\end{equation}
代入 $E_a=C_e\Phi n$，得\textbf{转速公式}
\begin{equation}
n=\frac{U-I_aR_a}{C_e\Phi}.
\end{equation}

\keyformula{n=\frac{U-I_aR_a}{C_e\Phi}}

这个式子说明：\textbf{转速由三个量决定}——端电压 $U$、电枢回路电阻压降 $I_aR_a$、每极磁通 $\Phi$。于是调速也只有三条路：调 $U$、调 $\Phi$、调 $R_a$（串电阻）。全部调速与起动分析都由此展开。

\subsection{转矩方程}
\begin{equation}
T=C_T\Phi I_a .
\end{equation}
稳态时 $T=T_2+T_0$（电磁转矩 $=$ 负载转矩 $+$ 空载制动转矩）。负载增大时，转速略降、$E_a$ 减小、$I_a$ 增大，转矩随之增大，直至与负载平衡——这是电动机的自动调节过程。

\subsection{各励磁方式的电流关系}
\begin{itemize}
  \item \textbf{并励}：$I_f=U/R_f$，$I_a=I-I_f$（$I$ 为电网输入电流）；
  \item \textbf{串励}：$I_f=I_a=I$；
  \item \textbf{他励}：$I_a=I$，$I_f$ 独立；
  \item \textbf{复励}：$I_a=I-I_{f\text{（并）}}$，串励绕组电流为 $I_a$。
\end{itemize}

\section{直流电动机的运行特性}

\subsection{机械特性}
机械特性 $n=f(T)$ 是电动机最重要的特性，由转速公式与 $I_a=T/(C_T\Phi)$ 联立得
\begin{equation}
n=\frac{U}{C_e\Phi}-\frac{R_a}{C_eC_T\Phi^{2}}T=n_0-\beta T ,
\end{equation}
其中 $n_0=U/(C_e\Phi)$ 为\textbf{理想空载转速}，$\beta=R_a/(C_eC_T\Phi^2)$ 为斜率。

\subsubsection{并励（他励）电动机}
$\Phi$ 基本恒定，故机械特性是一条\textbf{略微下倾的直线}，属\textbf{硬特性}：负载从空载到额定，转速下降仅 5\%～10\%。定义转速变化率
\begin{equation}
\Delta n\%=\frac{n_0-n_N}{n_N}\times100\% ,
\end{equation}
一般 $<10\%$，适合要求转速平稳的机床、造纸、轧钢（配合调压调速）。

\subsubsection{串励电动机}
$\Phi$ 随 $I_a$ 变化（磁路未饱和时 $\Phi\propto I_a$），故 $T\propto I_a^2$，转速
\begin{equation}
n\propto\frac{U}{\sqrt{T}} ,
\end{equation}
机械特性为\textbf{软特性}：轻载时转速极高，重载时转速迅速下降。串励电动机起动转矩大（$T\propto I_a^2$），适合电车、起重机、牵引；但\textbf{绝对不允许空载或轻载运行}，否则转速将升高到危险程度（"飞车"），故禁止用皮带传动。

\subsubsection{复励电动机}
介于两者之间，既有较大的起动转矩，又不致空载飞车。

\begin{center}
\begin{tabular}{@{}p{2.0cm}p{3.4cm}p{3.6cm}p{4.2cm}@{}}
\toprule
励磁方式 & 机械特性 & 转速变化率 & 典型用途 \\
\midrule
并励/他励 & 硬特性（直线略下倾） & 5\%～10\% & 机床、轧钢、造纸 \\
串励 & 软特性（$n\propto1/\sqrt T$） & 很大 & 电车、起重机（忌空载） \\
复励 & 介于两者之间 & 中等 & 需要大起动转矩又怕飞车 \\
\bottomrule
\end{tabular}
\end{center}

\subsection{工作特性}
并励电动机的工作特性指 $U=U_N$、$I_f=I_{fN}$ 时，$n$、$T$、$I$、$\eta$ 随输出功率 $P_2$ 的变化：
\begin{itemize}
  \item $n=f(P_2)$：略下降（硬特性）；
  \item $T=f(P_2)$：近似直线（$T=T_0+P_2/\Omega$）；
  \item $I=f(P_2)$：随负载近似线性增大；
  \item $\eta=f(P_2)$：先升后降，在 $(0.75\sim1.0)P_N$ 附近最高。
\end{itemize}

\section{直流电动机的启动、调速及电磁制动}

\subsection{启动问题与限流}
起动瞬间 $n=0$，$E_a=0$，故
\begin{equation}
I_{st}=\frac{U}{R_a}.
\end{equation}
电枢电阻 $R_a$ 很小，直接起动电流可达额定电流的 \textbf{10～20 倍}，将造成：
\begin{enumerate}
  \item 换向器产生强烈火花甚至环火；
  \item 电网电压骤降；
  \item 电磁转矩 $T=C_T\Phi I_{st}$ 过大，冲击机械负载与传动机构。
\end{enumerate}
因此直流电动机\textbf{不允许直接起动}（微型电机除外）。

\subsubsection{降压起动}
降低电源电压起动，随着转速上升逐步升压。起动电流小、起动平滑，需要可调压直流电源（现代多用晶闸管整流或直流斩波），是当前主流方案。

\subsubsection{电枢串电阻起动}
在电枢回路串入起动电阻 $R_{st}$，使
\begin{equation}
I_{st}=\frac{U}{R_a+R_{st}}\le\left(1.5\sim2\right)I_N ,
\end{equation}
随转速上升逐段切除电阻。设备简单，但起动过程中电阻上消耗大量能量，且为有级起动。

\subsection{调速}
由转速公式 $n=\dfrac{U-I_aR_a}{C_e\Phi}$ 可知三种调速方式。

\subsubsection{调压调速}
保持 $\Phi$ 不变（额定励磁），改变电枢电压 $U$。转速近似与 $U$ 成正比，可\textbf{无级、平滑}调节，调速范围宽（可达 10:1 以上），低速时特性仍硬。因 $\Phi$ 不变，允许的转矩不变，属\textbf{恒转矩调速}——最适合恒转矩负载（起重机、轧机）。这是直流电动机最优的调速方式。

\subsubsection{调磁（弱磁）调速}
保持 $U$ 不变，减小励磁电流以减小 $\Phi$，转速升高。因 $n\propto1/\Phi$，故为\textbf{升速}调速；受换向与机械强度限制，转速一般不超过 $1.2\sim2\,n_N$。由于 $\Phi$ 减小后允许的转矩减小，而功率近似不变，属\textbf{恒功率调速}，适合恒功率负载（机床主轴）。

\subsubsection{电枢串电阻调速}
在电枢回路串入电阻，使机械特性变软，同一转矩下转速降低。设备简单，但：① 为有级调速；② 低速时电阻上损耗大、效率低；③ 特性变软、转速随负载波动大。仅用于短时、小范围调速。

\begin{center}
\begin{tabular}{@{}p{2.0cm}p{2.6cm}p{3.2cm}p{3.4cm}p{2.4cm}@{}}
\toprule
调速方式 & 调节量 & 转速变化 & 调速性质 & 评价 \\
\midrule
调压 & $U$ & $n\propto U$，无级 & 恒转矩 & 最优，范围宽、特性硬 \\
调磁（弱磁） & $\Phi$ & $n\propto1/\Phi$，升速 & 恒功率 & 只能向上调，受换向限制 \\
电枢串电阻 & $R_a$ & 特性变软、转速降低 & 变转矩 & 有级、损耗大，少用 \\
\bottomrule
\end{tabular}
\end{center}

\subsection{电磁制动}
\begin{enumerate}
  \item \textbf{能耗制动}：切断电源，把电枢接到制动电阻上。惯性使电枢继续旋转，$E_a$ 作用下电流反向、转矩变为制动，动能消耗在电阻上。制动平稳、准确。
  \item \textbf{反接制动}：把电枢电压反接（或电枢反接），$E_a$ 与 $U$ 同方向，电流很大，必须串入制动电阻限流。制动强烈、迅速，用于快速停车与位能负载下放。
  \item \textbf{回馈（再生）制动}：当转速高于理想空载转速（如起重机下放重物、电车下坡），$E_a>U$，电流反向流入电网，电机作发电机运行，把机械能回馈电网。最经济，但只能在 $n>n_0$ 时使用。
\end{enumerate}

\begin{center}
\begin{tabular}{@{}p{2.2cm}p{4.2cm}p{3.6cm}p{3.2cm}@{}}
\toprule
制动方式 & 做法 & 能量去向 & 适用 \\
\midrule
能耗制动 & 电枢接制动电阻 & 消耗在电阻上 & 平稳停车 \\
反接制动 & 电压（或电枢）反接 & 消耗在电阻上 & 快速停车、下放 \\
回馈制动 & $n>n_0$ 自然进入发电态 & 回馈电网 & 下放重物、下坡 \\
\bottomrule
\end{tabular}
\end{center}

\subsection{稳定运行与拖动系统}
电动机的稳态工作点是机械特性 $n=f(T)$ 与负载特性 $n=f(T_L)$ 的交点。稳定运行的判据是：转速略有升高时，电动机转矩小于负载转矩（$\mathrm{d}T/\mathrm{d}n<\mathrm{d}T_L/\mathrm{d}n$）。

动态过程中，转速的变化由拖动系统的运动方程决定：
\begin{equation}
T-T_L=J\frac{\mathrm{d}\Omega}{\mathrm{d}t},
\end{equation}
其中 $J$ 为转动惯量。该式说明：电磁转矩大于负载转矩时加速，小于时减速，相等时匀速——这是分析起动时间、制动过程的基础。并励电动机机械特性为下垂直线，与恒转矩负载的交点唯一且稳定；串励电动机特性下垂很陡，轻载时交点上移极快，这正是"忌空载"的数学表现。

\subsection{并励与他励的区别}
两者机械特性形式相同，但并励电动机的励磁电流取自电枢端电压，故：
\begin{itemize}
  \item 调压调速时，$U$ 下降会同时使 $I_f$ 下降、$\Phi$ 减小，导致"弱磁升速"与"降压降速"相互抵消一部分，特性不如他励硬；
  \item 若电枢电压降低过多，励磁可能不足，需改用他励（或自励恒压）以保证励磁稳定。
\end{itemize}
因此要求宽范围调速的场合，多采用\textbf{他励}结构（励磁由独立可控电源供电）。

\corepoint{本章只需牢记一条主线：\textbf{一切从 $n=\dfrac{U-I_aR_a}{C_e\Phi}$ 出发}。起动时 $E_a=0$ 故电流巨大 $\Rightarrow$ 必须降压或串电阻；调速就是调 $U$、$\Phi$、$R_a$ 三个量中的某一个，而调哪个量决定了"恒转矩"还是"恒功率"；制动就是让 $E_a$ 与 $U$ 的关系反过来，使电流反流。串励电动机的"忌空载"也必须记牢：$\Phi\propto I_a$，空载时 $I_a$ 很小、$\Phi$ 很小，$n\propto U/\Phi$ 会飞车。}

% ============================================================
% 第 25 章 直流电机的换向（教材第六篇，印刷页 270--277）
% 节名：第一节 直流电机换向的电磁现象；第二节 改善换向的方法；第三节 补偿绕组
% ============================================================
\chapter{直流电机的换向}

换向是直流电机的"专利问题"：电枢元件从一条支路进入另一条支路时，元件中的电流必须从 $+i_a$ 变为 $-i_a$，这个过程在电刷下几十微秒内完成，极易产生火花。换向不良是直流电机最主要的故障源，也是限制其容量与转速提高的根本原因。

\section{直流电机换向的电磁现象}

\subsection{换向过程}
设电刷宽度等于一个换向片宽度，则被电刷短路的那个元件从开始被短路到脱离短路所经历的时间称为\textbf{换向周期} $T_k$：
\begin{equation}
T_k=\frac{b_k}{v_k},
\end{equation}
其中 $b_k$ 为换向片宽度，$v_k$ 为换向器表面线速度。$T_k$ 通常只有几毫秒到几十微秒，而元件电流却要从 $+i_a$ 变到 $-i_a$。

\subsection{换向元件中的电动势}
换向元件本身是一个电感线圈，电流的急剧变化会产生\textbf{自感电动势}（电抗电动势）：
\begin{equation}
e_L=-L_k\frac{\mathrm{d}i}{\mathrm{d}t}\approx L_k\frac{2i_a}{T_k},
\end{equation}
其方向\textbf{阻碍电流变化}，即力图维持原电流方向——这使换向电流"延迟"改变，在电刷后刷边引起火花。

此外，换向元件位于几何中性线附近，虽然主磁场在此处为零，但\textbf{交轴电枢磁场}仍然存在，换向元件切割该磁场产生\textbf{旋转电动势}（电枢反应电动势）$e_a$，方向同样阻碍换向。

\subsection{三种换向}
\begin{itemize}
  \item \textbf{直线换向}：$e_L+e_a=0$（理想情况），换向电流随时间线性变化，无火花，最理想；
  \item \textbf{延迟换向}：$e_L+e_a>0$ 占主导，电流变化滞后，后刷边电流密度大，\textbf{易产生火花}（实际电机多为此种）；
  \item \textbf{超前换向}：人为加强换向磁场使其过量，电流变化超前，前刷边电流密度大，也会产生火花。
\end{itemize}

\subsection{火花的危害}
火花会灼伤换向器与电刷表面，加剧磨损；严重时形成\textbf{环火}（电弧跨越全部换向片），烧毁电机。因此必须把火花控制在允许范围内（国家标准按火花等级分为 1、1¼、1½、2、3 级）。

\section{改善换向的方法}

改善换向的基本思路是：\textbf{用外加磁场（换向极）产生一个与 $e_L$ 方向相反的电动势，把延迟换向拉回直线换向}。

\subsection{装设换向极}
在相邻主磁极之间的几何中性线上装设\textbf{换向极}，其绕组与电枢绕组\textbf{串联}，使换向极磁动势随电枢电流成正比变化。换向极的极性应使它在换向元件中产生的电动势 $e_k$ 与 $e_L$ 相反，即
\begin{equation}
e_k\approx-\left(e_L+e_a\right).
\end{equation}
换向极极性的确定：对\textbf{发电机}，换向极极性与电枢旋转方向上前一个主磁极的极性\textbf{相同}；对\textbf{电动机}，则\textbf{相反}。换向极磁路应不饱和，以保证 $e_k\propto I_a$ 在整个负载范围内都能补偿。

\subsection{移刷（移动电刷位置）}
把电刷从几何中性线移开一个角度，使换向元件处于主磁场中，利用主磁场产生的电动势抵消 $e_L$。发电机应\textbf{顺转向}移刷，电动机应\textbf{逆转向}移刷。缺点是只对某一负载有效，负载变化时补偿效果变差，且移刷会引入直轴电枢反应，故工程上只作为辅助手段。

\subsection{减小电抗电动势}
\begin{itemize}
  \item 减小换向元件的匝数（每元件匝数取 1）；
  \item 减小元件的电感（缩短端部、增大槽口）；
  \item 适当增大换向片数、减小片间电压。
\end{itemize}

\subsection{选用合适的电刷}
电刷的接触特性对换向影响很大：石墨电刷接触电阻大、抑制火花能力强，但允许的电流密度小；金属石墨电刷导电好、电流密度大，但换向性能差。应按电流密度与换向要求选择。

\section{补偿绕组}

\subsection{作用}
换向极只能补偿几何中性线附近的换向，无法消除交轴电枢反应对主磁场的\textbf{畸变}。当电机负载重、过载大、或采用晶闸管供电（电流脉动大）时，畸变会使极靴下一侧磁密过高、另一侧过低，导致片间电压过高而环火。为此需装设\textbf{补偿绕组}。

\subsection{结构与连接}
补偿绕组嵌放在\textbf{主磁极极靴的槽内}，其导体中电流方向与该极下电枢导体电流方向\textbf{相反}，从而直接抵消交轴电枢反应磁动势。补偿绕组与电枢绕组\textbf{串联}，以保证补偿作用随负载自动成正比变化。

\subsection{特点}
\begin{itemize}
  \item 补偿效果好，可使气隙磁场基本不畸变，允许电机做大容量、高过载；
  \item 结构复杂、造价高（极靴需开槽、嵌线），故只用于大型或特殊要求的直流电机；
  \item 补偿绕组使换向极的负担减轻，二者常配合使用。
\end{itemize}

\begin{center}
\begin{tabular}{@{}p{2.2cm}p{3.8cm}p{4.2cm}p{3.4cm}@{}}
\toprule
措施 & 针对的问题 & 原理 & 适用范围 \\
\midrule
换向极 & 换向元件中的电抗电动势 $e_L$ & 串联绕组产生反向电动势 & 几乎所有直流电机 \\
移刷 & $e_L$ 引起的延迟换向 & 利用主磁场产生补偿电动势 & 辅助手段，负载固定时 \\
减小元件电感 & $e_L$ 本身 & 减少匝数、缩短端部 & 设计阶段 \\
补偿绕组 & 交轴电枢反应造成的磁场畸变 & 极靴槽内绕组直接抵消 & 大型、重载、晶闸管供电 \\
\bottomrule
\end{tabular}
\end{center}

\corepoint{本章的因果链要背熟：\textbf{换向 = 电流急剧反向 $\rightarrow$ 电感产生 $e_L$（阻碍变化）$\rightarrow$ 换向延迟 $\rightarrow$ 后刷边火花}。改善手段都是"制造一个反向电动势"：换向极（串联、随负载自动补偿）是主力，移刷是辅助，补偿绕组解决的是另一个问题——\textbf{交轴电枢反应引起的磁场畸变}（不是换向本身）。考试常问两者的区别，务必分清"换向极管换向、补偿绕组管畸变"。}

\end{document}